% The document class supplies options to control rendering of some standard
% features in the result.  The goal is for uniform style, so some attention 
% to detail is *vital* with all fields.  Each field (i.e., text inside the
% curly braces below, so the MEng text inside {MEng} for instance) should 
% take into account the following:
%
% - author name       should be formatted as "FirstName LastName"
%   (not "Initial LastName" for example),
% - supervisor name   should be formatted as "Title FirstName LastName"
%   (where Title is "Dr." or "Prof." for example),
% - degree programme  should be "BSc", "MEng", "MSci", "MSc" or "PhD",
% - dissertation title should be correctly capitalised (plus you can have
%   an optional sub-title if appropriate, or leave this field blank),
% - dissertation type should be formatted as one of the following:
%   * for the MEng degree programme either "enterprise" or "research" to
%     reflect the stream,
%   * for the MSc  degree programme "$X/Y/Z$" for a project deemed to be
%     X%, Y% and Z% of type I, II and III.
% - year              should be formatted as a 4-digit year of submission
%   (so 2014 rather than the academic year, say 2013/14 say).



\documentclass[ % the name of the author
                    author={Revalda Putawara},
                % the name of the supervisor
                supervisor={Dr. James Pope},
                % the degree programme
                    degree={MSc},
                % the dissertation    title (which cannot be blank)
                     title={Robustness of Bayesian Neural Networks Design Choices Against Weight and Bias Perturbation Induced by Single Event Upsets (SEUs)},
                % the dissertation subtitle (which can    be blank)
                %  subtitle={And those including an optional subtitle too, for %good measure},
                % the dissertation     type
                      type={},
                % the year of submission
                      year={2025}]{dissertation}

\usepackage{booktabs} %TODO added to accomodate bottomrule etc
\usepackage{float} %TODO add floating hard constraint placement
\usepackage{enumitem}
\begin{document}

\newcommand{\B}[2]{\ifnum#1=1 \textbf{#2}\else #2\fi} % added to modify how the text in a table is neatly bold

% =============================================================================

% This section simply introduces the structural guidelines.  It can clearly
% be deleted (or commented out) if you use the file as a template for your
% own dissertation: everything following it is in the correct order to use 
% as is.

\iffalse
\section*{Prelude}
\thispagestyle{empty}


A typical MSc thesis will be structured according to a fairly standard sequence of 
sections, and one example is shown here.  However, it is hard and perhaps even 
counter-productive to generalise: the goal of this document is {\em not}\/ to be prescriptive, to show you how you {\em must} structure your thesis; rather it is to simply to act as a guideline, to show you how you {\em could} structure your thesis.  In particular, each chapter's indicated page-count is
important but {\em not}\/ absolute: the aim is simply to highlight that a 
clear and concise description is better than a rambling alternative that makes
it hard to separate important content and facts from irrelevant trivia.

You can use this document as a \LaTeX-based 
template for your own thesis by simply deleting extraneous sections
and content, and adding your own text. If you do that, we recommend using BibTeX
to deal with the associated bibliography. 
You may want to use Overleaf (an online cloud-based collaborative \LaTeX platform).


The words ``thesis'' and ``dissertation'' are used interchangeably here. All active hyperlinks in the PDF file are rendered in the same shade of dark blue. 
\fi


% =============================================================================

% This macro creates the standard UoB title page by using information drawn
% from the document class (meaning it is vital you select the correct degree 
% title and so on).


\maketitle

% After the title page (which is a special case in that it is not numbered)
% comes the front matter or preliminaries; this macro signals the start of
% such content, meaning the pages are numbered with Roman numerals.

\frontmatter

% This macro creates the standard UoB declaration; on the printed hard-copy,
% this must be physically signed by the author in the space indicated.

\makedecl

% LaTeX automatically generates a table of contents, plus associated lists 
% of figures, tables and algorithms.  The former is a compulsory part of the
% dissertation, but if you do not require the latter they can be suppressed
% by simply commenting out the associated macro.

\tableofcontents
\listoffigures
\listoftables
\listofalgorithms
%\lstlistoflistings

% The following sections are part of the front matter, but are not generated
% automatically by LaTeX; the use of \chapter* means they are not numbered.

% -----------------------------------------------------------------------------

\chapter*{Abstract}

\iffalse
{\bf A compulsory section, of at most $1$ page} 
\vspace{1cm} 

\noindent
This section should summarise the project context, aims and objectives,
and main contributions (e.g., deliverables) and achievements.  The goal is to ensure that the 
reader is clear about what the topic is, what you have done within this 
topic, {\em and}\/ {\bf what your view of the outcome is.}

Essentially 
this section is a (very) short version of what is typically covered in more depth in the first 
chapter.  If appropriate, you should include here  
a clear statement of your research hypothesis.  This will obviously differ significantly
for each project, but an example might be as follows:

\begin{quote}
My research hypothesis is that a suitable genetic algorithm will yield
more accurate results (when applied to the standard ACME data set) than 
the algorithm proposed by Jones and Smith, while also executing in less
time.
\end{quote}

\noindent
The latter aspects should (ideally) be presented as a concise, factual 
list of the main points of achievement.  Again the points will differ for each project, but 
an might be as follows:

\begin{quote}
\noindent
\begin{itemize}
\item I spent $120$ hours collecting material on and learning about the 
      Java garbage-collection sub-system. 
\item I wrote a total of $5000$ lines of {\em Python} source code, and associated orchestration scripts. 
\item I designed a new algorithm for computing the non-linear mapping 
      from A-space to B-space using a genetic algorithm.
\item I implemented a version of the algorithm proposed by Jones and 
      Smith (2010), corrected a mistake in it, and 
      compared the results with several alternatives.
\end{itemize}
\end{quote}
\fi

\iffalse
Bayesian Neural Networks (BNNs) introduce uncertainty into deep learning models by treating weights and biases as probability distributions instead of exact values. This properties make BNNs a promising candidate for deployment in harsh environment where reliability under hardware faults is essential. Single Event Upsets (SEUs), an event of which hardwares experience bitflips due to radiation, are one of the most critical source of faults in a machine learning system. Previous research has been primarily focused on the robustness of deterministic Convolutional Neural Networks against SEU attacks or robustness of Bayesian Neural Networks against adversarial attacks. There are no significant research done on the robustness of Bayesian Neural Networks against SEU injection yet.

In this research, deterministic and Bayesian CNN models were trained under 4 different design choices, which are base model, model with modified pooling mechanism in convolutional layer, model trained with dropout layer, and model trained with weight decay for its optimizer. Each design then was further evaluated under 7 different activation functions. Bayesian Neural Networks in this project uses 3 different distribution priors, which are Gaussian, Laplace, and Uniform prior. These model variations are then subject to SEU injection targetting different layers and layer parameters, bit positions, and neuron positions, in order to comprehensively assess robustness. A new robustness metric, the Aggregate Robustness Index (ARIn), is developed to simultaneously measure change in accuracy and softmax difference, two of the robustness measure that was introduced in previous researches.

This experiment showed that deterministic CNN models achieved higher baseline test accuracy (97.55\% vs. 82.85\%), but Bayesian CNN models showed superior robustness with an ARIn score of 0.0779 compared to Deterministic CNNs' 0.1909. Our findings show that the model with dropout layers during training is the most robust for the Bayesian CNN model. We also find that the Gaussian prior is the most robust, while also finding that the spread parameter (variance) is more vulnerable against SEU injection compared to an attack on the center parameter (mean). Among all activation functions tested, Rectified Weighted Gaussian is the most robust, and the trade-off between test accuracy and SEU robustness is observed while looking at different activation functions. The fully connected layer is less robust due to the mechanism of how SEU-injected value can directly affect the logits and subsequently the final class prediction. On the bit location of the attack, the first exponent bit suffers the most devastating effect on our neural network robustness performance. In this study, flipping a bit from 1 to 0 is found to be more robust compared to the opposite. %Lastly, we found no significant difference in robustness in attacking different neuron locations, which, in our case, the first and the last neuron were tested.

Overall, this work concludes that Bayesian CNNs are not universally more robust than deterministic CNNs, but they provide greater stability in prediction confidence. This advantage make them especially suitable for applications where softmax probabilities are crucial, such as medical imaging and ensemble learning systems.
\fi

Bayesian Neural Networks (BNNs) introduce uncertainty into deep learning models by treating weights and biases as probability distributions instead of exact values, making them promising for deployment in harsh environments where reliability under hardware faults is crucial. Single Event Upsets (SEUs), caused by radiation, are a major fault source in machine learning systems. While prior work has examined deterministic CNNs under SEUs and BNNs under adversarial attacks, little research addresses BNN robustness against SEUs. 

In this research, deterministic and Bayesian CNN models were trained under 4 different design choices: base model, modified pooling, dropout, and weight decay. Each design was tested with seven activation functions. For BNNs, three priors (Gaussian, Laplace, Uniform) were evaluated. Models were subjected to SEU injections targetting different layers, layer parameters, bit positions, and neuron positions to assess robustness. A new metric, the Aggregate Robustness Index (ARIn), was introduced to jointly measure accuracy difference and softmax difference.

Results show that deterministic CNNs achieve a higher baseline test accuracy (97.55\% vs. 82.85\%), but Bayesian CNN models showed superior robustness with an ARIn score of 0.0779 compared to Deterministic CNNs' 0.1909. Dropout improved BNN robustness the most, with Gaussian priors outperforming Laplace and Uniform. SEU attacks on the variance parameters were more damaging than those on means. Among activation functions, Sigmoid proved to be the most robust overall, though with an accuracy-robustness trade-off. Fully connected layers were more vulnerable, as SEU directly impacted the logits. First exponent bit-flips are the most harmful SEU, and flipping 1 to 0 showed a better robustness.

Overall, this work concludes that Bayesian CNNs are not universally more robust than deterministic CNNs, but they provide greater stability in prediction confidence. This makes them suitable for applications where class probabilities are crucial, such as medical imaging and ensemble learning systems.



%Keywords: Bayesian Neural Networks, Robustness, Single Event Upsets, Bitflips, Softmax Difference

% -----------------------------------------------------------------------------

%\chapter*{Summary of Changes}

\iffalse
{\bf A conditional section, of at most $1$ page} 
\vspace{1cm} 

If (and only if) the dissertation represents a resubmission (e.g., as the result of
a resit), then this section is compulsory: the content should summarise all
non-trivial changes made to the initial submission.  Otherwise you can
omit it, since {\bf a summary of this type is only needed for resubmissions}.

When included, the section will ideally be used to highlight additional
work completed, and address criticism raised in any associated feedback.
Clearly it is difficult to give generic advice about how to do so, but
an example might be as follows:

\begin{quote}
\noindent
\begin{itemize}
\item Feedback from the initial submission criticised the design and 
      implementation of my genetic algorithm, stating ``there seems 
      to have been no attention to computational complexity during the
      design, and obvious methods of optimisation are missing within
      the resulting implementation''.  Chapter 3 now includes a
      comprehensive analysis of the algorithm, in terms of both time
      and space.  While I have not altered the algorithm itself, I
      have included a cache mechanism (also detailed in Chapter 3)
      that provides a significant improvement in average run-time.
\item I added a feature in my implementation to allow automatic rather
      than manual selection of various parameters; the experimental
      results in Chapter 4 have been updated to reflect this.
\item Questions after the presentation highlighted a range of related
      work that I had not considered: I have make a number of updates 
      to Chapter 2, resolving this issue.
\end{itemize}
\end{quote}
\fi

% -----------------------------------------------------------------------------

\chapter*{Supporting Technologies}

%{\bf A compulsory section, of at most $1$ page}
%\vspace{1cm} 

%\noindent
%This section should present a detailed summary, in bullet point form, 
%of any third-party resources (e.g., hardware and software components) 
%used during the project.  Use of such resources is always perfectly 
%acceptable: the goal of this section is simply to be clear about how
%and where they are used, so that a clear assessment of your work can
%result.  The content can focus on the project topic itself (rather,
%for example, than including ``I used \mbox{\LaTeX} to prepare my 
%dissertation''); an example is as follows:

This work relied on a range of supporting technologies/third-party resources. The main tools are summarized below:

\begin{quote}
\noindent
\begin{itemize}

\item Computation was carried on Amazon Web Services with Elastic Compute Cloud (EC2) to provision virtual machines with GPU.

\item {\em Python} is used as the main programming language, with {\em Anaconda} to provide packages and virtual environment.

\item Deep learning models were implemented with {\em PyTorch} (for deterministic CNN models) and  {\em Pyro} (for Bayesian CNN models).

\item Experiment result data analysis and visualization were supported by publicly available Python libraries such as {\em Pandas}, {\em Matplotlib}, and {\em Seaborn}.

\item Version control was maintained through {\em Github}.
 
\item This dissertation was formatted using \LaTeX\ using the cloud-based editor {\em Overleaf}. 
\end{itemize}
\end{quote}

% -----------------------------------------------------------------------------

\chapter*{Notation and Acronyms}

%{\bf An optional section, of roughly $1$ or $2$ pages}
%\vspace{1cm} 

%\noindent
%Any well written document will introduce notation and acronyms before
%their use, {\em even if} they are standard in some way: this ensures 
%any reader can understand the resulting self-contained content.  

%Said introduction can exist within the dissertation itself, wherever 
%that is appropriate.  For an acronym, this is typically achieved at 
%the first point of use via ``Advanced Encryption Standard (AES)'' or 
%similar, noting the capitalisation of relevant letters.  However, it 
%can be useful to include an additional, dedicated list at the start 
%of the dissertation; the advantage of doing so is that you cannot 
%mistakenly use an acronym before defining it.  A limited example is 
%as follows:

For clarity and ease of reference, this chapter summarizes the mathematical notation and acronyms used throughout the dissertation. The notation and acronyms are defined as follows and will be used consistently throughout this dissertation, unless explicitly stated otherwise.

\begin{quote}
\noindent
\begin{tabular}{lcl}
BNN                 &:     & Bayesian Neural Networks                                         \\
CNN                 &:     & Convolutional Neural Networks                                         \\
DNN                 &:     & Deterministic Neural Networks                                          \\
SEE                 &:     & Single Event Effect                                         \\
SEU                 &:     & Single Event Upset                                         \\
WG                 &:     & Weighted Gaussian                                         \\
RWG                 &:     & Rectified Weighted Gaussian                                         \\
ELBO                 &:     & Evidence Lower Bound                                    \\
VI                 &:     & Variational Inference                                        \\
SVI                 &:     & Stochastic Variational Inference                                        \\
MCMC                 &:     & Markov Chain Monte Carlo                                        \\
HMC                 &:     & Hamiltonian Monte Carlo                                      \\
AAAD                &:      & Average Absolute Accuracy Difference \\
ASD                 &:      & Average Softmax Difference \\
ARIn                &:      & Aggregate Robustness Index \\
%                    &\vdots&                                                                      \\
$p(\theta)$         &:     & posterior distribution over parameters \\   
$p(\theta \mid y)$        &:     & posterior distribution over parameters given data\\
$q_\phi(z)$         &:      & variational distribution \\
$p_\theta(z \mid x)$&:      & true posterior \\
$\mathrm{KL}\!\left(q_\phi(z)\,\|\,p_\theta(z \mid x)\right)$&:     & Kullback-Leibler divergence between variational distribution and true posterior\\   

\end{tabular}
\end{quote}

% -----------------------------------------------------------------------------

\chapter*{Acknowledgements}

\iffalse
{\bf An optional section, of at most $1$ page}
\vspace{1cm} 

\noindent
It is common practice (although totally optional) to acknowledge any
third-party advice, contribution or influence you have found useful
during your work.  Examples include support from friends or family, 
the input of your Supervisor and/or Advisor, external organisations 
or persons who  have supplied resources of some kind (e.g., funding, 
advice or time), and so on.

\vspace{1cm}
Dave Cliff writes here to say huge thanks to his colleague Dr Dan Page for sharing this \LaTeX\ thesis template, which was originally written by Dan, for Computer Science dissertations. Dave edited Dan's original to better suit the needs of the Data Science MSc: please don't hassle Dan about any of this, but do feel free to contact Dave if you have any questions or comments on it.  

\begin{enumerate}
\item Alleya Hanifathariane Naud'a
\item Indonesia Endowment Fund for Education Agency (LPDP)
\item Warga Bristol 2024
\item Dr. James Pope
\end{enumerate}
\fi

First and foremost, I would like to express my deepest gratitude to my wife, Alleya, whose unconditional love and unwavering support have accompanied me every day throughout this journey. Her encouragement to pursue this master's program was the spark that set everything in motion, and her presence has been my anchor, helping me persevere until the very end.

I am profoundly grateful to the Indonesia Endowment Fund for Education Agency (LPDP) for placing their trust in me through the scholarship program. This opportunity has truly been life-changing and has opened doors I could only have dreamed of before.

My heartfelt thanks also go to Warga Bristol 2024, a community of fellow Indonesian students who began this journey alongside me. Their companionship, shared laughter, coffee sessions, comforting food that reminded me of home, and countless conversations about everyday matters have made this experience richer and more meaningful.

Finally, I extend my sincere appreciation to my supervisor, Dr. James Pope, whose guidance and constructive advice have continually motivated me to deliver my best work. His support has been invaluable in shaping the progress and outcome of this project.



% =============================================================================

% After the front matter comes a number of chapters; under each chapter,
% sections, subsections and even subsubsections are permissible.  The
% pages in this part are numbered with Arabic numerals.  Note that:
%
% - A reference point can be marked using \label{XXX}, and then later
%   referred to via \ref{XXX}; for example Chapter\ref{chap:context}.
% - The chapters are presented here in one file; this can become hard
%   to manage.  An alternative is to save the content in seprate files
%   the use \input{XXX} to import it, which acts like the #include
%   directive in C.

\mainmatter

% -----------------------------------------------------------------------------

\chapter{Introduction}
\label{chap:introduction}

\iffalse
{\bf A compulsory chapter, roughly 10\% of the total page-count. Blackboard suggestion: 2}
\vspace{1cm} 

% putting a \noindent before the first para in each chapter looks nicer.
\noindent
This chapter should describe the project context, and motivate each of
the proposed aims and objectives.  Ideally, it is written at a fairly 
high-level, and easily understood by a reader who is technically 
competent but not an expert in the topic itself.

In short, the goal is to answer three questions for the reader.  First, 
what is the project topic, or problem being investigated?  Second, why 
is the topic important, or rather why should the reader care about it?  
For example, why there is a need for this project, who will benefit from the 
project and in what way (e.g., clients/end-users who needed some analysis
done, or other data scientists who might need the tools you have developed), what 
work does the project build on and why is the selected approach either
important and/or interesting (e.g., fills a gap in literature, applies
results from another field to a new problem).  Finally, what are the 
central challenges involved and why are they significant? 
 
The chapter should conclude with a concise bullet point list that 
summarises the aims, objectives, {\bf and achievements}\/ of your work. 

\noindent\makebox[\linewidth]{\rule{\paperwidth}{0.4pt}}
\fi

%This chapter will explain the background of the research, the motivation behind the research, along with the goal of the research and challenges involved in it. We will introduce the concept of Robust Neural Networks and Bayesian Neural Networks, and how the aforementioned system will be effected by weight/bias perturbation.

This chapter introduces the foundations of the research, including the research background, motivation, objectives, and challenges. It begins with a discussion on machine learning models' role in real-world deployment under normal and harsh environments. The motivation of this research will then be stated, especially highlighting the limited existing work in examining the effect of weight/bias perturbation against Bayesian Neural Networks. Finally, we will present the approach of the research, by clearly stating the aim of the research and also the possible challenges that will be faced during the research.

\section{Background} % answer question 1, what is the problem being investigated

\iffalse
in this section I will explain:
\begin{enumerate}
\item Machine learning models are used for image classification
\item Reliability etc. (intro to robust ML)
\item Attack can also happen outside of the input, which is in the weight and bias of the model itself
\item The example is Neural Networks used in safety-critical condition, which is prone to SEE
\item Deterministic Neural Networks vs Bayesian Neural Networks
\item There are some suggestions to "strengthen" Deterministic NN (refer to William's research) especially against SEE
\item While Bayesian Neural Networks has been proven to be robust against adversarial attacks, will that also be robust against SEEs?
\end{enumerate}
\fi

Machine learning models are widely used in our daily life. Most of machine learning systems can be found close to us, like a recommender system, large language model or even autonomous vehicle that utilizes Convolutional Neural Network (CNN) models to steer clear from obstacles. While most of those systems are deployed in a normal environment, some of machine learning models are also deployed in harsh environments with its own specific designated purpose. 

One of machine learning system that is deployed in a harsh environment is CloudScout \cite{Giuffrida2020-ry}. CloudScout, which is installed in Hyperscout-2 satellite, part of the European Space Agency (ESA) PhiSat-1 mission, is a first in-orbit demonstration of Deep Neural Network for data processing on edge. It uses CNN to classify hyperspectral satellite images as either cloudy or not cloudy directly on-board, and only transmit the non cloudy image to the ground, significantly reducing bandwidth consumption.

Machine learning systems deployed in harsh environments are exposed to risks during its inference, and a common term to associate this performance with Deep Neural Networks is \textit{robustness} \cite{shafique_2021}. Robustness are categorized into two subproperties, security and reliability. While the security aspect of DNNs focuses on preventing information inside the network to be stolen or protecting the system to get an incorrect input (adversarial attack), the reliability aspect of robustness focuses on keeping the model to display a consistent output, parameter, or behaviour, despite changes of environmental condition or fabrication.

\section{Motivations} %answer question 2 : Why it is important, and why should people care

\iffalse
in this section I will explain:
\begin{enumerate}
\item What is the impact of SEE?
\item Is there any real-life life-threatening accidents caused by SEE?\
\item What is SEE's impact to Neural Networks (by bitflips etc.)
\item What does Robustness means and how it can mitigate the impact of SEE
%\item Related Work
\end{enumerate}
\fi

Being deployed in the space, a harsh environment, we can assume that machine learning systems like CloudScout is prone to risks that exists due to being deployed outside of earth surface. One of the risks that could hamper the performance of machine learning models that are deployed in harsh environments is single event effects (SEEs).

One example of SEE that can happen to machine learning models deployed in a harsh environment is Single Event Upsets (SEUs). SEUs, explained in \cite{petersen_single_2011}, are a temporary change of memory or control bit. One phenomenon that occurs and has shown the characteristic of SEUs is the incident of bitflips.

Machine learning with a good quality of robustness, or Robust Machine Learning Systems, can mitigate the impact of SEUs/bitflips, rendering the output of the model to be consistent. In this research, we are motivated to understand the impact of SEUs/bitflips to weight/bias perturbation in Bayesian Neural Networks. 

This research is built upon work done in \cite{dennispoperobustcnn} that reviewed CNN design choices to develop robust machine learning models. This research is also motivated by previous research that has proved that Bayesian Neural Networks are more robust against adversarial attacks (security aspect of robustness). Although there are already a significant number of studies that addressed those two aforementioned fields, \textbf{there has been little to no investigation into the robustness of Bayesian Neural Networks against weight/bias perturbation} (reliability aspect of robustness). This dissertation seeks to address that gap.

%Building upon the hypotheses that Bayesian Neural Networks will also demonstrate robustness against weight/bias perturbation (reliability aspect of robustness), this research is trying to prove the aforementioned hypotheses, alongside a comparison of Bayesian Neural Networks design choice that can show a better robustness.

\section{Our Approach}

In this research, we will show multiple Bayesian Neural Network design choice and compare its performance measure. Subsequently, we will simulate Single Event Upset (SEUs) with bitflips across all Bayesian Neural Network design alternatives and measure its performance measures afterwards. Robustness is then measured from the difference on the result of the model, before and after the model is perturbed with SEUs.

\subsection{Research Aim} \label{subsection:researchaim}

\iffalse
in this section I will explain the aim of my research:
\begin{enumerate}
\item Build a Bayesian CNN model and observe its performance
\item Vary the prior distributions, distribution parameter, and activation function, and see what is the performance
\item simulate a SEU condition, then measure each model's performance. Compare all of them, which one give the most robust performance
\item (if possible) suggest a new method that is more robust
\end{enumerate}
\fi

To find a Bayesian Neural Network design choice that is robust against perturbation of weight/bias, we aim to complete several tasks in this research:
\begin{enumerate}
\item Train multiple deterministic and Bayesian Neural Networks with different design choices and note their performance
\item Do multiple experiment of SEU injection on the different deterministic and Bayesian Neural Network design and note its performance
\item Compare the robustness of each different design
\item Validate whether Bayesian Neural Networks are more robust compared to deterministic Neural Networks
\end{enumerate}

\subsection{Challenges} % answer question 3
\iffalse
in this section I will explain potential challenges to the project:
\begin{enumerate}
\item "Uncertainty factor" in bayesian neural networks means that performance from deterministic CNN does not directly translate to its bayesian counterpart
\item "Generalization" of weights and bias by approximating them with the same distribution and its parameter across layer might "make the model perform worse"
\item The "Real SEE" is not available to replicate in a lab with machines, thus "Simulated". It might not be 100\% similar to the real world.
\item Magnitude of SEE (bitflips) might differ between neurons (bitflip at index 2 for weight x might be resulting in a much smaller scale compared to bitflip at index 2 for weight x+1)
\end{enumerate}
\fi

The challenge that is imminent in this research is mainly due to the Bayesian nature of the BNNs. The challenge also come from simulating the process of SEU. The challenges are as follows:
\begin{enumerate}
\item Due to the uncertainty of model weight and bias selection in Bayesian Neural Networks, the performance observed from it does not replicate its deterministic counterpart completely.
\item All models are designed to have all their layers and neurons to be approximated with the same parameter distribution, potentially making the resulting model sub-optimal towards the solution.
\item Magnitude of the result of SEU/bitflips might differ between neurons, making the measurement of robustness not a straightforward process.
\end{enumerate}
% -----------------------------------------------------------------------------

\chapter{Background}
\label{chap:background}

\iffalse
{\bf A compulsory chapter, roughly 20\% of the total page-count. Blackboard suggestion: 5}
\vspace{1cm} 

\noindent
This chapter is intended to describe the technical basis on which execution
of the project depends.  The goal is to provide a detailed explanation of
the specific problem at hand, and existing work that is relevant such as an
existing algorithm that you use, alternative solutions proposed, supporting
technologies, and relevant literature. The literature you include should cover 
appropriate peer-reviewed academic publications and books, and maybe also 
news and current-affairs articles published in reputable sources such as 
{\em The Economist} magazine or {\em The Financial Times} newspaper. 

Your thesis should include complete set of references/bibliography: everything that you cite should be in there, in full. This means including the publisher's name for anything that is a book; the editors and title of books that a paper appears in as one item in a larger collection (e.g. proceedings volumes and/or thematic edited collections) and the relevant page-numbers; and so on. Put most simply, we require and expect your references to look like the references in a published peer-reviewed academic paper, so you can work out whether you still have work to do by comparing your references to the reference-list on the publications that you're working from. 

Note there is a subtle difference from
this and a full-blown literature {\em survey}.  The latter might try
to capture and organise (e.g., categorise somehow) {\em all}\/ related work,
potentially offering meta-analysis, whereas here the goal is simply to
ensure that your thesis is self-contained.  Put another way, after reading 
this chapter an intelligent non-expert reader with no prior knowledge of your project should have obtained enough background to 
understand what {\em you}\/ have done (by reading subsequent sections), and then 
accurately assess your work.  You might view an additional goal as giving 
the reader confidence that you are able to absorb, understand and clearly 
communicate highly technical material.

Just as there is no single ideal structure for an MSc thesis, there is no one correct name for this chapter. You could just call it {\em Background}\/ if you wish, or {\em Technical Background}. Or if you feel you have a lot to say, you could split this chapter into two: you might have your first three chapters being:
\begin{enumerate}
\item Introduction
\item Context
\item Related Work
\end{enumerate}
\noindent
But if you prefer to use a more concise structure, you could instead use:
\begin{enumerate}
\item Introduction: Contextual Background
\item Technical Background
\end{enumerate}
\noindent
The choice is yours. 

The key thing is that by the end of these first two or three chapters, you have told the reader everything they need to know so that they can understand the rest of your thesis. This is particularly important because at least one of the people who actually examines your thesis will be a UoB academic, one of our lecturers or professors, who you can assume knows almost nothing about the specifics of what work you have done, and who is given a copy of your thesis to mark. So what you write has to adequately explain things to that examiner. You can safely assume that whoever examines your thesis is numerate, intelligent, understands programming and data analytics, etc, but you cannot safely assume that their specific individual expertise is a perfect match to the topic of your thesis (it's in that sense that the examiner is a {\em non-expert}). So, the person you write for, your {\em audience}, should be that undefined group of academics who might possibly examine your thesis: if you write for that audience and do a good job, your thesis should be understandable by a wide range of people, including potential employers and colleagues in the world of work.
\fi

%\section{Numerical Representation in Computers}

\iffalse
\_in this section, I will write about:
\begin{enumerate}
\item pytorch uses fp32. fp32 uses IEEE754, 2019
\item sign bit, exponent bits, mantissa bits
\end{enumerate}

TODO\_How computer works like transistor etc, how numbers are represented. Then, how the standards are established.
\fi

\iffalse
Institute for Electrical and Electronic Engineers (IEEE) published a report called Binary Floating Point Arithmetic Standard 754 in 1985. This document, updated in 2008 and 2019, provides standards for binary and decimal floating numbers, formats for data interchange, algorithms for rounding arithmetic operations, and for handling exceptions \cite{BurdenFaires2010}. IEEE-754 \cite{IEEE754-2019} defines five basic formats to represent a finite subset of real numbers, which are Binary format with encodings of 32 bits (binary32), 64 bits (binary64), 128 bits (binary128) and Decimal format with encodings of 64 bits (decimal64) and 128 bits (decimal128).

In \cite{BurdenFaires2010}, binary format and decimal format are mentioned by binary machine numbers and decimal machine numbers respectively. For decimal machine numbers, a float can be notated as normalized floating-point form

$$
y = 0.d_1d_2d_3{\dots}d_{k-1}d_{k} \times 10^{n} , 1 \le d_1 \le 9 \ and \ 0 \le d_i \le 9
$$

While decimal machine numbers provide better accuracy in representing float numbers compared to binary machine numbers, that system is deemed inefficient due to the need of storing the bits in decimal format. Binary machine numbers is fully supported by modern CPUs and accelerated in multiple framework of matrix operations, like CUDA cores and tensor cores. Pytorch, one of the deep learning framework, stated that it follows floating point arithmetic and IEEE 754 standard, noting the limited accuracy of single precision floating point numbers (FP32) and double precision floating point numbers (FP64) \cite{pytorchNumericalAccuracy}. 

Binary machine numbers represent float numbers in 0s and 1s or Base-2. Floating-point data in the binary interchange format are encoded in $k$ bits, with these three fields: 1-bit sign $S$, $w$-bit biased exponent, and ($t=p-1$)-bit trailing significand (mantissa) field. For example, for binary 32, we will have $w=8$ exponent bits, and $t=23$ mantissa bits.

%\begin{figure}[h]
%\caption{Binary interchange floating-point, as shown in \cite{IEEE754-2019}}
%\includegraphics[scale=0.6]{binary-interchange-floating-point-format.png}
%\centering
%\label{fig:floatingpointIllustration}
%\end{figure}

\iffalse
The encoding parameters for each interchange format is shown in table below. 

\begin{table}[h]
\centering
\begin{tabular}{|c|c|c|c|}
\hline
parameter      & binary32      & binary64  & binary128     \\
\hline
%$w     $ & $0     $ & $0     $ & $baz  $\\
%\hline
$w$, \textbf{exponent} field width in bits      & $8     $ & $11     $ & $15 $\\
\hline
$t$, trailing significand (\textbf{mantissa}) field width in bits      & $23     $ & $52     $ & $112 $\\
\hline
\end{tabular}
\caption{Binary Interchange format parameter}
\label{Binary_Interchange_format_parameter}
\end{table}
\fi

This project uses Pyro, a probabilistic programming language built upon Python and PyTorch \cite{pyroIntroductionPyro}. PyTorch has a critical role in Pyro is the stochastic gradient descent process in solving concrete optimization problems from probabilistic computations converted by Pyro. Therefore, the default setting for floating point precision in Pyro relies on PyTorch's default setting, which is floating point 32 (fp32) / binary32. To conclude, in this research, we are dealing with 32 bits, or 32 combinations of 0s and 1s as illustrated in the Figure \ref{fig:floatingpointusedIllustration} below.

\begin{figure}[h]
\caption{Architecture for storing floating point numbers in this research, which is FP32}
\includegraphics[scale=0.4]{FP32.png}
\centering
\label{fig:floatingpointusedIllustration}
\end{figure}

This concept will be useful for us to understand how a Single Event Upset (SEU), or bitflips, can impact the weight/bias distribution parameter stored in our machine learning system, which will be explained in subsequent sub chapters.
\fi


\iffalse
\section{Deep Neural Networks}

\iffalse % THIS STARTS HERE
TODO\_in this section, I will write about:
\begin{enumerate}
\item a pinch of history of neural networks, include \cite{RumelhartHintonWilliams1986}
\item convolutional neural networks, add \cite{LeCun1989}. Add diagram of convolutional layer from \cite{bishoppatternrecognition2007}. 
\item layers in convolutional neural networks, include convolution, maxpool etc.
\item forward propagation and backpropagation, including gradient descent
\item z = wx+b illustration 
\item activation function, include all the 7 act function
\end{enumerate}
\fi % AND ENDS HERE

The earliest concept of "neural network" was introduced in \cite{mcculloch43a}, in which it models nervous systems as "net of neurons" that act in an "all-or-none" process. In \cite{Rosenblatt1958ThePA}, the concept is developed to introduce perceptron, a model of learnable weights. \cite{RumelhartHintonWilliams1986} works then popularized the concept of backpropagation for training a multi-layer networks. This breakthrough enabled the creation of deep neural networks that can be trained and effective to complete the specified tasks.

\subsection{Feed-Forward} \label{subsection:feedforward}

Work in \cite{RumelhartHintonWilliams1986} also standardised the basic neural network model, which consist of series of functional transformations. The process consists of linear transformation of input followed by a non-linear transformation in a neuron, or we can call it as a feed-forward process. This non-linearity enables deep neural networks to approximate complex, non-linear functions. This process, as mentioned in \cite{bishoppatternrecognition2007} is notated as:

\begin{equation}
a_j = \sum_{i=1}^{D} w^{(1)}_{ji} x_i + w^{(1)}_{j0}
\end{equation}

%page 227
where $x_1,x_2,\dots,x_D$ are inputs, $j=1,2,\dots,M$ and the superscript (1) represents the parameters in the first 'layer' of the network, which have a 'width' of $M$. $w^{(1)}_{j0}$ is  Each of $a_j$ is then transformed in the following equation:

\begin{equation}
z_j = h(a_j)
\end{equation}

where $h(\cdot)$ is a differentiable activation function. The output of $z_j$ would then become the input of the following layers until the output neurons are reached. An example of a simple neural network with a feed-forward topology is as follows.






\begin{figure}[h]
\caption{A neural network with feed-forward topology}
\includegraphics[scale=0.3]{feed-forward.png}
\centering
\label{fig:feedforwardIllustration}
\end{figure}

An illustration on how the calculation of input $x$ multiplied by weights $w$ added with bias $b$ and followed by an activation function $\sigma$ is shown in figure \ref{fig:zwxbIllustration}.

%TODO\_modify the image to be of your own making, especially the notation in the medium picture

\begin{figure}[h]
\caption{Illustration of linear transformation of input followed by an activation function, as shown in \cite{mediumJamesUnderstandingDNN}}
\includegraphics[scale=0.6]{zwxb_illustration.png}
\centering
\label{fig:zwxbIllustration}
\end{figure}

\subsection{Activation Function}

$h(\cdot)$ , differentiable activation function discussed earlier in \ref{{subsection:feedforward}} is to be determined by the narute of the data and the assumed distribution of target variables. This activation function is essential in neural networks due to its ability to introduce non-linearity, allowing the network to learn and model complex patterns in data. In this research, we will consider the following activation functions:

\begin{enumerate}
\item Rectified Linear Unit (ReLU)
\begin{equation}
h(x) = \max(0, x)
\end{equation}
\item Sigmoid
\begin{equation}
h(x) = \frac{1}{1+\exp(-x)}
\end{equation}
\item Tanh
\begin{equation}
h(x) = \frac{\exp(x)-\exp(-x)}{\exp(x)+\exp(-x)}
\end{equation}
\item Sinusoidal
\begin{equation}
h(x) = \sin(x)
\end{equation}
\end{enumerate}


We also want to show you the added activation function in the work of \cite{dennispoperobustcnn}: Weighted Gaussian (WG) and Rectified Weighted Gaussian (RWG)

\begin{enumerate}
\item Weighted Gaussian (WG)
\begin{equation}
h(x) = x \exp(-x^{2})
\end{equation}
\item Rectified Weighted Gaussian (RWG)
\begin{equation}
h(x) = \max(0, x \exp(-x^{2}))
\end{equation}
\end{enumerate}

\subsection{Backpropagation}

Backpropagation is the process of sending the information in neural network forward and backward to evaluate the error function $E$(w). The simplest process of backpropagation is gradient descent introduced in \cite{RumelhartHintonWilliams1986}. Suppose an error function of:

%page 242

\begin{equation}
E(\mathbf{w}) = \sum_{n=1}^{N} E_n(\mathbf{w}).
\end{equation}

Consider a simple linear model below

\begin{equation}
y_k = \sum_i w_{ki} x_i
\end{equation}

with an error function of

\begin{equation} \label{eq:book5.48}
E_n = \frac{1}{2} \sum_k (y_{nk} - t_{nk})^2
\end{equation}

where $y_{nk}=y_k(x_n,w)$, the prediction of the neural network of the true value $t_{nk}$. The gradient of the aforementioned error function with respect to weight $w_{ji}$ is given by

\begin{equation} \label{eq:book5.49}
\frac{\partial E_n}{\partial w_{ji}} = (y_{nj} - t_{nj}) x_{ni}
\end{equation}

With respect to the fact that $E_{n}$ depends on the weight $w_{ji}$ only via summed input of $a_j$ to unit $j$, to solve for the same equation above we can apply the chain rule for partial derivatives, to give

\begin{equation} \label{eq:book5.50}
\frac{\partial E_n}{\partial w_{ji}} = \frac{\partial E_n}{\partial a_j} \cdot \frac{\partial a_j}{\partial w_{ji}}.
\end{equation}

If we introduce a new notation

\begin{equation} \label{eq:book5.51}
\delta_j \equiv \frac{\partial E_n}{\partial a_j}
\end{equation}

where $\delta$ represents error. Due to the fact that:

\begin{equation} \label{eq:book5.52}
a_j = \sum_i w_{ji} z_i
\end{equation}

we can write 

\begin{equation}
\frac{\partial a_j}{\partial w_{ji}} = z_i.
\end{equation}

substituting (\ref{eq:book5.51}) and (\ref{eq:book5.52}) into (\ref{eq:book5.50}), we then get the following equation

\begin{equation}
\frac{\partial E_n}{\partial w_{ji}} = \delta_j z_i.
\end{equation}

The above equation shows us that the derivative is obtained by multiplying the value of $\delta$ for the unit at the output end of the weight by the value of $z$ for the unit at the input end of the weight (where $z$ = 1 in the case of a bias).

For the output unit, error ($\delta$)is notated as

\begin{equation}
{\delta j} = y_k - t_k
\end{equation}

To evaluate the $\delta$ for hidden units, chain rule for partial derivatives can also be used:

\begin{equation} \label{eq:book5.55}
\delta_j \equiv \frac{\partial E_n}{\partial a_j} = \sum_k \frac{\partial E_n}{\partial a_k} \frac{\partial a_k}{\partial a_j}
\end{equation}

If we substitute the definition of $\delta$ given by (\ref{eq:book5.51}) into  (\ref{eq:book5.55}) and use (\ref{eq:book5.48}) and (\ref{eq:book5.49}), we can define the backpropagation formula

\begin{equation} \label{eq:book5.56}
\delta_j = h'(a_j) \sum_k w_{kj} \delta_k
\end{equation}

For batch methods, the equation (\ref{eq:book5.50}), which is the derivative of the total error $E$, can then be obtained by repeating the backpropagation process above for each data in the training set and then summing for all data:

\begin{equation}
\frac{\partial E}{\partial w_{ji}} = \sum_n \frac{\partial E_n}{\partial w_{ji}}.
\end{equation}

\fi

%\subsection{Convolutional Neural Networks (CNN)}

\section{Convolutional Neural Networks (CNN)}

% bishop book 267
In machine learning models with images as its input, we can simply treat each pixel as channels of color intensity, and then build a fully connected network. In principle, this network could work and yield a good result, but ignore the key property of images, that nearby pixels are more correlated than pixels far away from it. Modern computer vision methods utilize this characteristic by extracting local features that depend only on small subregions of image. Convolutional neural network take advantage of this idea in three mechanisms: local receptive fields, weight sharing, and subsampling \cite{bishoppatternrecognition2007}.  

%A sample structure of convolutional neural network (CNN) is shown in figure \ref{fig:convolutionIllustration}. 

%\begin{figure}[h]
%\caption{Part of Convolutional Network, as shown in %\cite{bishoppatternrecognition2007}}
%\includegraphics[scale=0.3]{convolution.png}
%\centering
%\label{fig:convolutionIllustration}
%\end{figure}

Convolutional layers use local receptive fields, with each unit looking only at a small patch (for example, 4 x 4) of the input image. All of the units in that layer share the same weights, with the purpose of detecting the same pattern at different locations. The output of this process is a feature map with kernels of the image, which became the input for the sub-sampling layer.  

In the sub-sampling layer, the spatial resolution of feature maps is further reduced by pooling algorithms. There are some popular pooling algorithms that are used in the sub-sampling layer \cite{app12178643}:

\begin{enumerate}
\item Max Pooling Method: Identifies the biggest element in each pooling region. The biggest limitation of max pooling lies in how all other items are ignored completely for the sake of the biggest value observed \cite{Passricha2020}.
\begin{equation}
f_{max}({X}) = max \ \{x_1,\dots,x_N\}
\end{equation}
\item Average Pooling Method: Calculates the average values for each pooling region.
\begin{equation}
f_{\text{ave}}(X) = \frac{1}{N} \sum_{i=1}^{N} x_i
\end{equation}
%\item Mixed Pooling Method: A combination of max pooling and average pooling with a weighing factor of $\lambda$
%\begin{equation}
%S_j = \lambda \max \{a_1,\dots,a_N\} + (1 - \lambda) \frac{1}{|R_j|} \sum_{i \in R_j} a_i
%\end{equation}
\end{enumerate}

\iffalse
\begin{figure}[h]
\caption{Illustration of max pooling, as shown in \cite{app12178643}}
\includegraphics[scale=0.3]{maxpooling.png}
\centering
\label{fig:maxpoolingIllustration}
\end{figure}

\begin{figure}[h]
\caption{Illustration of average pooling, as shown in \cite{app12178643}}
\includegraphics[scale=0.3]{averagepooling.png}
\centering
\label{fig:avgpoolingIllustration}
\end{figure}
\fi

\begin{figure}[h]
\caption{Illustration of max and average pooling (2 x 2 filters and stride 2) as shown in \cite{app12178643}}
\includegraphics[scale=0.5]{poolingcombined.png}
\centering
\label{fig:maxavgpoolingIllustration}
\end{figure}

\section{Bayesian Neural Networks}

\iffalse
TODO\_in this section, I will write about:
\begin{enumerate}
\item  Callback to the Bayesian Theorem
\item History of the first research on bayesian neural networks. Include \cite{mackay1992a}. 
\item Two types of uncertainty in Bayesian neural networks: Aleatoric Uncertainty and Epistemic Uncertainty, mentioned in \cite{zhang2021a}
\end{enumerate}
\fi

\subsection{Bayes' rule and Bayesian Inference} \label{subsection:bayesruleandbayesianinference}

Joint probability distribution for parameter $\theta$ and observed value $y$ can be written as a product of two densities that are often referred to as the prior distribution $p(\theta)$ and the sampling distribution (or data distribution) $p(y|\theta)$ \cite{gelman2013bayesian}: $p(\theta, y) = p(\theta)\, p(y \mid \theta).$

Conditioning on the known value of the data $y$ using the basic property of Bayes' rule, posterior density, can be notated in Equation (\ref{eq:book_1.1}). An equivalent form of Equation (\ref{eq:book_1.1}), in which omits the factor of $p(y)$, which can be considered a constant due to its independency from $\theta$, yield the \textit{unnormalized posterior density} of Equation (\ref{eq:book_1.2}).

%\begin{equation} \label{eq:book_1.1}
%p(\theta \mid y) = \frac{p(\theta, y)}{p(y)} = \frac{p(\theta)\, p(y \mid \theta)}{p(y)}
%\end{equation}



%\begin{equation}  \label{eq:book_1.2}
%p(\theta \mid y) \propto p(\theta)\, p(y \mid \theta).
%\end{equation}

While Bayes' rule is often used to infer unknown model parameters $\theta$, Bayesian Inference is more general and can be applied to latent variables, evidence evaluations, and predictions. In this broader context, we often introduce latent variables $z$ and compute distributions over both the observed data $x$ and hidden variables $z$. $\theta$ defines the model, inference focuses on latent variables $z$, and $\theta$ is learned via evidence maximization. Drawing conclusions about latent variables from data can be done by computing posterior over latent variables $z$, as shown in Equation (\ref{eq:posteriorlatent})

%\begin{equation} \label{eq:posteriorlatent}
%p_{\theta}(z \mid x) = \frac{p_{\theta}(x, z)}{\int dz \; p_{\theta}(x, z)}
%\end{equation}

How well a model fits observed data $x$ can be quantified with Evidence (marginal likelihood), which can be defined with Equation (\ref{eq:marginallikelihood}). To make predictions for new data, we can utilize the posterior predictive distribution in Equation (\ref{eq:posteriorpredictivedistribution}). In order to learn the parameters $\theta$ of our models from observed data $x$, we can maximize the marginal likelihood as shown in Equation (\ref{eq:thetamle}).

%\begin{equation} \label{eq:marginallikelihood}
%p_{\theta}(x) = \int dz \; p_{\theta}(x, z)
%\end{equation}

%\begin{equation} \label{eq:posteriorpredictivedistribution}
%p_{\theta}(x' \mid x) = \int dz \; p_{\theta}(x' \mid z) \; p_{\theta}(z \mid x)
%\end{equation}

%\begin{equation}
%\theta_{\text{max}} = \arg\max_{\theta} \; p_{\theta}(x) = \arg\max_{\theta} \int dz \; p_{\theta}(x, z)
%\end{equation} 

\noindent
\begin{minipage}{0.48\linewidth}
\begin{equation}
p(\theta \mid y) = \frac{p(\theta, y)}{p(y)} 
= \frac{p(\theta)\,p(y \mid \theta)}{p(y)}
\label{eq:book_1.1}
\end{equation}
\end{minipage}\hfill
\begin{minipage}{0.48\linewidth}
\begin{equation}
p(\theta \mid y) \propto p(\theta)\,p(y \mid \theta)
\label{eq:book_1.2}
\end{equation}
\end{minipage}

\vspace{1em}

\noindent
\begin{minipage}{0.48\linewidth}
\begin{equation}
p_\theta(z \mid x) = \frac{p_\theta(x,z)}{\int dz\, p_\theta(x,z)}
\label{eq:posteriorlatent}
\end{equation}
\end{minipage}\hfill
\begin{minipage}{0.48\linewidth}
\begin{equation}
p_\theta(x) = \int dz\, p_\theta(x,z)
\label{eq:marginallikelihood}
\end{equation}
\end{minipage}

\vspace{1em}

\noindent
\begin{minipage}{0.48\linewidth}
\begin{equation}
p_\theta(x' \mid x) = \int dz\, p_\theta(x' \mid z)\, p_\theta(z \mid x)
\label{eq:posteriorpredictivedistribution}
\end{equation}
\end{minipage}\hfill
\begin{minipage}{0.48\linewidth}
\begin{equation}
\theta_{\max} = \arg\max_\theta p_\theta(x) 
= \arg\max_\theta \int dz\, p_\theta(x,z)
\label{eq:thetamle}
\end{equation}
\end{minipage}


\subsection{Prior Distribution} \label{section:backgroundpriordist}

Prior ($p(\theta)$) represents the belief about parameter $\theta$ before observing any data. If this prior shows genuine, context-specific information about the parameter, it is referred as an \textbf{informative prior}. The choice of prior not only shows our assumptions or existing knowledge but can also influence the posterior result. To see the effect on chosing a prior to the posterior, we observe three types of prior distributions.

%: Gaussian prior, which favors values near a central mean and is widely used due to its mathematical convenience and flexibility, Laplace prior, which favors values near a central mean but with a sharper tail, and Uniform prior, which assumes no prior preference over the parameter space (uninformative).

%There are two main perspectives on interpreting prior distributions \cite{gelman2013bayesian}, which are population interpretation and state of knowledge interpretation.

%In \textbf{population interpretation}, the prior represents a distribution over a population of possible parameter values from which the current value of interest is assumed to be drawn. On the other hand, in \textbf{state of knowledge interpretation}, the prior is treated as an expression of uncertainty or knowledge about the parameter, recognizing that the parameter value is fixed but unknown.

\textbf{Gaussian distribution/Normal distribution} is highly mentioned with the central limit theorem, which states that if a random variable is the sum of a large number of independent random variables, it is normally distributed \cite{ricemathstats2007}. The density function of parameter $\theta$ that depends of two parameters, mean $\mu$ and standard deviation, $\sigma$ (where $-\infty$ $<$ $\mu$ $<$ $\infty$, $\sigma$ $>$ 0) is shown in Equation (\ref{eq:gaussian}).

%\begin{equation}
%f(x) = \frac{1}{\sigma \sqrt{2\pi}} e^{-\frac{(x - \mu)^2}{2\sigma^2}}
%\end{equation}

%A statement of "X follows a normal distribution with parameters of $\mu$ and $\sigma$ can be notated as $X \sim N (\mu,\sigma^2)$.

\textbf{Laplace distribution} is a heavy-tailed distribution that is also known as double-sided exponential distribution \cite{murphy2012}. The density function of the Laplace distribution is given by Equation (\ref{eq:laplace}). $\mu$ is a location parameter and $b > 0$ is a parameter of scale. Laplace distribution is deemed to be more robust to outliers compared to the Gaussian distribution \cite{murphy2012}.

%\begin{equation}
%f(x) = \frac{1}{2b} \exp\left( -\frac{|x - \mu|}{b} \right)
%\end{equation}



%, as shown in figure \ref{fig:laplaceIllustration}.

%\begin{figure}[h]
%\caption{Illustration of the effect of outliers on fitting Gaussian, Student, and Laplace distributions, as shown in
%\includegraphics[scale=0.4]{laplace_against_outlier.png}
%\centering
%\label{fig:laplaceIllustration}
%\end{figure}

If we know nothing about $\theta$, we can use a \textbf{Uniform} prior, which is a kind of uninformative prior \cite{murphy2012}, as shown in Equation (\ref{eq:uniform}). The uniform distribution can be represented as a beta distribution with $a = b = 1$. 

%The density function of a uniform distribution, continuous over the interval of $[a,b]$ is as follows:

%\begin{equation}
%f(x) =
%\begin{cases}
%\frac{1}{b - a}, & \text{for } x \in [a, b] \\
%0, & \text{otherwise}
%\end{cases}
%\label{eq:uniformbound}
%\end{equation}

\noindent
\begin{minipage}{0.32\linewidth}
\begin{equation}
f(x) = \frac{1}{\sigma\sqrt{2\pi}} 
       \exp\!\left(-\frac{(x-\mu)^2}{2\sigma^2}\right)
\label{eq:gaussian}
\end{equation}
\end{minipage}\hfill
\begin{minipage}{0.32\linewidth}
\begin{equation}
f(x) = \frac{1}{2b} 
       \exp\!\left(-\frac{|x-\mu|}{b}\right)
\label{eq:laplace}
\end{equation}
\end{minipage}\hfill
\begin{minipage}{0.32\linewidth}
\begin{equation}
f(x) =
\begin{cases}
 \tfrac{1}{b-a}, & x \in [a,b], \\
 0, & \text{otherwise},
\end{cases}
\label{eq:uniform}
\end{equation}
\end{minipage}

%\begin{figure}[h]
%\caption{Illustration of different beta distributions parameter, which also %shows the uniform distribution \cite{murphy2012}}
%\includegraphics[scale=0.4]{beta_distributions.png}
%\centering
%\label{fig:betaIllustration}
%\end{figure}

Based on the explanation for the three priors above, we can see that Gaussian and Laplace distribution have a similar parameter, which represents the location of the mean ($\mu$) and the scale of how dispersed is the random variable of the parameters ($\sigma$ in Gaussian, and $b$ in Laplace). These paradigm are commonly used in Bayesian modelling modules like Pyro. The uniform distribution is parameterized differently compared to the other two distributions, as it is parameterized with $a$ (lows) and $b$ (highs). To have a similar sense accross all of the distributions, we introduced a parameter of $widths$, which is the difference between highs and lows, $\text{widths} = b - a$.


In this research, the two parameters for each distribution will be notated in a same notation, a and b. The table below shows the mapping of the parameters:

\begin{table}[h]
\centering
\begin{tabular}{|c|c|c|}
\hline
prior distribution      & a (in this research)      & b (in this research)      \\
\hline
%$w     $ & $0     $ & $0     $ & $baz  $\\
%\hline
Gaussian      & $\mu     $ & $\sigma     $\\
\hline
Laplace     & $\mu     $ & $b    $\\
\hline
%Uniform     & $a \ (\text{lows})     $ & $b - a \ (\text{widths})     $\\
Uniform     & $a \ (\text{lows})     $ & $\text{highs} - \text{lows} \ (\text{widths})     $\\
\hline
\end{tabular}
\caption{Prior distribution parameter used in this research}
\label{Prior_distribution_parameter}
\end{table}

\begin{figure}[h]
\caption{Prior Probability Distribution Function for $a=0$ and $b=1$ (Gaussian and Laplace) and $a=-0.5$ and $b$=0.5 (Uniform)}
\includegraphics[scale=0.6]{priorpdfs.png}
\centering
\label{fig:priorpdfs}
\end{figure}

%TODO\_ADD info on how it initializes and updates

%TODO\_ADD info ON 

\iffalse
TODO\_In this section, I will write about:
\begin{enumerate}
\item Gaussian ($\mu$, $\sigma$)
\item Laplace ($\mu$, $\sigma$)
\item Uniform ($a$, $b$)
\end{enumerate}
\fi

\subsection{Bayesian Neural Networks}

In this section, we will take a look at the concept of Bayesian Neural Networks. This research assumes that the reader is familiar with the basics of Neural Networks, including the concept of feed-forward and backpropagation. The idea of Bayesian Neural Networks (BNNs) dates back to some research done in the early 1990s. One of the early prominent works in Bayesian Neural Networks is the work done in \cite{Mackay1992a}. In that research, although not explicitly mentioning Bayesian Neural Network (BNN), David J.C. Mackay proposes a quantitative and practical Bayesian framework for learning in feedforward networks. This research revolutionizes how to think about weight and biases by giving them a probabilistic interpretation within a learning model. This work criticized the backpropagation framework's limitations, such as arbitrary hyperparameter settings, high data demands, and a lack of uncertainty estimates.

%\begin{enumerate}
%\item Objective setting of control parameters: There are not yet established ways of objectively setting the neural network parameters, such as hidden units or weight decay terms ($\alpha$), or noise level ($\beta$)
%\item Demand of data size: A Large test set is needed to reduce the signal-to-noise ratio in the test error, and cross-validation is computationally demanding.
%\item Lack of quantified uncertainty: standard backpropagation does not offer error bars on outputs or weights, which means the network can not tell which weights are more flexible and which ones are highly restrictive to maintain the model performance
%\end{enumerate}

To fix the issues on the above points, Mackay proposed a Bayesian way of doing backpropagation by the following:

\begin{enumerate}
%\item Objective setting of control parameters: The bayesian framework uses evidence maximization to set the model hyperparameters initialized with associated priors. The marginal likelihood penalizes overfitting with Occam's razor, selecting the best model complexity by maximizing Equation (\ref{eq:evidence}).
\item Objective setting of control parameters: In the Bayesian framework, model hyperparameters are determined by maximizing the marginal likelihood, which integrates over the parameter space with respect to the chosen priors. It automatically balances model fit and complexity, ensuring that the selected model achieves the best trade-off without overfitting.
%\begin{equation} \label{eq:occamsrazor}
%P(\mathcal{D} \mid \mathcal{M}) = \int P(\mathcal{D} \mid \mathbf{w}, \mathcal{M}) \, P(\mathbf{w} \mid %\mathcal{M}) \, d\mathbf{w}
%\end{equation}
\item Demand of data size: Bayesian Inference penalizes complexity automatically, as shown in Equation (\ref{eq:evidence}), by automatically trades off data fit vs model complexity without the need for a separate validation or test set to detect overfitting.
\item Lack of quantified uncertainty: Bayesian training infers a distribution over weights shown in Equation (\ref{eq:posterior}).
%\begin{equation}
%p(\mathbf{w} \mid \mathcal{D}) \propto p(\mathcal{D} \mid \mathbf{w}) \, p(\mathbf{w})
%\end{equation}
The effective number of well-defined parameters ($\gamma$), are those weights, given the data, that are actually being used by the model in a meaningful way. With the above equation, well-defined parameters can be spotted by seeing which parameters have less variance, showing the model's confidence in their value.
\end{enumerate}

\noindent
\begin{minipage}{0.48\linewidth}
\begin{equation}
P(\mathcal{D} \mid \mathcal{M}) 
= \int P(\mathcal{D} \mid \mathbf{w}, \mathcal{M}) \,
      P(\mathbf{w} \mid \mathcal{M}) \, d\mathbf{w}
\label{eq:evidence}
\end{equation}
\end{minipage}\hfill
\begin{minipage}{0.48\linewidth}
\begin{equation}
p(\mathbf{w} \mid \mathcal{D}) \propto 
   P(\mathcal{D} \mid \mathbf{w}) \, p(\mathbf{w})
\label{eq:posterior}
\end{equation}
\end{minipage}


In \cite{Mackay1992a}, approximations are made to the posterior by using Laplace approximation and use Hessian of the loss to estimate posterior curvature. The term Bayesian Neural Network first coined in \cite{Neal1996bayesianlearning}. In that work, Radford M. Neal take the posterior distribution approximation to a new direction, by using Markov Chain Monte Carlo (MCMC) to not just do evidence maximization, but also sampling from the full posterior over weights, to approximate $p(\mathbf{w} \mid \mathcal{D})$. Illustration on how Bayesian Neural Networks differ to the Deterministic Neural Networks, in which each weight is assigned a distribution, is shown in Figure \ref{fig:BNNIllustration}.

\begin{figure}[h]
\caption{Illustration of Bayesian Neural Networks, as shown in \cite{blundell2015bayesbybackprop}}
\includegraphics[scale=0.4]{bayesbybackprop.png}
\centering
\label{fig:BNNIllustration}
\end{figure}



\subsection{Posterior Distribution Approximation}

%SVI : https://pyro.ai/examples/svi_part_i.html
% https://pyro.ai/examples/intro_long.html#Inference-in-Pyro
% mean field variational approximation : 

% talk about variational inference
% then talk about stochastic variational inference

While Bayes' rule in Subsection \ref{subsection:bayesruleandbayesianinference} provide the base to compute posterior distribution over parameter or latent variables, in practice these posterior are often intractable to compute. The equations from (\ref{eq:posteriorlatent}) to (\ref{eq:posteriorpredictivedistribution}) involve high-dimensional integrals that is computionally challenging and impossible to evaluate in complex models such as Bayesian Neural Networks.

To mitigate this issue, we use approximate inference methods, and one of them is \textbf{Variational Inference}. The method first introduced in \cite{PetersonAndersonMeanField}, in which mean-field methods was introduced, to fit neural networks in minimizing a variational free-energy functional. Later work is done in \cite{HintonVanCampVariational} in employing mean-field variational inference in neural network training.

Instead of evaluating the true posterior $p_{\theta}(z|x)$ directly, Variational Inference (VI) introduces simpler tractable distributions $q_{\phi}(z)$ parameterized by $\phi$. The goal is to select $q_{\phi}(z)$ that is closest to the true posterior according to a divergence measure, typically the \textbf{Kullback-Leibler (KL) divergence}:

\begin{equation}
\mathrm{KL}\!\left(q_\phi(z)\,\|\,p_\theta(z \mid x)\right) = \mathbb{E}{q\phi(z)} \big[\log q_\phi(z) - \log p_\theta(z|x)\big]
\end{equation}

Since the true posterior is unknown, we optimize an equivalent and tractable objective called the \textbf{Evidence Lower Bound (ELBO)}. In maximizing ELBO, we ensure that the variational distribution $q_{\phi}(z)$ approximates the true posterior as close as possible, approximating the solution for Bayesian Inference in large-scale or high-dimensional models.

\begin{equation} \label{eq:genericelbo}
\text{ELBO} \equiv \mathbb{E}_{q_\phi(z)} \big[ \log p_\theta(x, z) - \log q_\phi(z) \big]
\end{equation}

When the dataset is large, computing ELBO over the entire dataset in each optimization step becomes expensive. \textbf{Stochastic Variational Inference (SVI)}, which was introduced in \cite{Hoffman2013}, combines VI with stochastic optimization. SVI subsamples mini-batches and computes estimate of the ELBO gradient.

Another method to approximate Bayesian Inference, besides Variational Inference, is \textbf{Markov Chain Monte Carlo (MCMC)}. First introduced in \cite{Metropolis:1953am} with Metropolis algorithm, the work is extended in \cite{Hastings1970} to generalize the former into Metropolis-Hastings algorithm making MCMC more broadly applicable to Bayesian inference problems. This algorithm works by constructing a Markov Chain that generates a sequence of samples that approximate draws from the true posterior distribution.

%TODO\_might be important to include the concept of mean field variational approximation


\section{Robust ML Systems}

\iffalse
TODO\_This section mainly cites the work in \cite{shafique_2021}, that will explain about:
\begin{enumerate}
\item Two threats of Machine Learning Systems, security threats and reliability threats
\item Explain about multiple examples of security threats and will deep dive on adversarial attacks
\item defence against adversarial attacks are blablablabla
\item Explain about multiple examples of reliability threats. Deep dive on transient faults. SEUs is a type of soft errors, which is a type of transient fault.
\item defense against adversarial attacks are including hardware redundancy, voting mechanism, error detection and correction codes. Maybe we can call a BNN "a system design with stronger error-correction mechanism". Redundancy based approach can incur large area overhead and cost
\end{enumerate}
\fi

Robustness in Machine Learning (ML) systems refers as a system's ability to maintain accuracy and correct behaviour under both \textbf{security threats} and \textbf{reliability threats} \cite{shafique_2021}.

\begin{enumerate}
\item Security threats

This type of threat comes from adversaries exploiting model weaknesses or implementation flaws. Common security threats are adversarial attacks, neural-level trojans, hardware attacks, and IP stealing. Adversarial attack work by adding imperceptible noise to the input, making the task of distinguishing clean and malignant inputs nearly impossible. This attack can cause misclassification, either by resulting in a random incorrect output class (untargeted attack), a specific designated output class (targeted attack), or simply reduce the confidence of the correct output class (confidence reduction). Some of the adversarial attack method are Fast Gradient Sign Method (FGSM) \cite{goodfellow2014explaining} and Jacobian saliency map approach \cite{2018-cs.LG-Wiyatno-Maximal}.

%TODO explain about adversarial attacks.
\item Reliability threats

Even without active attackers, ML systems can fail due to hardware or environmental factors. Some types of reliability threat are intermittent faults, permanent faults, and environmental noise. Soft errors are a type of transient fault mostly caused by a radiation induced change of  charge $Q_{rad}$ in the transistors across the chip. If that charge exceed the threshold value of $Q_{th}$, the state of the transistor would change, resulting in a bit-flip. This phenomenon, which is called as Single Event Upset (SEU), will be explained in the next sub chapter.


\end{enumerate}



\section{Single Event Effects}

Single Event Effects (SEEs), as described in \cite{petersen_single_2011}, happened through the action of a single ionizing particle as it penetrates sensitive nodes inside an electronic device. The effect of SEEs can be as small as a tiny glitch inside a monitor to a substantial error that is life threatening. On October 2008, a Qantas A330 flight from Singapore to Perth experienced a sudden and unexpected altitude change, resulting in 36 crew and passengers injured \cite{airsafeqantas}. The Australian Transport Safety Bureau (ATSB) considers SEE as a possible cause of the accident \cite{ATSB2009QF72}. Some variants of SEEs are Single Event Upsets (SEUs), Multibit Upsets (MBUs), and Single Event Burnouts (SEBs) \cite{petersen_single_2011}.

Illustrations of heavy ions passing through a semiconductor is shown in Figure \ref{fig:ionizationtIllustration} (a) and (b). In Figure \ref{fig:ionizationtIllustration} (a), the ion creates a dense ionization track that generates enough charge to trigger a Single Event Upset (SEU). In Figure \ref{fig:ionizationtIllustration} (b), a proton causes a nuclear reaction inside the device, producing secondary heavy ions that deposit charge and may also result in an SEU.

SEUs or bit-flips have been simulated either by injecting permanent faults (non modifiable 0s or 1s) into the weights of CNN layers at the software level to assess reliability and classify prediction deviations \cite{bosio2019}, or by injecting transient single and double bit flips into GPU registers and instruction outputs at the architecture level to study application-level Silent Data Corruptions (SDCs) \cite{hari2017sassifi}. Both studies show that SEUs can lead to incorrect or unsafe outputs in Deep Neural Networks, potentially degrading their reliability in safety-critical applications.  This research will follow a similar procedure to the SEU injection method in \cite{hari2017sassifi}, in which we will inject SEU (bitflip) on the weight/bias parameter of the model \textbf{when it is doing inference on a new data}. An example of the impact of bitflips to stored weight/bias is shown in Table \ref{table:robbothbitloc2}.

\begin{figure}[H]
\caption{Ionization path due to direct passage of heavy iron (a), and due to proton reaction in a device (b) as shown in \cite{petersen_single_2011}}
\includegraphics[scale=1.2]{ionization_path.png}
\centering
\label{fig:ionizationtIllustration}
\end{figure}

The IEEE-754 standard (\cite{BurdenFaires2010}, \cite{IEEE754-2019}) defines binary (32, 64, 128 bits) and decimal (64, 128 bits) floating-point (FP) formats. Modern CPUs and frameworks (e.g. PyTorch) support binary machine numbers (FP32, FP64) for efficiency despite limited precision \cite{pytorchNumericalAccuracy}. Since Pyro is built on PyTorch \cite{pyroIntroductionPyro}, this research uses FP32 (Figure \ref{fig:floatingpointusedIllustrationown}), making it important to understand floating-point representation for analyzing the effects of Single Event Upsets (SEUs) on weight and bias parameters stored in FP32.

\iffalse
To give you an illustration of how bitflips can impact the performance of Neural Networks during inference, we will look back at Figure \ref{fig:zwxbIllustration}. Assume that we have the following set of input, weight, and bias:

\begin{flushright}
$$
\begin{aligned}
w = [0.2415, 0.5, 0.75] \\
x = [2.0, 3.0, 4.0] \\
b = 1 \\
\\
z = w_1 . x_1 + w_2 . x_2 + w_3 . x_3 + b \\
z = 0.2415 \cdot 2.0 + 0.5 \cdot 3.0 + 0.75 \cdot 4.0 \\
z = 4.983
\end{aligned}
$$
\end{flushright}

A bitflip in the first component of the exponent in $w_1$ representation, changing it from 1 to 0, will result in the value of 0.2415 changing to $8.22 \times 10^{37}$. The value of $z$ will become:

\begin{flushright}
$$
\begin{aligned}
z = w_1 . x_1 + w_2 . x_2 + w_3 . x_3 + b \\
z = 8.22 \times 10^{37} \cdot 2.0 + 0.5 \cdot 3.0 + 0.75 \cdot 4.0 \\
z = 16.44 \times 10^{37}
\end{aligned}
$$
\end{flushright}

The change of the value of $z$ from 4.983 to $16.44 \times 10^{37}$ is a significant change, which will propagate through the network. Although activation functions such as sigmoid and tanh can partially reduce the impact by limiting the output range of $h(\cdot)$ result, the injected error can still cause deviation to the subsequent layers. Consistency of output is threaten in this case, raising the need to look at the robustness of the neural networks in dealing with this example of SEU/bitflip.
\fi

\iffalse
TODO\_ContinueTheExplanation

There are types of single-event effects :
\begin{enumerate}
\item Single Event Upset (SEU)
\item Multi Bit Upset (MBU)
\item Real lab simulation of SEUs and how this experiment is just a 'miniaturized' version of it
\end{enumerate}


As mentioned in 
TODO\_include research from Bosio, Hari

TODO\_Give Illustration on z = wx+b on how the seu injection can impact the subsequent layers
\fi

\section{Robustness of Neural Networks}

% should I move this to the end of the chapter?

\subsection{Robustness against Adversarial Attack}

\iffalse
\begin{enumerate}
\item \cite{Blaas2021} Blaas 2021 
\item augmented robustness and the result of that research \cite{essbai2024a}
\item SEBR Robust Bayesian Neural Networks by Spectral Expectation Bound
Regularization : 2021
\item SGLD Bayesian Learning via Stochastic Gradient Langevin Dynamic: 2011
\item Carbone's research : June 2020
\end{enumerate}
\fi

Before we take a look into robustness of neural networks against SEUs/bitflips, we will see what are the defenses against adversarial attack applied in recent research, especially for Bayesian Neural Networks.

%\textbf{Stochastic Gradient Langevin Dynamics (SGLD)} did not address adversarial attack directly, but focuses on how deterministic model is overconfident to point estimate of weights used. SGLD was proposed in \cite{sgld2011} by adding noise during training while behaving similar to the Stochastic Gradient Descent (SGD). The final model predictions come from averaging many sampled parameter values. The robustness that is addressed in this paper is more related to how the model does not overfit and provide uncertainty, not reliability-related robustness. 

%In \cite{carbone2020a}, the author observed how Bayesian Neural Network is robust against \textbf{gradient-based adversarial attacks} e.g. Fast Gradient Sign Method (FGSM) and Projected Gradient Descent (PGD). This work explains that BNNs naturally mitigate gradient-based attacks because of the Bayesian posterior averaging.

%On the other hand, work in \cite{Blaas2021} claims the contrary, showing that BNNs trained with with accurate posterior inference (e.g. Hamiltonian Monte Carlo, advanced VI) \textbf{do not automatically gain adversarial robustness} compared to standard deterministic neural networks. This work showed an alternative result from \cite{carbone2020a}, in which the work shows that bayesian averaging does not wash out adversarial directions as strongly as the theory predicts. One of the proof is the non-vanishing gradients on the HMC sampling and susceptibility to PGD and FGSM attacks.

%\textbf{Spectral Expectation Bound Regularization (SEBR)} is introduced in \cite{sebr2021}. SEBR is a regularization technique added during training to limit the Lipschitz constant of BNN layers, which in the end improve robustness. It works against gradient-based attacks like FGSM and PGD by penalizing the expected spectral norm of each weight matrix (mean + variance), constraining how much the model output can change for a small input perturbations. The penalty term in SEBR is shown in $L_{\text{total}} = L_{\text{ELBO}} + \lambda \cdot L_{\text{SEBR}}$, in which $\lambda$ is the trade-off factor for the SEBR loss.


\textbf{Stochastic Gradient Langevin Dynamics (SGLD)} \cite{sgld2011} introduces noise in training to capture uncertainty, although it is not specifically designed for adversarial robustness. BNNs also have been reported to mitigate gradient-based attacks, e.g., Fast Gradient Sign Method (FGSM) and Projected Gradient Descent (PGD), due to posterior averaging \cite{carbone2020a}. \textbf{Spectral Expectation Bound Regularization (SEBR)} \cite{sebr2021} also increases robustness against gradient-based attacks by regularizing the spectral norm of weights. Contradicting work includes \cite{Blaas2021}, which shows that a more accurate posterior inference (e.g., Hamiltonian Monte Carlo, advanced VI) does not guarantee robustness.



%\begin{equation}
%L_{\text{total}} = L_{\text{ELBO}} + \lambda \cdot L_{\text{SEBR}}
%\label{eq:sebrpenalty}
%\end{equation}

\subsection{Robustness against Weight Perturbation}

%The term Single Event Upset Induced Parameter Perturbation (SIPP) is introduced in \cite{yan_2020a}. This paper reiterates SEU effect in flipping bits in memory or registers, especially in SRAM where model parameters are stored. It also shows that single bit-flip can severely degrade DNN accuracy, dropping accuracy from 70\% to 1\% in the experiment on using ResNet56 on CIFAR-100. This paper proposes hardware-level protection schemes : Triple Modular Redundancy (TMR), Error-Correcting Codes (ECC), and Selective Protection on top of TMR and ECC.

Single Event Upset Induced Parameter Perturbation (SIPP) \cite{yan_2020a} shows that SEUs in memory or parameters can drastically degrade accuracy (from 70\% to 1\%). Proposed defense in this work includes Triple Modular Redundancy (TMR), Error-Correcting Codes (ECC), and Selective Protection on both.

%The work in \cite{dennispoperobustcnn} take a different approach in developing robust machine learning systems against weight perturbation by focusing on model-intrinsic robustness instead of hardware fault tolerance. This paper proposes two new bounded activation functions, namely Weighted Gaussian (WG) and Rectified Weighted Gaussian (RWG) that limits the effect of flipped bits. This paper also introduce smartpooling that ignores unreasonable outlier activations above threshold. Finally, this paper also shown that dropout introduced during training helped maintain accuracy while a smaller effect is also shown by introducing weight decay during training.

Model-intrinsic robustness is explored in \cite{dennispoperobustcnn}, where Weighted Gaussian (WG) and Rectified Weighted Gaussian (RWG) activation functions, smart pooling, dropout, and weight decay show improvements in deterministic CNNs under SEU-like perturbations.

While both of the works mentioned previously take an interest in the robustness of deterministic neural networks, especially variants of convolutional neural networks, against weight/bias perturbation, \textbf{no significant research is found on observing Bayesian Neural Networks (BNNs) robustness against weight perturbation yet}. This area of research is what this paper aims to look into.

%That research focuses on replicating the weights and bias (like by a factor of 3 etc.). Methods that has been done focuses more on the replication etc., not strengthening the model architecture.



\section{Measure of Robustness} \label{section:robmeasuretheory}

To evaluate how well a neural network withstands adversarial attack or weight/bias perturbation, we need a quantitative metric that can measure its robustness. This measure will help us in comparing the variation of models or strategies with each other and decide which ones are the most robust against perturbations. Different alternatives of robustness, both for input and parameter perturbations, are as follows:

\begin{enumerate}
    \item Accuracy Delta

    %Accuracy delta is the most commonly used measure of robustness in this research. In research observing adversarial attacks, accuracy, which is the performance measure of the model in predicting the new data, is measured in at least two conditions, for clean data (clean accuracy) and data with adversarial attack (adversarial accuracy) \cite{zhang2021a}. In general, the accuracy delta is used to detect the change in accuracy before treatment and after treatment. A smaller accuracy delta build the perception of a more robust model, as it shows the stability of the model against the attack. Accuracy delta is used in multiple works, as in: \cite{dennispoperobustcnn}, \cite{feng2024attackingbayesadversarialrobustness}, \cite{pang2021evaluatingrobustnessbayesianneural}.

    Accuracy delta is the most commonly used measure of robustness in this field of research. It measures the performance difference before/after perturbation. Smaller values are perceived to show a better robustness \cite{dennispoperobustcnn}, \cite{feng2024attackingbayesadversarialrobustness}, \cite{pang2021evaluatingrobustnessbayesianneural}.

    \begin{equation} \label{eq:accuracydelta}
    \Delta \text{Accuracy} = \text{Accuracy}_{after} - \text{Accuracy}_{before}
    \end{equation}
    
    \item Difference in softmax

    %In the work of \cite{carbone2020a}, for $N$ test points, average difference in the softmax prediction between the output of the original data ($\mathbf{x}_j$), and the crafted adversarial input ($\tilde{\mathbf{x}}_j$) is notated by:

    Softmax Difference \cite{carbone2020a} captures prediction confidence change. Robustness is $1 - \text{softmax difference}$. Later work \cite{isrobustnessdongsu2018} highlights the trade-off between accuracy and robustness.

    \begin{equation} \label{eq:softmaxdifference}
    \text{Softmax Difference} = \frac{1}{N} \sum_{j=1}^{N} 
    \Big|
    \big\langle f(\mathbf{x}_j, \mathbf{w}) \big\rangle_{p(\mathbf{w}|D)}
    -
    \big\langle f(\tilde{\mathbf{x}}_j, \mathbf{w}) \big\rangle_{p(\mathbf{w}|D)}
    \Big|_{\infty}
    \end{equation}

    %In that work, robustness is also stated as (1 - softmax difference). While in \cite{isrobustnessdongsu2018} it is mentioned that robustness is a trade-off accuracy, work in \cite{carbone2020a} shown that BNNs trained with HMC became more accurate while maintaining robustness against FGSM attack. A smaller softmax difference makes a model to be deemed more robust, as it reflects less deviation in class confidence.
    
    %\item Augmented Robustness

    %Introduced in \cite{essbai2024a}, this measure of robustness includes the properties of Bayesian Neural Networks (BNNs) in its \textbf{uncertainty} towards its prediction. If the output of the network does not achieve a certain threshold, a new class that can be called as "unknown" is introduced to the result of the network. 

%Augmented robustness incorporates efficiency ($eff$), tolerance ($tol$), and penalization ($dep$). Efficiency ($eff$) quantifies the classifier's performance considering both accuracy and indecision (unknown predictions), balancing correct predictions and uncertainty. Tolerance ($tol$) rewards performance levels that meet or exceed a predefined acceptable threshold, while Penalization ($dep$) penalizes it. 

    %\begin{equation}
    %rob_{Aug}^{A}(C,D) = \frac{\int_{L_A}^{U_A} [tol(eff(C,D^{A_i}), \beta) - %dep(eff(C,D^{A_i}), \beta)] \cdot p_{A}(i)\,di}{2} + \frac{1}{2}
    %\end{equation}

    %Augmented robustness measures how the network can achieve high accuracy while minimizing indecision (the rate of unknown predictions due to uncertainty).
   
    
    %\item Sensitivity to SIPP
\end{enumerate}

While in general the measure of robustness is the same, some measure of robustness is exclusive to those aiming to minimize the impact of adversarial attacks. Other research includes a metric called Augmented Robustness, which was proposed in \cite{essbai2024a}, that takes into account the "unknown" class if the output of the network does not achieve a certain threshold confidence.


\iffalse
TODO\_in this section, I will explain multiple measure of robustness used in this field. I will write each measure of robustness and which author are including them in their research.

\begin{enumerate}
\item Accuracy delta : \cite{dennispoperobustcnn} (non adversarial attack). \cite{zhang2021a} also used Aleatoric Uncertainty and Epistemic Uncertainty.
\item Difference in softmax : \cite{carbone2020a}
\item Augmented robustness : \cite{essbai2024a}
\item Adversarial Global Accuracy (AGA): \cite{Blaas2021}
\item Sensivity to SIPP (SSIPP) : \cite{yan_2020a}, similar to accuracy delta, but focuses on the MAX across all bits.
\end{enumerate}

In this same subsection, I will also write the formula for any interesting measure of robustness, including the one that is offered in Augmented Robustness.

TODO\_add explanation on the mechanism of SSIPP, add the flowchart if you have to. Emphasize your statement in a way that shows that SSIPP focus on flipping all bits and get the max of it.
\fi


\iffalse
\section{Method of increasing robustness in neural networks}

\subsection{Pooling Layer (Architecture)}

TODO\_William's research suggest smartpooling from \cite{santos_2019a}. 
TODO\_Give the algorithm on how smartpooling works

\subsection{Activation Function (Architecture)}

TODO\_William's research suggests WG and RWG(0.01)

\subsection{Weight Decay (Training Process)}

TODO\_William's research suggests weight decay of 0 and 0.0001

\subsection{Dropout (Training Process)}

TODO\_William's research suggests dropout rate of 0 and 0.5

\subsection{Work on SIPP}

\cite{yan_2020a} suggests Triple Modular Redundancy (all parameters protected) to protect parameters in SRAM. Error-correcting code with Hamming code is used to make the area overhead redundancy exponentially smaller (only 3.5\% of the area overhead could be achieved at a cost of approximately 3.5× of protection logic)

TODO\_understand more about SRAM and PROTECTION LOGIC difference.

\cite{yan_2020a}

\section{Method of increasing robustness in bayesian neural networks against adversarial attacks}

TODO\_CHECK blaas, SEBR etc.
TODO\_add claims from Y.Gal and Carbone that any gradient-based adversarial attacks are invalid against BNN, in an extremely idealized conditions.

\begin{enumerate}
\item Spectral Expectation Bound Regularization (SEBR) : \cite{zhang2021a} TODO\_explain how it work, which is still a type of regularization in SVI or VI. also mention a little bit about Lipschitz constant.
\end{enumerate}
\fi

% -----------------------------------------------------------------------------

\chapter{Design}
\label{chap:design}

\iffalse
{\bf A topic-specific chapter, roughly 30\% of the total page-count. Blackboard suggests 10 pages, including chapter implementation.}
\fi
\vspace{1cm} 
\iffalse

\section{Example Section}

This is an example section; 
the following content is auto-generated dummy text.
\lipsum

\subsection{Example Sub-section}

\begin{figure}[t]
\centering
foo
\caption{This is an example figure.}
\label{fig}
\end{figure}

\begin{table}[t]
\centering
\begin{tabular}{|cc|c|}
\hline
foo      & bar      & baz      \\
\hline
$0     $ & $0     $ & $0     $ \\
$1     $ & $1     $ & $1     $ \\
$\vdots$ & $\vdots$ & $\vdots$ \\
$9     $ & $9     $ & $9     $ \\
\hline
\end{tabular}
\caption{This is an example table.}
\label{tab}
\end{table}

\begin{algorithm}[t]
\For{$i=0$ {\bf upto} $n$}{
  $t_i \leftarrow 0$\;
}
\caption{This is an example algorithm.}
\label{alg}
\end{algorithm}

\begin{lstlisting}[float={t},caption={This is an example listing.},label={lst},language=C]
for( i = 0; i < n; i++ ) {
  t[ i ] = 0;
}
\end{lstlisting}

This is an example sub-section;
the following content is auto-generated dummy text.
Notice the examples in Figure~\ref{fig}, Table~\ref{tab}, Algorithm~\ref{alg}
and Listing~\ref{lst}.
\lipsum

\subsubsection{Example Sub-sub-section}

This is an example sub-sub-section;
the following content is auto-generated dummy text.
\lipsum

\paragraph{Example paragraph.}

This is an example paragraph; note the trailing full-stop in the title,
which is intended to ensure it does not run into the text.

\fi

\iffalse
\section{Overview of Our Approach} \label{section:approach}

TODO\_MIGHT AS WEL INCLUDE THE flowchart of our work from start to end
\fi

\section{Dataset Description} \label{section:designdataset}

In this research, we are using an image dataset called "Ships in Satellite Imagery" \cite{robert_hammell_2018}, or "Shipsnet" in short. The dataset is satellite imagery captured from a commercial company called Planet. Those satellite images are taken from over the San Francisco Bay and San Pedro Bay areas of California. Initially, the pictures are orthorectified 3-meter pixel size, but for this task in detecting ships, 4000 80 x 80 RGB images are derived from it, labeled with either a "ship" (1) or "no-ship" (0) classification. Both classes are explained below:

\begin{enumerate}
\item Ship: This class includes 1000 images. To be qualified to be included in this class, ships are supposed to be centered in the image. The images consist of ships of different sizes, orientations, and atmospheric collection conditions.
\item No-ship: This class includes 3000 images. 
    \begin{enumerate}
    \item Land cover features, randomly sampled: buildings, bare earth, vegetation, water
    \item Images that only cover a portion of ships: these images, although they contain part of ships, do not qualify to be categorized as "ships"
    \item Images that were previously mislabeled by machine learning models, typically caused by bright pixels or strong linear features
    \end{enumerate}
\end{enumerate}


\begin{figure}[h]
\caption{Sample images from the Shipsnet dataset.}
\includegraphics[scale=0.4]{shipsnet_samples.png}
\centering
\label{fig:shipsnetsamples}
\end{figure}



\section{Deterministic CNN Model Architecture}  \label{section:designdeterministic}

\iffalse
\begin{enumerate}
\item each layer used
\item how big is the dimension, input output
\item any dropouts?
\item what are the activation functions used?
\end{enumerate}
\fi

The deterministic CNN model that we will set as a benchmark for this research is designed for binary image classification as observed from section \ref{section:designdataset}. Here are the layers used, along with the input and output shape for each layer. 


\begin{table}[h]
\centering
\begin{tabular}{|c|c|c|c|c|}
\hline
Layer      & Input Shape      & Output Shape & Trainable Weight & Trainable Bias      \\
\hline
Conv1 + Pooling & [B,3,64,64] & [B,32,32,32] & 864 &32\\
\hline
Conv2 + Pooling & [B,32,32,32] & [B,64,16,16] &18,432&64\\
\hline
Flatten + Pooling & [B,64,16,16] & [B,16384] &32,768&2\\
\hline
Linear (Output) & [B,16384] & [B,2] &-&-\\
\hline
\end{tabular}
\caption{Layers in our model}
\label{table:deterministicdesign}
\end{table}

Conv1 and Conv2 represent the convolutional layer (torch.nn.Conv2d), followed by an activation function and a max pooling (torch.nn.MaxPool2d), then flattened to be processed in a fully connected layer. The training process of this model will use the Adam optimizer with a learning rate of $10^{-3}$ and a cross-entropy loss function. In this deterministic implementation of our model, with \textbf{52,162 trainable weights and biases}, do note that all the weights and biases for each neuron are fixed after the training is done. The model output is logits of the two classes, and the chosen class is the argmax of these logits, which is the class with the highest logits among all.

\section{Bayesian Implementation of the Deterministic CNN Model} \label{subsection:bayesianimplementationdist}

\iffalse
\begin{enumerate}
\item Distributions used and the chosen parameter of gaussian(0,10). laplace(0,10) and uniform(-10,10). I should also explain about each of the distribution in the background.
\item AutoNormal and custom guide that I am using, and how the monte carlo simulation of this guide shows the intended distribution
\end{enumerate}
\fi

The Bayesian CNN model designed in this experiment preserves the core structure of the previously mentioned deterministic CNN in subsection \ref{section:designdeterministic}. The first thing that is different from the deterministic CNN is how each weight and bias parameter is modeled as a random variable drawn from a prior distribution. As already mentioned in subsection \ref{section:backgroundpriordist}, we will be using the Gaussian prior, Laplace prior, and Uniform prior with parameters $a$ and $b$.

Bayesian Inference for these random variables will be done with \textbf{Stochastic Variational Inference (SVI)}. The variational family used in this experiment is the \textbf{mean-field assumption}. In this type of variational family, all weights and biases are assumed to be independent.


%TODO : Consider adding the concept of "Variational Family" to the background. explain some other type of it like the normalizing flow and the research behind it etc.



\section{Training Process}

\iffalse
In this section, I will explain about SVI that's done in the pyro:
\begin{enumerate}
\item what's the loss, (ELBO) configured (num\_samples)
\item what's the optimizer, do I use weight decay
\item what's the learning rate
\item what's the batch size
\end{enumerate}
\fi

The loss function used to train the Bayesian CNN model is the negative Evidence Lower Bound (ELBO). While the generic version of it mentioned in Equation \ref{eq:genericelbo} is used in generic Variational Inference, the BNN-specific ELBO that we will use in this experiment is defined as:

\begin{equation}
\textbf{ELBO} \;\equiv\; 
\mathbb{E}_{q_\phi(W)}\big[\log p(D \mid W)\big]
\;-\;
\text{KL}\big(q_\phi(W) \;\|\; p(W)\big)
\label{eq:elbo_bnn}
\end{equation}

\iffalse
\begin{equation}
\log p_\theta(x, z) = \log p(x \mid z) + \log p(z)
\end{equation}
\fi


Given the dataset $D$ with observed inputs and labels, the ELBO loss aims to maximize the expected log-likelihood of the data from the sampled network weights while minimizing the KL divergence between the learned variational distribution $q(W)$ and the prior $p(W)$, trying to balance the data fit and regularization. This exact formula is the loss function that can be found in the Pyro package in Python, specifically in \texttt{pyro.infer.trace\_elbo}.

The optimizer used in this experiment is \textbf{ClippedAdam}, taken from the Pyro Python package, with a learning rate of $10^{-3}$. The default setting of weight decay is 0, but it can be adjusted for the purpose of understanding the impact on the overall robustness.


\section{Inference Process}

\iffalse
\begin{algorithm}[t]
\For{$i=0$ {\bf upto} $n$}{
  $t_i \leftarrow 0$\;
}
\caption{This is an example algorithm.}
\label{alg}
\end{algorithm}

In this section, I will tell that:
\begin {enumerate}
\item the inference process for image is done 10 times
\item get the logits for each class
\item average the logits (TODO: find out the literature that motivates this)
\item class is the argmax of the logits mean
\end {enumerate}
\fi

Inference for the Bayesian CNN model, in this experiment, is done in a Monte Carlo Sampling fashion, as shown in Algorithm \ref{alg:mc_inference}.

\begin{algorithm}[h]
\caption{Monte Carlo Sampling for Bayesian CNN Inference}
\label{alg:mc_inference}
\KwIn{Trained Bayesian model $f$, guide $q_\phi(W)$, data loader $\mathcal{D}$, number of samples $S=10$}
\KwOut{Predicted class labels $\hat{y}$ for all inputs}
\BlankLine

\ForEach{mini-batch $(x, y)$ in $\mathcal{D}$}{
    Initialize logits accumulator: $L \leftarrow 0$\;
    \For{$s = 1$ {\bf to} $S$}{
        Sample weights: $W^{(s)} \sim q_\phi(W)$\;
        Compute logits: $\ell^{(s)} \leftarrow f(x; W^{(s)})$\;
        Accumulate: $L \leftarrow L + \ell^{(s)}$\;
    }
    Compute mean logits: $\bar{L} = \frac{L}{S}$\;
    Predicted class: $\hat{y} = \arg\max(\bar{L})$\;
    Store $\hat{y}$\;
}
\Return all predicted labels $\hat{y}$\;
\label{alg:msamplingbnninference}
\end{algorithm}

With the algorithm below, the default Monte Carlo samples per batch is set to be equal to 10. The number is arbitrary and can be adjusted depending on the needs, with a higher number will take from more posterior weight samples. Each iteration samples new weights ($W^{(s)}$) from a variational posterior family. Logits are averaged before applying argmax, reducing prediction variance. This mechanism shows the aggregation of Bayesian predictive uncertainty, which can not be found in the deterministic CNN.

%TODO add explanation to the implementation chapter that actually, SEU will hit all the S=10 monte carlo sampling. In the analysis, tell the readers that actually if the hit only happen in one of the S=10, the robustness will be better.

\section{Activation Function Variations}

\begin{figure}[h]
\caption{Activation Functions Plot}
\includegraphics[scale=0.5]{activation_functions_new.png}
\centering
\label{fig:activationfunctions}
\end{figure}

In this research, we will vary the activation functions to form different model variations. We have 7 activation functions that we want to compare in terms of their performance, which some widely known activation functions: ReLU, Sigmoid, Tanh, and Sine. We also want to test the Weighted Gaussian (WG) and Rectified Weighted Gaussian introduced in \cite{dennispoperobustcnn} in this experiment, which are shown in Equation (\ref{eq:wg}) and Equation (\ref{eq:rwg}). As previous research shows that ReLU is the least robust activation function, we also want to look at a different version of it called ReLU6, which is shown in Equation (\ref{eq:relu6}). The plots of each activation function are shown in Figure \ref{fig:activationfunctions}.

%TODO explain WG and RWG formula
%TODO add relu 6
%TODO make the graph to be your own, possibly merging them into one or group them : (WG RWG, relu relu 6, and the rest)

\noindent
\begin{minipage}{0.48\linewidth}
\begin{equation}
\text{WG}(x) = x e^{-x^{2}}
\label{eq:wg}
\end{equation}
\end{minipage}\hfill
\begin{minipage}{0.48\linewidth}
\begin{equation}
\text{RWG}(x) = \max(0, x e^{-x^{2}})
\label{eq:rwg}
\end{equation}
\end{minipage}

\noindent
\begin{minipage}{0.48\linewidth}
\begin{equation}
\text{ReLU6}(x) = \min(\max(0, x), 6)
\label{eq:relu6}
\end{equation}
\end{minipage}\hfill
\begin{minipage}{0.48\linewidth}
%begin{equation}
%\text{RWG}(x) = \max(0, x e^{-x^{2}})
%\label{eq:rwg}
%\end{equation}
\end{minipage}


\section{Prior Distribution Parameter Variations}

In this research, for each prior distribution that we want to test, as mentioned in sub-section \ref{subsection:bayesianimplementationdist}, we will vary the parameter of the distribution to see the effect on model accuracy and SEU robustness. We will fix the value of $a$ ($\mu$ or $lows$) in our prior to be equal to 0, as we have no strong belief about the exact value that best fits the model, but we still want to have a prior distribution that is pleasantly centered. We will vary the value of $b$ with the following:
\begin{enumerate}[noitemsep, topsep=0pt]
    \item $0.1$: very narrow prior that shows we are supposedly confident that the weights are deviating only slightly from 0. As the weight is kept in a small window, this discourages flexibility in the model weight, which leads to underfitting.
    \item $1$: a moderate prior that is commonly used. This value is not expected to be underfitting nor overfitting.
    \item $10$: very wide prior that shows we have weak belief towards the true value. The model is too free to fit the data, which poses the risk of overfitting.
\end{enumerate}

%With the above-mentioned values of $b$, we would like to observe the learned distributions of weights and biases after training, and how closely they align with their respective priors.

\section{Model Variations} \label{subsection:modelvariations}

To observe different model architectures and their impact to the robustness of our Bayesian CNN model, we will explore a few options that we can choose to apply to that model, on top of our base model. 

\begin{enumerate}[noitemsep, topsep=0pt]
\item SmartPooling

In \cite{santos_2019a}, the author evaluates how single radiation-induced soft errors in GPUs could propagate through CNN layers and impact detection accuracy. They propose mechanisms like ECC (error-correcting codes) and ABFT (algorithm-based fault tolerance) for matrix multiplication operations. In the same work, they also propose a modified maxpooling layer called "\textbf{Reliable Maxpooling}". This layer detects unreasonable values, typically chosen in maxpooling layers, and either halts frame processing or propagates the second-largest value to reduce error propagation.

Examination of how reliable maxpooling is working in a weight perturbed condition, \cite{dennispoperobustcnn} implemented reliable maxpooling in the name of SmartPooling, in which the proposed layer can be a useful mitigation for unbounded activation functions like ReLU. Their work shows that 16\% robustness is achieved in Relu-based deterministic CNNs equipped with SmartPooling.

\begin{algorithm}[h]
\caption{SmartPool: Reliable Max-Pooling with Outlier Detection and Correction}
\KwIn{
Input tensor $X \in \mathbb{R}^{N \times C \times H \times W}$; \\
Kernel size $k$; Stride $s$; Threshold $\theta$; \\
}
\KwOut{Pooled tensor $Y \in \mathbb{R}^{N \times C \times H' \times W'}$}
\BlankLine
\textbf{1.} Extract all sliding windows (patches) of size $k \times k$ from $X$ with stride $s$, producing:
\[
P = \text{Unfold}(X) \quad \text{where}\; P \in \mathbb{R}^{N \times C \times k^2 \times L}
\]
where $L$ is the number of windows per feature map.\;

\textbf{2.} For each patch $P_{n,c}$, compute the top-2 largest values:
\[
(\text{max1}, \text{max2}) = \text{TopK}(P_{n,c}, k=2)
\]
where $\text{max1}$ is the largest and $\text{max2}$ is the second-largest value in each window.\;

\textbf{3.} Detect abnormal activations (spikes):
\[
S_{n,c,l} =
\begin{cases}
1 & \text{if } \text{max1}_{n,c,l} > \theta,\\
0 & \text{otherwise}
\end{cases}
\]
for every batch $n$, channel $c$, and window index $l$.\;

\textbf{4.} If the algorithm is set to only see the spiked values without altering the output, set:
\[
Y_{n,c,l} = \text{max1}_{n,c,l}
\]
If the reliable maxpooling is used to overwrite the spiked value, replace detected spikes with the second-largest value:
\[
Y_{n,c,l} =
\begin{cases}
\text{max2}_{n,c,l}, & S_{n,c,l}=1 \\
\text{max1}_{n,c,l}, & S_{n,c,l}=0
\end{cases}
\]
\textbf{5.} Reshape $Y$ back to spatial dimensions:
\[
Y \in \mathbb{R}^{N \times C \times H' \times W'}, \quad
H' = \frac{H}{k},\quad W' = \frac{W}{k}
\]
\KwRet{$Y$}
\label{alg:smartpooling}
\end{algorithm}

In this research, we typically set the value of Threshold $\theta=10$. In activating this SmartPooling, we replace every MaxPooling that was applied in our paper with it. The smartpooling algorithm is shown in Algorithm \ref{alg:smartpooling}.

%TODO\_give an illustration how the smartpooling works 

\item Dropout layer

Work in \cite{dennispoperobustcnn} also explored how the dropout layer can benefit from mitigating the SEU-induced performance drop. In this work, the author explored two dropout probability $p$ values, which are $p=0$ (off state) and $p=0.5$. (on state). Dropout layer is activated during training, which hopefully will help the model regularize against SEU by deactivating random neurons during training, so the network cannot rely too much on any single pathway, forcing redundancy and better generalization.

%TODO\_ILLUSTRATION OF DROPOUT LAYERS

\item Weight decay

Weight decay is a common regularization technique in CNNs to prevent overfitting. Regularization allows complex models to be trained on data sets of limited size without severe over-fitting \cite{bishoppatternrecognition2007}. The motivation in modifying the weight decay to inspect the impact of SEU/bitflips lies in the ability of weight decay to constrain the weights, hopefully reducing the impact of bitflips on any of the weights/bias. In our deterministic CNN model, we will use the Adam optimizer from PyTorch \cite{kingma2014adam}, and the Bayesian implementation will use the ClippedAdam optimizer from Pyro \cite{computation11050095}. Updated weight parameter, considering weight decay, is formulated in the equation below:

\begin{equation}
\label{eq:clippedadam_update}
w_{t+1} = w_t - \eta \Big( \nabla L(w_t) + \lambda \, w_t \Big)
\end{equation}

In that formula, $w_t$ represents the weight parameter at step $t$, $\eta$ is the learning rate, $\nabla L(w_t)$ is the gradient of the loss with respect to the weights, and $\lambda$ is the weight decay coefficient. In this research, weight decay is set in two alternative values, which are $\lambda=0$ and $\lambda = 10^{-4}$.
\end{enumerate}

These model variations, added with the base model, will make a total of 4 model variations, which will be explained more in Subsection \ref{subsection:trainingdeterministiccnn}.

\section{SEU process}

\iffalse
In this section, I will tell about:
\begin{enumerate}
\item SEU process in this section is a simulated one (done on software level)
\item Among a lot of SEE, I will try to simulate SEU only, not anything like MBU etc.
\item bitflips is done to the fp32 representation
\item which bit index to choose from?
\item which weight/bias index that I want to attack? is it all? is it only beginning and end? is it necessary to attack all?

\end{enumerate}
\fi

The process of simulating SEU in this research is done with the help of Python's \texttt{struct} module. Here are some rules that are imposed during the simulation of SEU injection:
\begin{enumerate}
\item SEU is done before the weight/bias is sampled. Thus, the SEU process is attacking the parameter of either $a$ or $b$ from the guide $q_\phi(W)$.  The weight/bias is sampled from the newly perturbed guide $q'_\phi(W)$. Example of sampled weight distribution before and after SEU is shown in Figure \ref{fig:sampledweight}
\item Single Event Upset/bitflip is attacking only one of the bit locations within the 32-bit IEEE 754 representation at a time. 
\item SEU is bitflipping a 0 value to 1, or 1 value to 0. Hence, it's bitflipping the value into another value regardless of the starting value of the bit.
\item SEU is only done in a single neuron's weight/bias at a time.
\item SEU might change the value of any of the distribution of the parameter into a different sign (+ to -, or vice versa). The value of parameter $b$, which is the variance of the Gaussian/Laplace distribution, and the width of the Uniform distribution is not expected to be negative. In this experiment, we assume that we does not enforce any constraint to keep the value $b$ positive after the SEU, that may make the model unable to return a valid prediction.
\end{enumerate}

In this research, for our deterministic CNN models, we will do 2,352 SEU injections. This SEU injections are the combination of 28 different models that we have (4 model variations and 7 activation functions), with each model get 84 different SEU injections, which comprises a combination of 2 neuron attack locations (first/last neuron), 3 layer (conv1/conv2/fc1), 2 layer parameter (weight/bias), and 7 bit index. 

For Bayesian CNN models, 28 different models will be attacked with 468 different SEU injections, 252 ($7 \times 36$) for the prior center and 216 ($6 \times 36$) for the prior spread, for a total of 13,104 SEU injections. The attack comprises of 36 different attacks, which combination comprises of 3 prior (Gaussian, Laplace, Uniform), 2 neuron attack locations, 3 layers and 2 layer parameters. We can only attack 6 variations of bit index into the prior spread, as hitting the bit index 0 for prior spread will change the prior spread parameter sign from + to -, which is the limitation of our prior as mentioned in Subsection \ref{section:backgroundpriordist}.

\section{Robustness Measure}

The robustness measure that we will use in this research will be based on the two metrics from Section \ref{section:robmeasuretheory} , which are Accuracy Delta and Softmax Difference. Accuracy of the model after the SEU injection does not necessarily lower than the accuracy before the SEU injection is done, and we want to look both for positive or negative change of accuracy and penalize them equally, without cancelling each other. To do that, we will modify the formula of Accuracy Delta shown in Equation \ref{eq:accuracydelta} to be the one in Equation \ref{eq:absaccdiff} below.

%TODO explain the formula for each. Accuracy delta will use absolute difference. Softmax will be done for the logits.

    \begin{equation} \label{eq:absaccdiff}
    \text{Absolute Accuracy Difference} = \left|\text{Accuracy}_{after} - \text{Accuracy}_{before} \right|
    \end{equation}


Our research questions rely on comparing multiple model variations and deciding which one is the best based on the given result. While both robustness measures are useful in achieving our goals on its own, interpreting both of the robustness measure at the same time may result in confusion, e.g., one model might have a low absolute accuracy difference, but a high softmax difference. To make the model comparison easier, we devise a new aggregate robustness measure called the Aggregate Robustness Index (ARIn), which takes Equation (\ref{eq:absaccdiff}) and Equation (\ref{eq:softmaxdifference}).

\begin{equation}
%\text{ARIn} = \sqrt{\frac{\left(\tilde{\Delta}_{\text{acc}}\right)^2 + \left(\tilde{\Delta}_{\text{softmax}}\right)^2}{2}}
\text{ARIn} = \sqrt{\frac{\left(\text{Absolute Accuracy Difference}\right)^2 + \left({\text{Softmax Difference}}\right)^2}{2}}
\label{eq:ARIn}
\end{equation}

%TODO Explain how it works. Does it penalize high value or what. Make link to the RMS function.


% -----------------------------------------------------------------------------

\chapter{Implementation}
\label{chap:implementation}

\iffalse
{\bf A topic-specific chapter, roughly 30\% of the total page-count. Blackboard suggests 10 pages, including chapter design.}
\vspace{1cm} 

\noindent
This chapter is intended to describe what you did: the goal is to explain
the main activity or activities, of any type, which constituted your work 
during the project.  The content is highly topic-specific. For some 
projects it will make sense to split the content into two main sections, or maybe even into two separate chapters: one 
will discuss the design of something, including any rationale or decisions made, 
and the other will discuss how this design was realised via some form of 
implementation.  You could instead give this chapter the title ``Design and Implementation"; or you might split this content into two chapters, one titled ``Design" and the other ``Implementation" .

Note that it is common to include evidence of ``best practice'' project 
management (e.g., use of version control, choice of programming language 
and so on).  Rather than simply a rote list, make sure any such content 
is useful and/or informative in some way: for example, if there was a 
decision to be made then explain the trade-offs and implications 
involved.
\fi

\section{Experimental Environment}

For this project, we use \texttt{Python 3.11} as the main programming language to do every tasks, from preparing the dataset, training and testing deterministic and Bayesian CNN models, up to the simulation of SEU/bitflips and its analysis. \texttt{Anaconda 24.7.1} is used in this project to provide package and virtual environment management. For deep learning tasks, especially the deterministic CNN models, we use \texttt{PyTorch 2.7.1}, \texttt{torchaudio 2.7.1}, and \texttt{torchvision 0.22.1}. All of these modules are installed with support for \texttt{Cuda 12.8}, enabling GPU acceleration.

For the implementation of Bayesian CNN models, we chose \texttt{Pyro 1.9.1}. Pyro is a universal probabilistic programming language (PPL) built in Python and powered by PyTorch as its backend \cite{bingham2018pyro}. It provides a flexible and expressive framework for constructing deep probabilistic models, combining the strengths of modern deep learning with Bayesian inference techniques. GitHub is used in this project for the purpose of code versioning.



\section{Computation Resource}

This project is run on the Elastic Compute Cloud (EC2) instance on Amazon Web Services (AWS), with credits provided by the University of Bristol. Specifically, this project is using the G4dn EC2 instance family. This instance family is equipped with NVIDIA T4 GPU and custom Intel Cascade Lake CPUs \cite{amazonAmazonInstances}. The instance type chosen is the G4dn.xlarge, with specification as follows \cite{amazonCloudCompute}: 1 GPU, 16 GiB of GPU memory, 4 vCPU, 16 GiB memory, and 1 x 125 GB NVMe SSD storage. The instance is operated on the Ubuntu 22.04 operating system on top of AWS Deep Learning Base GPU AMI. 

\section{Dataset Preparation}

The dataset mentioned in Section \ref{section:designdataset} is then preprocessed with \texttt{torchvision.transforms} before we train them to our models:
\begin{enumerate}
\item Resize : Image is resized from 80 x 80 pixels to 64 x 64 pixels. This will help us in fewer computation needed, and convenience of using image with a power of two that will make it easier for downsampling with pooling layers.
\item Scaling : From the image with RGB channels that have a value of intensity between 0 and 255, the value is converted to be ranged from 0 to 1.
\item Standardization : For each channel, we calculate the mean and standard deviation of pixel density, then normalize each values with:

\begin{equation}
x_{\text{norm}} = \frac{x - \mu}{\sigma}
\end{equation}

This will ensure that the value of each pixel ranges from [-1,1] and centered around 0. This will help our deterministic and Bayesian CNN models in learning, especially as optimizers assumes the input to be centered and scaled similarly. 

%TODO\_Might also include the before and after of the transformation

\end{enumerate}

\section{Training of Deterministic CNN Model} \label{subsection:trainingdeterministiccnn}

Model is trained in 20 epochs and saved each time best validation accuracy is reached. Model names are coded as follows:
\begin{enumerate}
\item Model 0 : This is the base model as mentioned in Section \ref{section:designdeterministic}
\item Model 1 : Instead of using \texttt{torch.nn.MaxPool2d}, this model uses smartpooling, as mentioned in subsection \ref{subsection:modelvariations} point 1.
\item Model 2 : Dropout layer (\texttt{torch.nn.Dropout}) is applied after conv2 layer with $p=0.5$
\item Model 3 : Weight decay is implemented to the \texttt{Adam} optimizer with $\lambda = 10^{-4}$
\end {enumerate}

\section{Training of Bayesian CNN Model}

The model trained in the Bayesian CNN setting have the similar architecture to the one in the deterministic setting. Model is trained in 100 epochs, model was saved each time best train accuracy is reached. Model names follows the nomenclature of the models trained in the deterministic CNN model (subsection \ref{subsection:trainingdeterministiccnn}). Model is trained on three different $b$ parameter : $b = [0.1, 1, 10]$.


\section{Implementations in Pyro}
\iffalse
TODO\_In this section, I will tell about adjustments made in Pyro
\begin{enumerate}
\item Explain about the usage of PyroModule and PyroSample, and how I utilize function make\_prior
\item Guide AutoNormal, AutoLaplace, AutoUniform. How for AutoLaplace and AutoUniform, you have to tweak it a bit for this research.
\item explain the SVI training process of the code. What is the output of the training, and how the training checkpoints are saved. (best train accuracy for every 10 epochs, first epoch, and last epoch)
\item how the guide (taking from the pyro param store) updated because of the training
\end{enumerate}
\fi

Here are some adjustments made in the Bayesian CNN models in Python by utilizing \texttt{Pyro}:
\begin{enumerate}
\item \texttt{pyro.nn.module.PyroModule} is used to embed probabilistic constructs (sampling, parameters with priors, constraints) directly into the \texttt{torch.nn.Module} paradigm. In our case, it transforms \texttt{torch.nn.conv2d} and \texttt{torch.nn.Linear} into Bayesian layers, by enabling their weights and biases to be treated as random variables rather than fixed variables.
\item \texttt{pyro.nn.module.PyroSample} is used as a descriptor inside PyroModule to define random variables that will be sampled from a prior distribution. For each layer's weight and bias we define each from what prior distribution is used for each. In this research, we define 3 distribution types : Gaussian, Laplace, and Uniform inside this descriptor, along with each parameter of $a$ and $b$.
\item The \textbf{BNN model} and its prior distribution has been created by the two classes above. For Stochastic Variational Inference (SVI) to work, we need the parameterized approximation to the posterior distribution over the parameters sampled from \texttt{pyro.nn.module.PyroSample}, which is the \textbf{guide}. To automatically initialize the guide, we use \texttt{pyro.infer.autoguide.guides.AutoNormal}. This automatic guide creates a separate Normal/Gaussian distribution (with learnable mean and scale/variance) for each PyroSample site. In this implementation, each weight and bias is assumed to be \textbf{independent} with each other.
\item While \texttt{pyro.infer.autoguide.guides.AutoNormal} can be used for the Gaussian Prior, it is not suitable for the Laplace and Gaussian prior, as we want to design our research to have a posterior approximation similar to the shape of our prior. Pyro implementation for automatic guide for Laplace and Uniform is not provided by the Pyro module; therefore, in this research, we modify \texttt{AutoNormal} to match the Laplace and Uniform distribution. In this research, we are naming them \texttt{AutoLaplace} and \texttt{AutoUniform}.
\item \texttt{pyro.infer.svi.step} is used to do the training step in SVI for our Bayesian CNN model instead of \texttt{torch.optim.Adam.step} in deterministic CNN. The loss value returned is the negative ELBO for that batch.
\item Learned posterior from Autonormal, e.g. \texttt{Autonormal.locs.conv1.weight}, is stored inside an object called \texttt{pyro.get\_param\_store()}. To ensure no parameter leak is happening in-between experiments, \texttt{pyro.clear\_param\_store()} is used to clear all global parameters.
\item \texttt{pyro.poutine.trace(guide).get\_trace(images)} samples the weight from the guide distribution, and saving the sampled deterministic value to a \texttt{trace}.
\item \texttt{pyro.poutine.replay(model, trace)} is then used to inject the sampled weights from the earlier to the model. This makes the model uses them deterministically in the forward pass process.
\end{enumerate}


\iffalse
\section{Dataset Description}

TODO, In this section, I will explain about the dataset used in this research:
\begin{enumerate}
\item The name of the Dataset, where did I get it from, and the licenses for usage (creative commons etc.) SHIPSNET. Maybe some motivation in my decision on choosing them.
\item Sample Images per class, along with the image dimension
\item Pixel numbers in RGB, and how I normalize it. Give before normalization and after normalization
\item Train Test Split
\end{enumerate}


\section{Evaluation Metrics}
TODO, in this section, I will explain the evaluation metrics used in this research:
\begin{enumerate}
\item Explain about what is normally used in evaluating NN models (accuracy)
\item Explain about Softmax Difference, and how important it is to capture our notion of robustness
\item Explain about delta accuracy
\end{enumerate}
\fi




% -----------------------------------------------------------------------------

\chapter{Critical Evaluation}
\label{chap:evaluation}

\iffalse
{\bf A topic-specific chapter, roughly 30\% of the total page-count. Blackboard suggests 10 pages.} 
\vspace{1cm} 

\noindent
This chapter is intended to evaluate what you did.  The content is highly 
topic-specific, but for many projects will have flavours of the following:

\begin{enumerate}
\item functional  testing, including analysis and explanation of failure 
      cases,
\item behavioural testing, often including analysis of any results that 
      draw some form of conclusion wrt. the aims and objectives,
      and
\item evaluation of options and decisions within the project, and/or a
      comparison with alternatives.
\end{enumerate}

\noindent
This chapter often acts to differentiate project quality: even if the work
completed is of a high technical quality, critical yet objective evaluation 
and comparison of the outcomes is crucial.  In essence, the reader wants to
learn something, so the worst examples amount to simple statements of fact 
(e.g., ``graph X shows the result is Y''); the best examples are analytical 
and exploratory (e.g., ``graph X shows the result is Y, which means Z; this 
contradicts [1], which may be because I use a different assumption'').  As 
such, both positive {\em and}\/ negative outcomes are valid {\em if} presented 
in a suitable manner.
\fi


\section{Analysis of Model Performance Metric}

%This includes test accuracy before SEU. This can also include how the training process going by explaining the LOSS per epoch, val accuracy, and how the best model is chosen.

In this section, we will monitor how each model, before SEU injection, performs on a given test set. The performance metric used in this experiment is test accuracy, which was mentioned in Section \ref{section:robmeasuretheory}. For deterministic CNN model, we will compare all model variants across all activation functions, compare the results, and decide which model variant and activation function are the best compared to others. For the Bayesian CNN model, in addition to that, we will also take a look at prior $b$ parameter variations, and see which one yields in a better result.

\subsection{Training Result of Deterministic CNN}

\iffalse
In this section, I will explain:
\begin{enumerate}
\item The training loss for each epoch for each activation function
\item Compare train, val, test accuracy
\item Spot if anything is weird
\end{enumerate}
\fi


\iffalse
\begin{table}[h]
\centering
\begin{tabular}{|c|c|c|c|c|c|c|c|c|}
\hline
Model Variant & RWG & WG & relu & relu6 & sigmoid & sin & tanh & \textbf{average} \\
\hline
Model 0 & 0.9938 & 0.9844 & 0.9906 & \textbf{0.9906} & \textbf{0.9781} & 0.9844 & 0.9875 & 0.9871\\
Model 1 & \textbf{0.9969} & 0.9812 & 0.9875 & \textbf{0.9906} & \textbf{0.9781} & \textbf{0.9875} & 0.9812 & 0.9861\\
Model 2 & \textbf{0.9969} & \textbf{0.9875} & \textbf{0.9906} & \textbf{0.9906} & 0.9719 & 0.9906 & \textbf{0.9938} & \textbf{0.9888}\\
Model 3 & \textbf{0.9969} & 0.9844 & 0.9875 & 0.9875 & 0.9750 & 0.9812 & 0.9844 & 0.9853\\
\hline
\end{tabular}
\caption{Validation accuracy of different activation functions across model variants in training Deterministic CNN Models}
\label{table:activation_performance_deterministic_train}
\end{table}

%TODO Analysis of Training Result

Based on Table \ref{table:activation_performance_deterministic_train}, we can see that the Rectified Weighted Gaussian activation function results in the highest validation accuracy among all activation functions. Model 2 is observed to result in a better validation accuracy too.
\fi
%shipsnet_train_analysis_deterministic.ipynb

\begin{table}[h]
\centering
\begin{tabular}{c|rrrrrrr|r}
\toprule
Model Variant  & RWG & WG & relu & relu6 & sigmoid & sin & tanh & \textbf{average} \\
\midrule
Model 0 &   \B{1}{.9838} &  \B{0}{.9750}  & \B{0}{.9762} &  \B{0}{.9712} &    \B{1}{.9700}   & \B{1}{.9825} & \B{0}{.9712} &          \B{0}{.9757} \\
Model 1 &   \B{1}{.9838} &  \B{0}{.9712} & \B{1}{.9788} &  \B{0}{.9788} &    \B{1}{.9700}   & \B{0}{.9712} & \B{0}{.9725} &          \B{0}{.9752} \\
Model 2 &   \B{0}{.9825} &  \B{1}{.9788} & \B{1}{.9788} &  \B{1}{.9800}   &    \B{0}{.9638} & \B{0}{.9762} & \B{1}{.9738} &          \B{1}{.9763} \\
Model 3 &   \B{0}{.9812} &  \B{0}{.9738} & \B{0}{.9738} &  \B{0}{.9775} &    \B{0}{.9675} & \B{0}{.9800}   & \B{0}{.9700}   &          \B{0}{.9748} \\
\midrule
Average & \B{1}{.9823} & \B{0}{.9747} & \B{0}{.9769} & \B{0}{.9769} & \B{0}{.9678} & \B{0}{.9775} & \B{0}{.9719} & \B{0}{.9755} \\
\bottomrule
\end{tabular}
\caption{Test accuracy of different activation functions across model variants in evaluating Deterministic CNN Models}
\label{table:activation_performance_deterministic_test}
\end{table}

%TODO Analysis of Test Result

Based on our test set, it is seen that the Rectified Weighted Gaussian is the best activation function for our Bayesian CNN Model. Model 2 is observed to result in a better test accuracy compared to all other model variants.

% TODO add explanation if possible

\iffalse
\begin{tabular}{rrrrrrrrr}
\hline
   Model Variant &   RWG &   WG &   relu &   relu6 &   sigmoid &    sin &   tanh &   mean\_test\_acc \\
\hline

\hline
\end{tabular}
\fi

%TODO add the testing result of Deterministic CNN using the test set

\subsection{Training Result of Bayesian CNN} \label{subsection:bayesiantrainresult}

\iffalse
In this section, I will explain:
\begin{enumerate}
\item The training loss for each epoch for each activation function
\item Compare train test accuracy
\item Spot if anything is weird
\item Highlight how the Bayesian CNN on the uniform might get an infinity Loss ==> part of future work suggestion
\end{enumerate}
\fi

Summary of test accuracy for Bayesian CNN Models, to find the best activation function and model variant, is provided in Table \ref{table:baytestactprior} and Table \ref{table:baytestactmodels}. The full result of test accuracy for each model variant, activation, prior, and prior $b$ are shown in Table \ref{table:activation_performance_bayesian_test-1} \ref{table:activation_performance_bayesian_test-0} and \ref{table:activation_performance_bayesian_test+1}.

Based on the test accuracy result, there are some highlights that we can make about the result. First, for models with prior $b=0.1$ (Table \ref{table:activation_performance_bayesian_test-1}) , relu6 yields better test accuracy compared to other activation function. Across all models, Gaussian and Laplace prior have a test accuracy around 75\%, while Uniform prior results in test accuracy around 90\%. Best accuracy among all models achieved by model 0, with average test accuracy of $0.8039$. Second, with prior $b=1$ (Table \ref{table:activation_performance_bayesian_test-0}), relu dominantly yields better test accuracy compared to other activation function. Best accuracy among all models achieved by model 1, with average test accuracy of $0.8332$. Lastly, with prior $b=10$ (Table \ref{table:activation_performance_bayesian_test+1}), tanh is observed to have a signifantly better test accuracy across almost all prior and activation function. Best accuracy among all models achieved by model 3, with average test accuracy of $0.8117$.

\iffalse
Based on the test accuracy result, there are some highlights that we can make about the result:
\begin{enumerate}
    \item With $b=0.1$, relu6 yields better test accuracy compared to other activation function. 
    \item With $b=0.1$, accross all models, gaussian and laplace prior have a test accuracy around 75\%, while uniform prior results in an accuracy around 90\%.
    \item With $b=0.1$, best accuracy among all models achieved by model 0, with average test accuracy of $0.8039$
    \item With $b=1$, relu dominantly yields better test accuracy compared to other activation function. 
    \item With $b=1$, best accuracy among all models achieved by model 1, with average test accuracy of $0.8332$
    \item With $b=10$, tanh dominantly yields better test accuracy compared to other activation function. 
    \item With $b=10$, best accuracy among all models achieved by model 3, with average test accuracy of $0.8117$
\end{enumerate}
\fi

While this research does not necessarily demand the model with the best accuracy to be the focus of our SEU injection process, we will proceed the SEU injection process to focus on all the priors with $\text{prior\_b}=1$, as shown in Table \ref{table:baytestactmodels} that also shows Model 1 is the best model among others. Its also observed from our experiment that across all model variations that we have, relu activation function shows the best average test accuracy, as shown in Table \ref{table:baytestactprior}. 

% Topline test accuracy (mean across models)
\begin{table}[H]
\centering
\begin{tabular}{l|rrrrrrr|r}
\toprule
Prior b & RWG & WG & relu & relu6 & sigmoid & sin & tanh & average \\
\midrule
0.1 & \B{0}{.8058} & \B{0}{.7859} & \B{1}{.8270} & \B{0}{.8249} & \B{0}{.7788} & \B{0}{.8034} & \B{0}{.7909} & \B{0}{.8024} \\
1 & \B{0}{.7618} & \B{0}{.7788} & \B{1}{.9123} & \B{0}{.8671} & \B{0}{.7892} & \B{0}{.8071} & \B{0}{.8832} & \B{1}{.8285} \\
10 & \B{0}{.7540} & \B{0}{.7674} & \B{0}{.8618} & \B{0}{.8090} & \B{0}{.8218} & \B{0}{.7489} & \B{1}{.8851} & \B{0}{.8068} \\
\midrule
Average & \B{0}{.7739} & \B{0}{.7774} & \B{1}{.8670} & \B{0}{.8336} & \B{0}{.7966} & \B{0}{.7865} & \B{0}{.8531} & \B{0}{.8126} \\
\bottomrule
\end{tabular}
\caption{Average test accuracy (mean across models) by prior and activation function. Best per prior b in bold.}
\label{table:baytestactprior}
\end{table}

\begin{table}[H]
\centering
\begin{tabular}{l|rrrrrrr|r}
\toprule
Model Variant & RWG & WG & relu & relu6 & sigmoid & sin & tanh & average \\
\midrule
Model 0 & \B{0}{.7733} & \B{0}{.7817} & \B{1}{.9046} & \B{0}{.8700} & \B{0}{.7979} & \B{0}{.7958} & \B{0}{.8783} & \B{0}{.8288} \\
%\midrule
Model 1 & \B{0}{.7800} & \B{0}{.7762} & \B{1}{.9125} & \B{0}{.8750} & \B{0}{.7979} & \B{0}{.8054} & \B{0}{.8854} & \B{1}{.8332} \\
%\midrule
Model 2 & \B{0}{.7471} & \B{0}{.7796} & \B{1}{.9117} & \B{0}{.8875} & \B{0}{.7746} & \B{0}{.8171} & \B{0}{.8754} & \B{0}{.8276} \\
%\midrule
Model 3 & \B{0}{.7467} & \B{0}{.7775} & \B{1}{.9204} & \B{0}{.8358} & \B{0}{.7862} & \B{0}{.8100} & \B{0}{.8938} & \B{0}{.8243} \\
\midrule
Average & \B{0}{.7618} & \B{0}{.7787} & \B{1}{.9123} & \B{0}{.8671} & \B{0}{.7892} & \B{0}{.8071} & \B{0}{.8832} & \B{0}{.8285} \\
\bottomrule
\end{tabular}
\caption{Average test accuracy by model for prior $b=1$. Best per model variant in bold.}
\label{table:baytestactmodels}
\end{table}

\subsubsection{A finding on training the Bayesian CNN model with uniform prior}

%TODO show the infinity loss, which epoch did they decide to go infinity. early? late?

%TODO explain the infinity loss

ELBO loss, which was explained in Equation (\ref{eq:genericelbo}), includes the joint probability $p_\theta(x, z)$, which can be break down in Equation (\ref{eq:uniformjp}). The ELBO loss becomes Equation (\ref{eq:genericelbo2}).

\noindent
\begin{minipage}{0.48\linewidth}
\begin{equation}
\log p_\theta(x, z) = \log p(x|z) + \log p(z)
\label{eq:uniformjp}
\end{equation}
\end{minipage}\hfill
\begin{minipage}{0.48\linewidth}
\begin{equation} \label{eq:genericelbo2}
\text{ELBO} \equiv \mathbb{E}_{q_\phi(z)} \big[ \log p(x|z) + \log p(z) - \log q_\phi(z) \big]
\end{equation}
\end{minipage}


As stated in Equation (\ref{eq:uniform}), the value of $p(z)$ depends on the bound of the uniform prior. For the value of $a=0$ and $b=1$, in this research, the value of $p(z)$ is stated as Equation (\ref{eq:uniformpzbound}). As the variational guide (in this case, the \texttt{AutoUniform}) samples a $z$ value such that $z \notin [0,1]$, the $\log p(z)$ becomes $-\infty$, making the ELBO loss from Equation (\ref{eq:genericelbo2}) becomes as Equation (\ref{eq:uniformboundelbo}).

% as shown in the bayesian training script, def _make_prior states that base = UniformReal(-self.prior_b, self.prior_b)

\noindent
\begin{minipage}{0.48\linewidth}
\begin{equation}
\log p(z) =
\begin{cases}
\log(0.5), & \text{if } z \in [0, 1] \\
\log(0) = -\infty, & \text{otherwise}
\end{cases}
\label{eq:uniformpzbound}
\end{equation}
\end{minipage}\hfill
\begin{minipage}{0.48\linewidth}
\begin{equation}
\text{ELBO} = \mathbb{E}_{q_\phi(z)} \big[ \log p(x|z) -\infty  - \log q_\phi(z) \big] = -\infty
\label{eq:uniformboundelbo}
\end{equation}
\end{minipage}

The epoch of which the training process start to experience $-\infty$ varies between activation functions and model variants, but an observable phenomenon can be seen across the variation of prior $b$ values. For $b=0.1$, as the range for the sampled weight can go beyond $[0,0.1]$ quick, the first epoch already encounters the $-\infty$ ELBO loss. For $b=1$, it's observed that most models can't get beyond 10 epoch, while for $b=10$ can go to epoch 60 to 70, or even finish the training in 100 epochs.

As the ELBO loss becomes $-\infty$, the process of the optimizer in updating the guide parameter breaks because of failing backpropagation process due to gradients that become $\infty$ or \texttt{NaN}. One way to prevent this from happening is by applying clipping mechanism to keep the sampled weights in the bound of the uniform prior. For this research we want to see the effect of SEU bitflips to the result of robustness score. Therefore, clipping is intentionally avoided to allow out-of-bound values to show their effects.
%nicely written!
%\subsection{Compare Training Result of Deterministic and Bayesian CNN}

\subsection{Test Accuracy comparison between Deterministic CNN Model and Bayesian CNN model}

Based on our findings in this chapter, \textbf{Bayesian CNN model is showing an inferior performance in test accuracy} compared to the deterministic CNN model. This observation is consistent across all activation functions, model variants, and Bayesian model activation function. The average accuracy of test data that can be achieved by the deterministic CNN model is 97.55\%, while the Bayesian CNN model can only achieve 82.85\% for prior $b$ value equal to 1. it's suspected that the performance difference is inherent to whether we are implementing Bayesian or deterministic CNN model. Here are some factors that can possibly result in this observation:
\begin{enumerate}[noitemsep, topsep=0pt]
\item All weights and bias in the Bayesian CNN neural network are initialized with a similar prior parameter ($a$ and $b$), assuming that all weight and bias' distributions are following the same distribution. 
\item Optimization of Bayesian Neural Networks with ELBO loss is harder compared to cross-entropy in deterministic CNN model. KL divergence in ELBO loss can pull the variational posterior distribution away from the optimal point to stay closer to the prior (posterior collapse \cite{posteriorcollapsevae2019}).
\end{enumerate}

\section{Analysis of Performance Against SEU} \label{section:SEUperformance}

In this section, most table will provide the aggregate value of SEU robustness measure by grouping the experiment result based on several criteria (e.g. model variation, activation function) and then averaging them. For simplicity, we denote these aggregated metrics as Average Absolute Accuracy Difference (AAAD) and Average Softmax Difference (ASD). Aggregated Robustness Index (ARIn) will be calculated based on these two values.

\subsection{Experiment Output} \label{subsection:experimentoutput}

%This includes plots of log difference of weight/bias before and after bitflips, compared to mean softmax difference and test delta.

%This also will explore how different activation function is more robust compared to others.

%TODO : Weird thing that is happening : Softmax Difference and Accuracy Difference is Significantly Higher in Bayesian, which is  => W E I R D

%\subsection{Compare SEU Result of Deterministic and Bayesian CNN}

%TODO Important : Create a graphic on how error propagates through the layers. Create the graph to be similar with the convolutional layer section.

For each SEU injection, the following metrics are recorded:
\begin{enumerate}[noitemsep, topsep=0pt]
    \item Accuracy difference : This metric is calculated as mentioned in  Equation (\ref{eq:accuracydelta}). Then, we compute the Absolute Accuracy Difference as shown in Equation (\ref{eq:absaccdiff}).
    \item Softmax difference, as mentioned in Equation (\ref{eq:softmaxdifference}).
    \item Absolute difference : This is calculated by substracting the newly observed value after SEU with the original value before SEU.
\end{enumerate}

%In this section, most of our robustness analysis will come from the absolute accuracy difference and softmax difference. We will also present our newly defined Aggregate Robustness Measure explained in Equation (\ref{eq:ARIn}) alongside the two, if presented in a tabular form. In most of the robustness presentation in a figure, we will plot the result in a scatter plot, which shows the softmax difference on the y axis, and the absolute accuracy difference on the x axis. In this visualization, we want to plant the intuition that the models that end up in the bottom left of the plot is more robust, and models that end up in the top right of the plot is the least robust. This plot is inspired by work in \cite{carbone2020a}, with some tweaks, especially on the usage of softmax difference instead of (1 - softmax) on the y axis.

In this section, the robustness analysis primarily relies on two metrics: the absolute accuracy difference and the softmax difference. In addition, we introduce our newly defined Aggregate Robustness Index (ARIn), as formalized in Equation (\ref{eq:ARIn}), which will be reported alongside these metrics when results are presented in tabular form. For most visualizations, robustness will be illustrated using scatter plots, where the x-axis represents the absolute accuracy difference and the y-axis represents the softmax difference. This representation is intended to convey the intuition that models located in the bottom left region of the plot exhibit greater robustness, whereas those in the top right region demonstrate the least robustness. The visualization approach is inspired by the work of \cite{carbone2020a}, with modifications, most notably the use of softmax difference rather than $1-\text{softmax}$ on the y axis.

The main visual cue used in these scatter plots will be mostly rely on colors, to show the different categories of our experiment result and how they manifested in the SEU robustness. An example of the aforementioned figure, that shows the unaggregated result of SEU injection Robustness, colored by the type of model used (Bayesian or deterministic), is shown in Figure \ref{fig:robboth} below. 


\begin{figure}[H]
\caption{Robustness of Deterministic and Bayesian CNN model}
\includegraphics[scale=0.5]{robboth2.png}
\centering
\label{fig:robboth}
\end{figure}

%Summary of absolute accuracy difference and softmax difference Deterministic vs Bayesian CNN
\begin{table}[H]
\centering
\begin{tabular}{lrrr}
\toprule
Model Type & AAAD & ASD & ARIn \\
\midrule
BNN & \B{0}{.0244} & \B{1}{.1074} & \B{1}{.0779} \\
DNN & \B{1}{.0149} & \B{0}{.2696} & \B{0}{.1909} \\
\bottomrule
\end{tabular}
\caption{Summary of average absolute accuracy difference and average softmax difference Deterministic vs Bayesian CNN}
\label{table:robbothcompare}
\end{table}


\subsubsection{Robustness of Deterministic CNN models}

From the aforementioned figure, we can infer that the impact of SEU on Absolute Accuracy Difference and Softmax Difference exhibits an approximately linear relationship, showing a clear correlation between the two performance metrics. This observation is similar from what has been observed from \cite{carbone2020a}, in which the author spotted a trade-off between pure accuracy and robustness (1 - softmax difference) for high performing deterministic networks against adversarial attack (FGSM).

\subsubsection{Robustness of Bayesian CNN models}

Contrary to the Deterministic CNN model SEU injection result that has a clear correlation between the two performance metric in a linear straight line, the performance shown by the Bayesian CNN model does not exhibit a same pattern. In Bayesian CNN Model SEU Injection from the figure above, some noticeable points make up a linear correlation of absolute accuracy difference and softmax difference, but there are also observations that does not support this relationship. Some SEU Injection results show high softmax difference, while still keeping the absolute accuracy difference low. 

\subsubsection{Robustness of Deterministic CNN models vs Bayesian CNN Models}

%TODO Explain why DNN better in absolute accuracy differenve, BNN better in softmax difference

Based on Table \ref{table:robbothcompare} it is observed that the deterministic CNN model is showing more robustness against SEU injection in terms of absolute accuracy difference. It shows that the average absolute accuracy difference if a deterministic CNN model is attacked by SEU injection is 1.49\%, while the Bayesian CNN model is expected to have an absolute accuracy difference of 2.44\%. 

This phenomenon can be triggered by the way that Bayesian Neural Network is sampling the weight/bias distribution parameter from the guide. In deterministic CNN model, the SEU only change one fixed weight/bias. In Bayesian CNN Model, as the distribution parameter is perturbed by the SEU, all Monte Carlo (MC) samples are drawn from the same perturbed distribution parameter. This makes every MC sample to be perturbed in the same way. This, will in turn disturb the average logits that is needed to classify the data with softmax. The resulting condition would be the data that is more prone to misclassification due to disturbed sampled weight and bias.

On the other hand, in terms of average softmax difference, Bayesian CNN model can be seen to demonstrate a better robustness against SEU injection. Softmax difference, which indicates the change in the predicted class probability, is observed to be 0.2696 on the Deterministic CNN Model, and 0.1074 on the Bayesian CNN Model. This means that Bayesian CNN Model has a smaller shift in output probabilities compared to the Deterministic CNN Model. 

The observed condition can happen due to how logits are calculated after MC sampling. As we run over the network 10 times, or as much as the number of samples is determined, logit shifts from the first run might be diluted by the subsequent run, as we are taking the average of all 10 runs to be the input of the softmax, that will classify our data. Although the chance of misclassifying some data sample is increased, Bayesian CNN model offer a more stable prediction probability against SEU.

To conclude our comparison, we employ the Aggregate Robustness Measure (ARIn) defined in Equation (\ref{eq:ARIn}). Based on this metric, we find that \textbf{Bayesian Neural Network exhibits greater robustness than the Deterministic Neural Network}, as evidenced by its smaller ARIn value (0.0779 compared to 0.1909 for the latter). In the applications where class confidence or predictive probability constitutes a critical factor, the Bayesian CNN model is therefore more suitable to use compared to the deterministic CNN model. Some applications can include medical imaging (e.g. If the model uncertainty exceeds the threshold, the case is "rejected" or flagged for human evaluation \cite{10.3389/fninf.2024.1346723} ) and ensemble modeling (in which multiple models are combined to produce aggregate predictions \cite{bishopbarber98}). 

\subsubsection{Possible explanation on the "high softmax difference but low accuracy difference" on the Bayesian CNN}

% TODO Explain why it could happen. - answered in this subsection

%Bayesian CNN model shown a better softmax difference robustness against SEU injection. This statement can be supported by the scatter plot in the Figure \ref{fig:robboth}, that shown so many Bayesian CNN models SEU injection result that is shown on the bottom left on the plot, indicating lower absolute accuracy difference and softmax difference, under the seemingly straight line made by the observation of the Deterministic model. Beside that, there are also a significant amount of result from the Bayesian CNN model that falls above the Deterministic CNN model, showing a high difference in softmax difference, but low accuracy difference (top left of the plot). 

The Bayesian CNN models demonstrate superior robustness in terms of softmax difference under SEU injection. This observation is supported by the scatter plot in Figure \ref{fig:robboth}, where a substantial number of SEU-injected Bayesian CNN results are located in the bottom-left region on the plot. This positioning indicates lower absolute accuracy differences and lower softmax differences, in contrast to the nearly linear distribution exhibited by the deterministic model. In addition, a considerable number of Bayesian CNN results appear in the top-left region of the plot, characterized by \textbf{high softmax differences but relatively low accuracy differences}, further distinguishing their behavior from that of the deterministic CNN.

%This observation results from the fact that in the Monte Carlo sampling process during inference (explained in Algorithm \ref{alg:msamplingbnninference} with number of samples = $10$), the variability of the logits among the samples, are increased by the SEU injection, but the mean logits before and after SEU injection remain similar. This makes the softmax function, which is nonlinear in the logits, provided a similar prediction on the classes, thus, no observable change in absolute accuracy difference is observed. On the other hand, our Deterministic CNN model only produce a single fixed set of logits, resulting logits and predicted class will likely be the same, unless the SEU moves the logits pass the decision boundary.

This observation arises from the Monte Carlo sampling process during inference (Algorithm \ref{alg:msamplingbnninference}, with number of samples = $10$). Under SEU injection, the variability of the logits across MC samples increases, while the mean logits before and after injection largely unchanged. Consequently, the nonlinear softmax function produces distributions with higher variability, which may lead to noticeable changes in softmax differences without altering the predicted class. This explains why many Bayesian CNN results fall in the bottom-left region of the scatter plot (low absolute accuracy difference and low softmax difference). In contrast, the deterministic CNN generates only a single fixed logit vector, so SEU injection affects the prediction only when the perturbation is sufficient to push the logits across a decision boundary, producing the near-linear distribution observed in the plot.

%To conclude, this is the uncertainty of Bayesian CNN model under SEU that came into play. The averaging process in inference can probably smooth out the bitflip, but there's also a chance that the averaging process amplified it. Our result has shown us that the prior statement is more likely.

To conclude, these results reflect the role of uncertainty in the Bayesian CNN under SEU injection. During inference, the averaging process across Monte Carlo samples can either smooth out the effect of a bitflip (low absolute accuracy difference and low softmax difference), or amplify it, resulting in a higher softmax difference while accuracy remains unchanged, as observed in the top-left of the plot. Our findings indicate that the smoothing effect dominates, which explains why the majority of Bayesian CNN SEU injection outcomes cluster in the bottom-left region.

%TODO add the table that shows Deterministic CNN vs Bayesian CNN

\subsection{Robustness by model variation}

%Summary of absolute accuracy difference and softmax difference by Model Variant
\iffalse
\begin{table}[H]
\centering
\begin{tabular}{lrrr}
\toprule
Model Name & AAAD & ASD & ARIn \\
\midrule
Model 0 & \B{0}{0.0145} & \B{0}{0.2694} & \B{0}{0.1908} \\
Model 1 & \B{1}{0.0142} & \B{0}{0.2699} & \B{0}{0.1911} \\
Model 2 & \B{0}{0.0157} & \B{0}{0.2702} & \B{0}{0.1914} \\
Model 3 & \B{0}{0.0152} & \B{1}{0.2687} & \B{1}{0.1903} \\
\bottomrule
\end{tabular}
\caption{Robustness of Deterministic CNN model. Metric is the average of values of each model variants.}
\label{table:robdetmodelgran}
\end{table}
\fi

%Summary of absolute accuracy difference and softmax difference by model variant for Deterministic and Bayesian CNN
\begin{table}[H]
\centering
\begin{tabular}{l|rr|rr|rr}
\toprule
Model Variant & \multicolumn{2}{c|}{AAAD} & \multicolumn{2}{c|}{ASD} & \multicolumn{2}{c}{ARIn} \\
 & DNN & BNN & DNN & BNN & DNN & BNN \\
\midrule
Model 0 & \B{0}{.0145} & \B{0}{.0243} & \B{0}{.2694} & \B{0}{.1106} & \B{0}{.1908} & \B{0}{.0800} \\
Model 1 & \B{1}{.0142} & \B{0}{.0254} & \B{0}{.2699} & \B{0}{.1109} & \B{0}{.1911} & \B{0}{.0805} \\
Model 2 & \B{0}{.0157} & \B{1}{.0221} & \B{0}{.2702} & \B{1}{.0944} & \B{0}{.1914} & \B{1}{.0685} \\
Model 3 & \B{0}{.0152} & \B{0}{.0257} & \B{1}{.2687} & \B{0}{.1138} & \B{1}{.1903} & \B{0}{.0825} \\
\bottomrule
\end{tabular}
\caption{Robustness of Deterministic and Bayesian CNN model by model variant}
\label{table:robbothmodelgran}
\end{table}

\subsubsection{Robustness of Deterministic CNN models}

%Figure \ref{fig:robdetmodelgran} shows the performance metric of each SEU injection for Deterministic CNN Model, colored by model variant, without any aggregation. Table \ref{table:robbothmodelgran} shows an aggregated version from Figure \ref{fig:robdetmodelgran}, which group all the SEU Injection by model variants and averages the Absolute Accuracy Difference and Softmax Difference. We can conclude from this table that Model 3, which is the model with weight decay during training, has the best robustness compared to other Deterministic CNN model variant in terms of our newly defined Aggregate Robustness Index (ARIn). While having a worse absolute accuracy difference compared to the base model (Model 0), Model 3 is better in the softmax difference. Lower softmax difference is expected in model trained with weight decay during training, as it drives the weight into a smaller magnitude, which in turn produce smaller logits and reduce the softmax difference. In this occasion, we would also like to compare the result that we have with previous research. While both research does not have an identical experiment settings, work in \cite{dennispoperobustcnn} implicitly shown that Model 2 shown the best absolute accuracy difference, which is the opposite of what we found on our research that shows Model 2 has the worst absolute accuracy difference. In our research, model with smartpooling (Model 1) is the model with best absolute accuracy difference.

Figure \ref{fig:robdetmodelgran} shows the performance metric of each SEU injection for Deterministic CNN Model, colored by model variant, without any aggregation. Table \ref{table:robbothmodelgran} summarizes these results by averaging the Absolute Accuracy Difference (AAAD) and Softmax Difference (ASD) across model variants. Based on the Aggregate Robustness Index (ARIn), \textbf{Model 3 (trained with weight decay) demonstrates the strongest robustness}. Although it exhibits slightly worse AAAD compared to the base model (Model 0), it achieves a noticeably lower ASD. This is consistent with the expectation that weight decay reduces the magnitude of the weights, leading to a smaller logits and hence smaller softmax differences.

In comparison to prior work \cite{dennispoperobustcnn}, our findings diverge. While it suggests that a variant equivalent to our Model 2 achieved the best AAAD, our experiments indicate that Model 2 has the least robust AAAD, while Model 1 (with smartpooling) performs best in AAAD. 

\subsubsection{Robustness of Bayesian CNN models}

%The resulting performance metric from each SEU injection for our Bayesian CNN model is shown in Figure \ref{fig:robbaymodelgran}. Each point in the scatter plot represent the absolute accuracy difference and softmax difference for each injection without aggregation, colored by model variant. 

%Summary of absolute accuracy difference and softmax difference by model variant for Bayesian CNN
\iffalse
\begin{table}[H]
\centering
\begin{tabular}{llrrr}
\toprule
model variant & AAAD & ASD & ARIn \\
\midrule
Model 0 & \B{0}{0.0243} & \B{0}{0.1106} & \B{0}{0.0800} \\
Model 1 & \B{0}{0.0254} & \B{0}{0.1109} & \B{0}{0.0805} \\
Model 2 & \B{1}{0.0221} & \B{1}{0.0944} & \B{1}{0.0685} \\
Model 3 & \B{0}{0.0257} & \B{0}{0.1138} & \B{0}{0.0825} \\
\bottomrule
\end{tabular}
\caption{Robustness of Bayesian CNN model. Metric is the average of values of each model variants}
\label{table:robbaymodelgran}
\end{table}
\fi

%Table \ref{table:robbothmodelgran} exhibits that the Bayesian CNN model variant is the most robust, based on our Aggregated Robustness Metric, is Model 2. Model 2, which was designed with dropout layers during training, shown the lowest absolute accuracy difference and lowest softmax difference. While it is not comparable, the result in which the Model 2 is the more robust among all model variants is similar to what has been found in \cite{dennispoperobustcnn}. While the absolute accuracy difference is similar across all model variants, it is noticeable that the softmax difference for Model 2 is comparably low compared to other model variants.

Figure \ref{fig:robbaymodelgran} shows the performance metrics of SEU injections for the Bayesian CNN models, with points colored by model variant. As shown in Table \ref{table:robbothmodelgran}, \textbf{Model 2 (trained with dropout) is the most robust among the Bayesian variants}, exhibiting the lowest values for both AAAD and ASD. Although AAAD is generally similar across Bayesian variants, the reduced ASD of Model 2 makes its particularly notable. This outcome aligns with findings in \cite{dennispoperobustcnn}, despite differences in experimental setup.

\subsubsection{On the robustness of Bayesian CNN model trained with dropout layer}

The significant robustness of Model 2 compared to other model variants in the Bayesian setting can be attributed by the nature of Dropout layers that demonstrate redundancy during training. As some neurons are inactive during training, our neural network can not just too reliant on any single neuron to carry critical information/infer a feature. In our Bayesian CNN model without dropout layer, there might be a chance that a SEU/bitflips can hit a dominant neuron and shift our prediction significantly. With a regularized model with dropout during training, the resulting model has no dominant neuron, so during inference, other neurons are as contributing as others, resulting in less shift in prediction if one of the neurons is being hit by bitflip.

%bayesian vs deterministic per metric


\subsection{Robustness by prior and prior parameter}

\iffalse
%Summary of absolute accuracy difference and softmax difference by prior and prior parameter for Deterministic and Bayesian CNN
\begin{table}[H]
\centering
\begin{tabular}{ll|rr|rr|rr}
\toprule
prior & param\_type\_group & \multicolumn{2}{c|}{AAAD} & \multicolumn{2}{c|}{ASD} & \multicolumn{2}{c}{ARIn} \\
 &  & DNN & BNN & DNN & BNN & DNN & BNN \\
\midrule
No prior & No param type group & \B{1}{.0149} & \B{0}{nan} & \B{1}{.2696} & \B{0}{nan} & \B{1}{.1909} & \B{0}{nan} \\
gaussian & center & \B{0}{nan} & \B{1}{.0177} & \B{0}{nan} & \B{1}{.0794} & \B{0}{nan} & \B{1}{.0575} \\
gaussian & spread & \B{0}{nan} & \B{0}{.0293} & \B{0}{nan} & \B{0}{.1281} & \B{0}{nan} & \B{0}{.0929} \\
laplace & center & \B{0}{nan} & \B{0}{.0193} & \B{0}{nan} & \B{0}{.0823} & \B{0}{nan} & \B{0}{.0598} \\
laplace & spread & \B{0}{nan} & \B{0}{.0287} & \B{0}{nan} & \B{0}{.1306} & \B{0}{nan} & \B{0}{.0946} \\
uniform & center & \B{0}{nan} & \B{0}{.0231} & \B{0}{nan} & \B{0}{.1054} & \B{0}{nan} & \B{0}{.0763} \\
uniform & spread & \B{0}{nan} & \B{0}{.0303} & \B{0}{nan} & \B{0}{.1277} & \B{0}{nan} & \B{0}{.0928} \\
\bottomrule
\end{tabular}
\caption{Robustness of Deterministic and Bayesian CNN model by prior and prior parameter}
\label{table:robbothpriorpriorparam}
\end{table}
\fi

%Summary of absolute accuracy difference and softmax difference by prior and prior parameter for Bayesian CNN
\begin{table}[H]
\centering
\begin{tabular}{l|rr|rr|rr}
\toprule
Prior & \multicolumn{2}{c|}{AAAD} & \multicolumn{2}{c|}{ASD} & \multicolumn{2}{c}{ARIn} \\
Prior Parameter & Center & Spread & Center & Spread & Center & Spread \\
\midrule
Gaussian & \B{1}{.0177} & \B{0}{.0293} & \B{1}{.0794} & \B{0}{.1281} & \B{1}{.0575} & \B{0}{.0929} \\
Laplace & \B{0}{.0193} & \B{1}{.0287} & \B{0}{.0823} & \B{0}{.1306} & \B{0}{.0598} & \B{0}{.0946} \\
Uniform & \B{0}{.0231} & \B{0}{.0303} & \B{0}{.1054} & \B{1}{.1277} & \B{0}{.0763} & \B{1}{.0928} \\
\bottomrule
\end{tabular}
\caption{Robustness of Bayesian CNN model by prior and prior parameter}
\label{table:robbaypriorpriorparam2}
\end{table}

In this section, we will only discuss the findings from SEU attack on Bayesian CNN models, as the deterministic CNN models does not have prior and prior parameter in their architecture. Table \ref{table:robbaypriorpriorparam2} aggregates the absolute accuracy difference and softmax difference of our experiment, grouping them by the prior distribution (Gaussian, Laplace, and Uniform) and the prior parameters, which are center ($a$) and spread ($b$). 

\subsubsection{Which prior is the most robust?}

For the absolute accuracy difference, both Gaussian and Laplace prior have a similar result on the center and spread, with the Uniform distribution is the worst among all prior distributions. In terms of softmax difference, the Gaussian and Laplace distribution also performs similarly on the prior center and are both better compared to the Uniform distribution. A different observation can be found on the spread, which shows the Uniform and Gaussian perform similarly and better compared to the Laplace distribution. Overall, we can conclude that the \textbf{Gaussian prior is the most robust prior} among all, with Laplace prior shown a better accuracy robustness on prior spread attack, and the uniform distribution shown a better softmax robustness on prior spread attack.

\subsubsection{Gaussian prior vs Laplace prior}

%TODO explain why in general Gaussian is better

As shown in Figure \ref{fig:priorpdfs}, each prior distribution have a different Probability Distribution Function (PDFs). The Gaussian distribution has lighter tails compared to the Laplace distribution with the same parameter. As we shift the mean for both distribution, we also pull the tail samples toward or accross the decision boundary. Because Laplace distribution have a heavier tail, the probability of tail samples to get sampled is higher, thus increasing the likelihood of prediction changes. The Uniform distribution is unsurprisingly doing worst, as the probability density of Uniform distribution is flat, making the samples evenly spread across the specified range. A shift in $a$ value for uniform distribution move a larger fraction of samples across the decision boundary, with the same probability of being sampled, making the prediction accuracy worse.

%TODO explain why the other two are good on their each acc diff and softmax

Laplace distribution is less affected in absolute accuracy difference if the spread parameter is being attacked compared to the Gaussian distribution, due to the fact that Gaussian distribution is more concentrated in the center. If we increase or do bitflip on the scale parameter for Gaussian distribution, the distribution of sample near the center widen and possibly move the sampled weight/bias across the decision boundary. For the Laplace distribution, modifying the scale parameter does not change the sampled value much, as the probability density of Laplace is already wide relative to the Gaussian prior. 

\subsubsection{Uniform prior robustness in softmax difference against the spread SEU injection}

Since changing the scale parameter for the Gaussian and Laplace prior affect the peak height and decay rate in the tails, if we do Monte Carlo sampling for weight/bias sampling, the difference of sample values is expected to be higher compared to the Uniform prior that have a flat PDF. This is due to the fact that the high-density center is now less dense and mass shifts into the tails, which will amplify the variability in logits. The uniform distribution can yield a better softmax difference, due to the fact that the each value of weight/bias have a similar probability to be sampled. A slight change in Uniform's scale parameter, for a latent weight $x$, will result in a smaller change of probability, compared to the change for the same latent weight $x$ in Gaussian/Laplace prior.

% should I mention about the changing of width

\subsubsection{On the robustness of prior center vs prior spread against SEU}

%TODO add explanation. Why the center parameter is more robust than the spread

Based on the Table \ref{table:robbaypriorpriorparam2} above, we can also retrieve a conclusion about how the SEU interacts with the center prior parameter and spread prior parameter. For all metrics, \textbf{the center prior parameter is consistently more robust than the spread prior parameter}. Robustness of Bayesian Neural Networks against SEU injection on the prior center parameter comes from the nature of shifting the mean of the prior that moves all the weight samples by the same absolute amount, preserving the relative positions and allowing them to cancel out during the Monte Carlo averaging. Shift in the spread parameter will change the weight sampled to be significantly different from the mean (either pushing some further into the tails or other closer to the center), making the resulting logits to be much different than before the SEU injection happened.

\subsection{Robustness by activation function}

%bayesian vs deterministic per metric
\begin{table}[H]
\centering
\begin{tabular}{l|rr|rr|rr}
\toprule
Activation Function & \multicolumn{2}{c|}{AAAD} & \multicolumn{2}{c|}{ASD} & \multicolumn{2}{c}{ARIn} \\
 & DNN & BNN & DNN & BNN & DNN & BNN \\
\midrule
RWG & \B{1}{.0078} & \B{1}{.0168} & \B{0}{.2651} & \B{0}{.0958} & \B{0}{.1875} & \B{0}{.0688} \\
WG & \B{0}{.0178} & \B{0}{.0213} & \B{0}{.2716} & \B{0}{.0964} & \B{0}{.1924} & \B{0}{.0698} \\
relu & \B{0}{.0234} & \B{0}{.0379} & \B{0}{.2835} & \B{0}{.1626} & \B{0}{.2012} & \B{0}{.1180} \\
relu6 & \B{0}{.0114} & \B{0}{.0248} & \B{0}{.2653} & \B{0}{.1012} & \B{0}{.1878} & \B{0}{.0737} \\
sigmoid & \B{0}{.0136} & \B{0}{.0190} & \B{1}{.2598} & \B{1}{.0800} & \B{1}{.1839} & \B{1}{.0581} \\
sin & \B{0}{.0141} & \B{0}{.0294} & \B{0}{.2702} & \B{0}{.1056} & \B{0}{.1913} & \B{0}{.0775} \\
tanh & \B{0}{.0163} & \B{0}{.0214} & \B{0}{.2714} & \B{0}{.1102} & \B{0}{.1923} & \B{0}{.0794} \\
\bottomrule
\end{tabular}
\caption{Robustness of Deterministic and Bayesian CNN model by activation function}
\label{table:robbothact}
\end{table}

Based on Table \ref{table:robbothact}, we can observe that the Rectified Weighted Gaussian activation function has the most robust absolute accuracy difference among all activation function for the deterministic CNN model (which has been proven in \cite{dennispoperobustcnn}) and Bayesian CNN Model. On our newly introduced softmax difference measure, the sigmoid activation function achieves the lowest Average Softmax Difference (ASD) in both settings.

%TODO why the RWG might be the best

Relu, which is the most expressive activation function that we have in this experiment (as seen from Figure \ref{fig:activationfunctions}), shown a comparably better test accuracy (Table  \ref{table:activation_performance_bayesian_test-1}) despite showing the worst robustness. RWG, which yields the lowest test accuracy, achieves the strongest robustness, especially in terms of absolute accuracy difference. This observation shown us a symptom of a \textbf{trade-off between test accuracy and SEU robustness}, which is also observed in the work of \cite{carbone2020a}.

It can also be seen from Table \ref{table:robbothact} that the sigmoid activation function have the second best robustness in absolute accuracy difference on BNNs, while also having the best softmax difference. In terms of SEU, which is a change of weight and bias before activation function is implemented, let's get back on how the neuron is calculated in the feed-forward process. This process, as mentioned in \cite{bishoppatternrecognition2007} is notated as:

\begin{equation}
\label{eq:neuronfeedforward}
a_j = \sum_{i=1}^{D} w^{(1)}_{ji} x_i + w^{(1)}_{j0}
\end{equation}

%page 227
where $x_1,x_2,\dots,x_D$ are inputs, $j=1,2,\dots,M$ and the superscript (1) represents the parameters in the first 'layer' of the network, which have a 'width' of $M$. $w^{(1)}_{j0}$ is  Each of $a_j$ is then transformed in the following equation:

\begin{equation}
z_j = h(a_j)
\end{equation}

where $h(\cdot)$ is a differentiable activation function.

$w^{(1)}_{ji}$ which represents weight, or $w^{(1)}_{j0}$ which represents bias, are the component in this feed-forward process that is prone to the SEU injection. As one of both values changed into SEU injected value $a'_j$, we would like to know how this change can effect the value of $z$ into $z'_j$, the lower the change, the difference in softmax value will be also lower.

In this experiment, the weight/bias is expected to be near zero with our setting of prior parameter $a=0$ and $b=1$. The majority of bitflips that we do will result in a change of weight/bias that is near zero, when we hit bit index that is not 0 (sign bit) and 1 (first exponent bit). Therefore, let's see the impact of changing $a$ to $z$ by observing the differential of our activation function $h'()$. We will take a look into sigmoid and the seemingly similarly shaped tanh, and its differential near 0.

$$
\sigma(z) = \frac{1}{1 + e^{-z}} ,\qquad\frac{d\sigma(z)}{dz} = \sigma(z)\big(1 - \sigma(z)\big) ,\qquad
\sigma(0) = 0.5 \quad \Rightarrow \quad \sigma'(0) = 0.25
$$

$$
\tanh(z) = \frac{e^z - e^{-z}}{e^z + e^{-z}} ,\qquad \frac{d}{dz}\tanh(z) = 1 - \tanh^2(z) ,\qquad
\tanh(0) = 0 \quad \Rightarrow \quad \tanh'(0) = 1
$$

From our calculation of diferential for both activation functions above, we can see that the differential on sigmoid at 0 is 0.25, while the tanh is at 1. A slight change in input for tanh have a bigger shift in the value after activation function. Therefore, we can conclude that the sigmoid function gained robustness in SEU because of a smaller gradient near $z=0$. 

It will also be interesting to look at the gradient of our activation function on extreme values, especially in looking in the effect of SEU injection on bit index 1 (first exponent). \cite{bishoppatternrecognition2007} mentioned that the sigmoid decay asymptotically like $e^{-z}$ for $z \to +\infty$, while tanh, which has $tanh'(z) = 1- tanh^2(z)$, will decay in $4e^{-2z}$ for $z \to +\infty$, which is faster than the sigmoid. We can assume that the tanh function will dampen the $z$ value better on the extreme values.


%While the shape of sigmoid activation function is similar to tanh, the sigmoid's output is bounded in $[0,1]$ while tanh is bounded in [-1,1].

%TODO : Include if there are any research papers that support on why RELU is not robust. or why Relu is popular in the first place.

%TODO why Relu is the worst, then pinpoint that relu 6 can keep it

%TODO Explain why sigmoid might be robust : might be related on how it's bounded on [0,1] and has easier slope compared to tanh

\subsection{Robustness by layer and layer parameter}

\iffalse
\begin{figure}[H]
\caption{Robustness of Bayesian and Deterministic CNN model by layer and layer parameter}
\includegraphics[scale=0.6]{robbothlayermodule.png}
\centering
\label{fig:robbothlayermodule}
\end{figure}
\fi

%Summary of absolute accuracy difference and softmax difference by bit layer and module for Deterministic and Bayesian CNN
\begin{table}[H]
\centering
\begin{tabular}{ll|rr|rr|rr}
\toprule
Layer & Layer Parameter & \multicolumn{2}{c|}{AAAD} & \multicolumn{2}{c|}{ASD} & \multicolumn{2}{c}{ARIn} \\
 &  & DNN & BNN & DNN & BNN & DNN & BNN \\
\midrule
conv1 & bias & \B{0}{.0048} & \B{0}{.0095} & \B{0}{.2651} & \B{1}{.0641} & \B{0}{.1875} & \B{1}{.0458} \\
conv1 & weight & \B{0}{.0070} & \B{0}{.0129} & \B{0}{.2696} & \B{0}{.0762} & \B{0}{.1907} & \B{0}{.0546} \\
\midrule
conv2 & bias & \B{1}{.0021} & \B{1}{.0089} & \B{1}{.2622} & \B{0}{.0723} & \B{1}{.1854} & \B{0}{.0515} \\
conv2 & weight & \B{0}{.0058} & \B{0}{.0136} & \B{0}{.2640} & \B{0}{.0987} & \B{0}{.1868} & \B{0}{.0704} \\
\midrule
fc1 & bias & \B{0}{.0326} & \B{0}{.0459} & \B{0}{.2773} & \B{0}{.1720} & \B{0}{.1974} & \B{0}{.1259} \\
fc1 & weight & \B{0}{.0372} & \B{0}{.0555} & \B{0}{.2792} & \B{0}{.1611} & \B{0}{.1991} & \B{0}{.1205} \\
\bottomrule
\end{tabular}
\caption{Robustness of Deterministic and Bayesian CNN model by layer and layer parameter}
\label{table:robbothlayermodule}
\end{table}

%Table \ref{table:robbothlayermodule} summarizes the SEU Injection result by averaging the value of Absolute Accuracy Difference and Softmax Difference across layer and layer parameter of the CNN architecture. Based on that table, followed by its illustration in Figure \ref{fig:robbothlayermodule}, it is observed that convolutional layer 2 is the most robust, followed by convolutional layer 1. The fully connected layer is the least robust, as it exhibists the highest mean of softmax difference and highest mean absolute accuracy difference. This findings also align with the result shown in \cite{dennispoperobustcnn}, where the further convolutional layer (conv2 and conv3) have the lowest average test delta and fully connected layer has the highest.

Table \ref{table:robbothlayermodule} summarizes the SEU Injection result by averaging the value of Absolute Accuracy Difference and Softmax Difference across layers and layer parameter of the CNN architecture. Based on that table, it is observed that convolutional layer 2 is the most robust, followed by convolutional layer 1. The fully connected layer is the least robust, as it exhibits the highest mean of softmax difference and the highest mean absolute accuracy difference. These findings also align with the result shown in \cite{dennispoperobustcnn}, where the further convolutional layer (conv2 and conv3) have the lowest average test delta and the fully connected layer has the highest.

\subsubsection {Layer SEU Robustness, propagation depth}

%TODO explain how error propagates through the layers. Error in early layers are likely to be 'dampened' by the activation function etc. Error on the far end of the layer will not be dampened well and captured by the prediction layer.

Fully connected layer, which is the furthest layer we have in our architecture, is the least robust due to how the SEU injection impacts the final logits. If we do the SEU injection early in the convolutional 1 layers, there are multiple occassion that nonlinear process will dampen the effect of SEU. 

First, every convolution layer will be followed by an activation function. Nonlinear activation function can help reduce the impact of SEU or bitflips. After the activation function, the result is taken into pooling layer, in which the multiple values can are aggregate into a smaller dimension. If the pooling layer is utilizing the average pooling layer, the value of the SEU will be averaged with other values, reducing the impact. In our case study, we are using max pooling. The value of the SEU might just be propagated to the next layer if the resulting value is the highest among all. Therefore, as the work in \cite{dennispoperobustcnn} suggested, smart pooling is used, to remove values that are above specified threshold, especially if the value was hit by SEU.

Let's compare this condition with fully connected layer from our model architecture. In this layer, the SEU injected values will not undergo any non-linear transformation like what we had in the convolutional layers. Therefore, this value will directly impact the calculation of logits and class probability, thus explaining the higher absolute accuracy difference and softmax difference.

%TODO why the most robust in conv2, not conv 1?

%The illustration on how SEU injection effects can propagate through the network and its impact is shown in Figure \ref{fig:errorpropillust}. An SEU injected in early convolutional layers, as shown in Figure \ref{fig:errorpropillust} (a) is symbolized as red lightning to show its devastating impact. The subsequent impact is marked yellow to show that SEU effect is weakened, followed by small green lightning on the FC1 layer, showing that the impact of SEU is further diminished. The sinusoidal wave in the neuron is intended to show that the attack undergo non-linear transformation, dampening the impact of the SEU injected value. The colored cloud at the end of the network represents the impact of SEU on the logits, with red means that the logits is impacted the worst (as shown in Figure \ref{fig:errorpropillust} (b)), and green is the least.

Figure \ref{fig:errorpropillust} illustrates how the effects of an SEU injection propagate through a CNN model. In panel (a), an SEU injected into an early convolutional layer (indicated by the red lightning bolt) undergoes attenuation as it passes through subsequent layers. This weakening is shown by yellow and then smaller green lightning symbols in later layers, reflecting the diminishing influence of the fault. The sinusoidal symbols within the neurons represent non-linear transformations, which further dampen the perturbation. Consequently, the effect on the final logits is minimal, represented by the green cloud at the output.

In contrast, panel (b) shows an SEU injected directly into the fully connected (FC1) layer. Here, the perturbation is not mitigated by earlier non-linear transformations and instead propagates strongly to the output layer, producing the most severe distortion in the logits, indicated by the red clouds.

\subsubsection{Layer component SEU Robustness, weight vs bias}

Based on the result of our experiment, we can see that the weight is less robust compared to the bias, when we inject them with SEU/biftlips. This is due to the nature of the neural networks itself. Looking back at Equation (\ref{eq:neuronfeedforward}), we can see that the weight is multiplied by the input value $x_i$ to compute $a_j$. \textbf{Change in weight to the value of $a_j$ is amplified by the $x_i$}, compared to the change in bias that is merely a constant.

%TODO add illustration on how error propagates.
%TODO add illustration on error in weight is worse than in bias. show the z=wx+b,



\begin{figure}[H]
\caption{Illustration of SEU injection and its propagated effect on conv1 layer (a) and fc1 layer (b) in a CNN model}
\includegraphics[scale=0.6]{error-propagation-illustration.png}
\centering
\label{fig:errorpropillust}
\end{figure}

\subsection{Robustness by attacked bit location} \label{robattackedbit}

%Summary of absolute accuracy difference and softmax difference by bit location for Deterministic and Bayesian CNN
\begin{table}[H]
\centering
\begin{tabular}{l|rr|rr|rr|rr}
\toprule
Bit Index & \multicolumn{2}{c|}{Average Log Absolute Difference} & \multicolumn{2}{c|}{AAAD} & \multicolumn{2}{c|}{ASD} & \multicolumn{2}{c}{ARIn} \\
 & DNN & BNN & DNN & BNN & DNN & BNN & DNN & BNN \\
\midrule
0 & \B{0}{-1.1804} & \B{0}{-0.7909} & \B{0}{.0005} & \B{0}{.0075} & \B{0}{.2613} & \B{0}{.0533} & \B{0}{.1848} & \B{0}{.0381} \\
\midrule
1 & \B{0}{37.0504} & \B{0}{37.8829} & \B{0}{.1023} & \B{0}{.1206} & \B{0}{.3193} & \B{0}{.4080} & \B{0}{.2371} & \B{0}{.3009} \\
3 & \B{0}{-1.4815} & \B{0}{-0.6454} & \B{0}{.0002} & \B{0}{.0070} & \B{0}{.2612} & \B{0}{.0527} & \B{1}{.1847} & \B{0}{.0376} \\
6 & \B{0}{-0.5849} & \B{0}{-0.4291} & \B{0}{.0013} & \B{0}{.0070} & \B{0}{.2617} & \B{0}{.0529} & \B{0}{.1850} & \B{0}{.0378} \\
\midrule
10 & \B{0}{-2.2308} & \B{0}{-1.3945} & \B{0}{.0001} & \B{1}{.0067} & \B{1}{.2611} & \B{1}{.0525} & \B{1}{.1847} & \B{1}{.0374} \\
15 & \B{0}{-3.7360} & \B{0}{-2.8997} & \B{1}{.0000} & \B{1}{.0067} & \B{1}{.2611} & \B{1}{.0525} & \B{1}{.1847} & \B{1}{.0374} \\
21 & \B{1}{-5.5422} & \B{1}{-4.7058} & \B{1}{.0000} & \B{1}{.0067} & \B{1}{.2611} & \B{0}{.0528} & \B{1}{.1847} & \B{0}{.0376} \\
\bottomrule
\end{tabular}
\caption{Robustness of Deterministic and Bayesian CNN model by attacked bit location}
\label{table:robbothbitloc}
\end{table}

\iffalse
\begin{figure}[H]
\caption{Log mean absolute difference against robustness metrics, excluding the data from bit index=1}
\includegraphics[scale=0.4]{logabsdiff.png}
\centering
\label{fig:logabsdiff}
\end{figure}
\fi

%To comprehend more about how this could happen, and simplify our calculation in defining the difference of values before and after SEU, we calculate the log of absolute difference, which takes the input from Subsection \ref{subsection:experimentoutput}, and formulated as follows:

%To aid our understanding of the impact of SEU injection in different attacked bit location to the value of original weight/bias, we introduce  the log of absolute difference, which takes the input from Subsection \ref{subsection:experimentoutput}, and formulated as follows:

To analyze the effect of SEU injections at different bit positions, we compute the log absolute difference of weight or bias values before and after an SEU, defined as:

\begin{equation}
\text{Log Absolute Difference} = \log_{10}(|\text{value after SEU}-\text{value before SEU}|)
\end{equation}

%Attacked bit index in this experiment are sign (index 0), exponent (index 1,3,6), and mantissa (index 10, 15, 21). Based on Table \ref{table:robbothbitloc}, the average log absolute difference of attack against bit index 1 is around 37, which is is the worst among all SEU injection bit index. We can see that bit index 1 is the significantly least robust in terms of AAAD, ASD, which consequently shows the worst ARIn. 

In this experiment, the attacked bit index in this experiment are sign (index 0), exponent (index 1,3,6), and mantissa (index 10, 15, 21). As shown in Table \ref{table:robbothbitloc}, flipping bit index 1 yields an average log absolute difference of approximately 37 and results in the worst robustness across all metrics (AAAD, ASD, and ARIn).

%It is observed that the attack on mantissa have almost no effect on the observed absolute accuracy difference, with the value of absolute accuracy difference observed is almost 0. Bit index 3 also observed to have a similar impact on both absolute accuracy difference and softmax difference, compared with the mantissa bits. 

By contrast, perturbations to the mantissa bits have negligible impact on AAAD, with values close to zero, and produce robustness metrics nearly identical to those from bit index 3 in the exponent. Attacks on bit index 6 also yield slightly higher AAAD than the mantissa bits, through still minor compared to bit index 1.

%Based on Table \ref{table:robbothbitloc}, which is summarized in Figure \ref{fig:logabsdiff2}, we can see that a correlation between log absolute difference with average absolute accuracy difference and average softmax difference cannot be established. SEU in bit index 3 and 6 yield higher average accuracy difference compared to other bit index, excluding bit index 1. It is also observed that attack in bit index 0 results in a higher average softmax difference.

Notably, no clear correlation can be established between the log absolute difference and robustness metrics such as AAAD or ASD. For example, bit index 0 (sign bit) produces a moderate log absolute difference yet yields slightly higher ASD compared to mantissa flips. Overall, these results confirm that SEUs affecting the exponent, particularly at bit index 1, are the most harmful to model robustness.

% SHOULD I mention about the MC logits averaging again, here?

\subsection{Robustness by original bit attacked}

%Summary of absolute accuracy difference and softmax difference by original bit condition for Deterministic and Bayesian CNN
\begin{table}[H]
\centering
\begin{tabular}{l|rr|rr|rr}
\toprule
Original Bit Value & \multicolumn{2}{c|}{AAAD} & \multicolumn{2}{c|}{ASD} & \multicolumn{2}{c}{ARIn} \\
 & DNN & BNN & DNN & BNN & DNN & BNN \\
\midrule
0 & \B{0}{.0275} & \B{0}{.0455} & \B{0}{.2767} & \B{0}{.1734} & \B{0}{.1966} & \B{0}{.1268} \\
1 & \B{1}{.0002} & \B{1}{.0069} & \B{1}{.2612} & \B{1}{.0529} & \B{1}{.1847} & \B{1}{.0378} \\
\bottomrule
\end{tabular}
\caption{Robustness of Deterministic and Bayesian CNN model by original bit attacked}
\label{table:robbothorbit}
\end{table}

In this subsection, we also compare how the robustness is affected by the original bit value/state before the SEU.Table \ref{table:robbothorbit} shows that bitflip from 1 to 0 exhibits an extremely more robust model compared to bitflip from 0 to 1, especially on the absolute accuracy difference. It's not aligned with the claim from \cite{Hanif2021} that stated the opposite. 

%TODO explain why and how, especially on how the bitflip from 1 to 0 is extremely more robust.

%As the weight and bias we have is near 0 and 1, the profile of the bit index before the SEU injection also play a big role in our experiment result. For example, it is observed that for all SEU injection on bit index 1, the original bit value/state is always 0. Because our previous findings about bit index 1 SEU injection in Subsection \ref{robattackedbit} shows its extreme vulnerability towards SEU injection, the result from Table \ref{table:robbothorbit} are consistent with these findings.

This difference can be explained by the distribution of weights and biases in our models, which are typically close to 0 or 1. Consequently, the original bit value/state play an important role in determining the impact of SEUs. For instance, in all SEU injections involving bit index 1, the original bit value/state was always 0. Since Subsection \ref{robattackedbit} already established that bit index 1 is highly vulnerable to SEUs, the observed lack of robustness in bitflip from 0 to 1 are consistent with these findings.

%TODO can explain more the bit distribution in the original bit and show the reader that it's always 0, thus flipping to 1 will have a huge impact.

\iffalse
\subsection{Robustness by neuron position}

%Summary of absolute accuracy difference and softmax difference by neuron position for Deterministic and Bayesian CNN
\begin{table}[H]
\centering
\begin{tabular}{l|rr|rr|rr}
\toprule
Neuron Position & \multicolumn{2}{c|}{AAAD} & \multicolumn{2}{c|}{ASD} & \multicolumn{2}{c}{ARIn} \\
 & DNN & BNN & DNN & BNN & DNN & BNN \\
\midrule
First neuron & \B{1}{.0146} & \B{1}{.0211} & \B{1}{.2689} & \B{0}{.1092} & \B{1}{.1904} & \B{0}{.0786} \\
Last neuron & \B{0}{.0152} & \B{0}{.0276} & \B{0}{.2702} & \B{1}{.1056} & \B{0}{.1914} & \B{1}{.0772} \\
\bottomrule
\end{tabular}
\caption{Robustness of Deterministic and Bayesian CNN model by neuron position}
\label{table:robbothneuloc}
\end{table}

Table \ref{table:robbothneuloc} shows the performance metrics from Figure \ref{fig:robdetmodelgran} grouped by neuron attack location of the SEU Injection. There is not a significant difference of the performance metric change in attacking the first neuron of the layer or the last neuron of the layer. The difference of performance for first neuron attack and last neuron attack might be attributed to the feature learning in our Bayesian CNN model and Deterministic CNN model. In our result, as the last neuron is more fragile to SEU, we can conclude that our model is trained to have the last neuron contribute more to the network.
\fi
% in design/implementation, please make sure you included how you define "first neuron" or "last neuron"

\iffalse
TODO explain how the first and last neuron of attack does not matter. You have suspicion on the probability of first/last neuron to be nullified by the dropout layer, so you bring this table. But you are unsure whether to include this table, as no significant finding is found. read more on how dropout layer is used in the inference, not training. they should be there right? they are not gone when you are doing inference, they are just gone during training to increase regularization.

%Summary of absolute accuracy difference and softmax difference by neuron location and model variant for Deterministic CNN
\begin{table}[H]
\centering
\begin{tabular}{lllll}
\toprule
location\_index & folder\_name & AAAD & ASD & ARIn \\
\midrule
First neuron & Model 0 & \B{0}{0.0150} & \B{0}{0.2701} & \B{0}{0.1913} \\
Last neuron & Model 0 & \B{0}{0.0141} & \B{0}{0.2688} & \B{0}{0.1903} \\
\midrule
First neuron & Model 1 & \B{1}{0.0135} & \B{0}{0.2687} & \B{0}{0.1902} \\
Last neuron & Model 1 & \B{0}{0.0148} & \B{0}{0.2711} & \B{0}{0.1920} \\
\midrule
First neuron & Model 2 & \B{0}{0.0144} & \B{0}{0.2687} & \B{0}{0.1903} \\
Last neuron & Model 2 & \B{0}{0.0170} & \B{0}{0.2717} & \B{0}{0.1925} \\
\midrule
First neuron & Model 3 & \B{0}{0.0156} & \B{1}{0.2680} & \B{1}{0.1898} \\
Last neuron & Model 3 & \B{0}{0.0148} & \B{0}{0.2694} & \B{0}{0.1908} \\
\bottomrule
\end{tabular}
\caption{TODO this table might be removed due to insignificance in result}
\label{table:mightberemoved}
\end{table}
\fi

%\section{Explana}

%a


% -----------------------------------------------------------------------------

\chapter{Conclusion}
\label{chap:conclusion}

\iffalse
{\bf A compulsory chapter,  roughly 10\% of the total page-count. Blackboard suggests 3 pages.}
\vspace{1cm} 

\noindent
The concluding chapter(s) of a dissertation are often underutilized because they're 
too often left too close to the deadline: it is important to allocate enough time and 
attention to closing off the story, the narrative, of your thesis.

Again, there is no single correct way of closing a thesis. 

One good way of doing this is to have a single chapter consisting of three parts:

\begin{enumerate}
\item (Re)summarise the main contributions and achievements, in essence
      summing up the content.
\item Clearly state the current project status (e.g., ``X is working, Y 
      is not'') and evaluate what has been achieved with respect to the 
      initial aims and objectives (e.g., ``I completed aim X outlined 
      previously, the evidence for this is within Chapter Y'').  There 
      is no problem including aims which were not completed, but it is 
      important to evaluate and/or justify why this is the case.
\item Outline any open problems or future plans.  Rather than treat this
      only as an exercise in what you {\em could} have done given more 
      time, try to focus on any unexplored options or interesting outcomes
      (e.g., ``my experiment for X gave counter-intuitive results, this 
      could be because Y and would form an interesting area for further 
      study'' or ``users found feature Z of my software difficult to use,
      which is obvious in hindsight but not during at design stage; to 
      resolve this, I could clearly apply the technique of Bloggs {\em et al.}.
\end{enumerate}

Alternatively, you might want to divide this content into two chapters: a penultimate chapter with a title such as ``Further Work" and then a final chapter ``Conclusions". Again, there is no hard and fast rule, we trust you to make the right decision. 

And this, the final paragraph of this thesis template, is just a bunch of citations, added to show how to generate a BibTeX bibliography. Sources that have been randomly chosen to be cited here include:
%\cite{miller_etal_2018_clojure,webber_marwan_2015,touretzky_2013_lisp,eckmann_etal_1987,marwan_2011,vach_2015,shiller_2017,vytelingum_2006,tesfatsion_2002,rust_etal_1992}.
\fi

%TODO CONCLUSION. Make sure it answers all your research questions, along any interesting findings in chapter 5. Make sure you add future works.


The purpose of this research is to examine the robustness of Bayesian Neural Networks designs against Single Event Upsets. The research starts by training both deterministic and Bayesian CNN Models, followed by SEU Injection on different model variants for each setting. In this chapter, we will reflect on our findings, the limitations encountered in this research, and suggestions of future work.


\section{Project Status}
%The research aims of this dissertation, as stated in Subsection \ref{subsection:researchaim} are to: (1) Train multiple Deterministic and Bayesian Neural Networks with different design choice and note its performance,(2) Do multiple experiment of SEU injection on the different Deterministic and Bayesian Neural Network design and note its performance, (3) Compare the robustness of each different design and (4) Validate whether Bayesian Neural Networks are more robust compared to Deterministic Neural Networks.

The research aims of this dissertation, as stated in Subsection \ref{subsection:researchaim} are as follows:

\begin{enumerate}
\item Train multiple deterministic and Bayesian Neural Networks with different design choice and note its performance
\item Do multiple experiment of SEU injection on the different deterministic and Bayesian Neural Network design and note its performance
\item Compare the robustness of each different design
\item Validate whether Bayesian Neural Networks are more robust compared to deterministic Neural Networks.
\end{enumerate}


\subsubsection{Training result}

Across all model variants and activation functions, \textbf{deterministic CNN model results in a better model performance}, which is test accuracy. The average test accuracy for Deterministic CNN model is 97.55\% and Bayesian CNN model only achieve 82.85\% for prior $b$ value equal to 1. This observation is suspected to happen due to the simplification of our Bayesian CNN model design in assuming that all weights and biases are initialized with the same prior and prior parameter. The other factor that can also contribute to this examination is due to the fact that KL Divergence, which is part of the ELBO loss used to train deterministic CNN model, is harder to optimize than cross entropy which is used in the deterministic CNN model.

For all model, activation, and prior variation, our \textbf{Bayesian CNN model with prior parameter $b$ equal to 1 results in a better test accuracy} of 82.85\% compared to the prior $b=0.1$ models that have the average test accuracy of 80.24\% and $b=10$ of $80.68\%$. The observed activation function that is better is either relu, relu6, or tanh.

\subsubsection{Robustness of different model designs against SEU Injection}

In this research, we test multiple model designs and compare its robustness for both deterministic CNN model and Bayesian CNN model. Here are the results as stated in Subsection \ref{section:SEUperformance}.

\begin{enumerate}
%\item In Figure \ref{fig:robdetmodelgran}, the characteristic of Deterministic CNN model that shown a trade-off between pure accuracy and robustness (1- softmax difference) against adversarial attack (FGSM) is also observed in our research by looking at the SEU robustness of Deterministic CNN Model.
\item \textbf{Model variations}: Dropout during training enhanced robustness against SEU in Bayesian CNNs, with model 2 showing the best both absolute accuracy difference and softmax difference. Weight decay (model 3) contributed to reduced softmax sensitivity in deterministic CNN, while the model with the best absolute accuracy difference is model 1 (model wih smartpooling).
\item \textbf{Prior and prior parameter}: The Gaussian prior is found to be the most robust, while Laplace and Uniform each showed better robustness under spread perturbations for absolute accuracy difference and softmax difference, compared to other priors, respectively. Across all priors, the center parameter proved to be more robust than the spread parameter. This is due to the fact that shifting the mean preserves relative sample positions, while changes in spread increase variability in logits.
\item \textbf{Activation functions}: Among all activation functions, Rectified Weighted Gaussian is the most robust in terms of absolute accuracy difference. The phenomenon of trade-off between test accuracy (which relies on the expressiveness of activation function) and SEU robustness is observed, as seen in the Relu activation function. Relu6, which is the bounded version of relu, shows better robustness against SEU injection. Sigmoid activation function shown better robustness in terms of softmax difference and ARIn.
\item \textbf{Layer and layer parameter}: Convolutional layers are observed to be more robust to SEU than the fully connected layer, as non-linear activations and pooling partially dampen perturbations before reaching the logits. On the other hand, SEUs in the fully connected layer directly affect the logits, resulting in larger accuracy and softmax difference. Additionally, weights are less robust than biases, since weight perturbations are amplified by the inputs, while biases act only as additive shifts.
\item \textbf{Attacked bit location}: SEU attacks on the first exponent bits caused the worst robustness quality, while the mantissa bit flips had a minimum effect.
\item \textbf{Original bit attacked}: Flipping a bit from 1 to 0 is found to be far more robust, especially in terms of accuracy, which is not aligned with a finding from previous research. The result is influenced by the fact that most weights and biases are closer to 0 or 1, meaning that the original bit state strongly determine the effect of SEU. In particular, first exponent bit flips, which is the most devastating bit flips, are impacting 0 bit values on all occassions, which explain the result that we have.
%\item \textbf{Neuron position}: SEU attacks on the first or last neurons showed little difference overall, although the last neuron is slightly more fragile. This suggests that in this specific case of our model and dataset, the later neurons may carry out more critical features, making the more influential in this SEU injection experiment.
\end{enumerate}

\subsubsection{Is Bayesian Neural Networks are more robust against deterministic Neural Networks against SEU injection?}

To conclude, our research question on whether Bayesian CNN model is more robust than the deterministic CNN model can be answered with Table \ref{table:robbothcompare}. Both model variation have their own strength, with deterministic CNN model that have better average absolute accuracy difference of 1.49\% compared to Bayesian CNN model that has 2.44\%. The Bayesian CNN Model is more robust in softmax difference (0.1074) against the deterministic CNN model (0.2686). Our newly defined Aggregate Robustness Index (ARIn), which aggregates both measures, shows that \textbf{Bayesian CNN model (0.0779) is more robust compared to deterministic CNN model (0.1909)}. Bayesian CNN model implementation benefits the scenario of which our Neural Networks takes class confidence/probability as one of the major aspect of its implementation.

\section{Limitations}

The ELBO loss used in the Bayesian CNN models training consists of the expected log-likelihood and KL divergence. In implementation, Pyro's \texttt{pyro.infer.trace\_elbo} does not expose them separately, preventing direct observation of their individual changes during training. This is especially useful if we want to modify the KL divergence term that went to $-\infty$ for uniform prior in Subsection \ref{subsection:bayesiantrainresult}.

This research assumes simplified priors with the same parameter across all weights and biases. To address this matter that results in lower performance compared to the deterministic CNN model, we could employ a warm start technique, in which we train the CNN models in the deterministic way first and use the resulting weight/bias as a starting point for our Bayesian CNN model. 

This work also assumes that all weights and biases do not correlate/depends with each other, which is the main assumption in mean-field variational guide. If memory and storage permits, instead of utilizing Pyro's \texttt{AutoNormal}, we can use \texttt{AutoMultivariateNormal} to have an automatic guide that also stores a Cholesky factorization, which covariance matrix can capture the correlations between latent variables.

While this research has covered 84 different Bayesian CNN models that each cover 168 different SEU injection occurences, the SEU injection we conducted with single bit-flips at a time, which may understate the real-world condition, in which Multi Bit Upset (MBU) can happen.

\section{Future Work}

\iffalse
\begin{enumerate}
\item Suggest in the future work to do the same with MCMC BNNs, SWAG, MC Dropout
\item Try more activation functions
\item Try more prior distribution
\item multi class problems
\item other metrics like f1 score for unbalanced dataset
\end{enumerate}
\fi

The CNN model architecture used in this experiment is quite small (2 convolutional layers and 1 fully connected layer with 52,162 parameters) on a dataset with 2 classes (Shipsnet). Dataset with more classes will offer a new perspective on how the softmax values are not symmetric, while a more complex architecture (e.g. ResNet, AlexNet, EfficientNet) resemble a more likely real-world application of models deployed in harsh environments.

%[Suggestion for other activation functions or prior distribution]

Testing our Bayesian Neural Networks with more variation of prior distribution is also a promising area of research. One of an interesting prior distribution to try on is the horseshoe prior, which is claimed to express robustness at handling unknown sparsity and large outlying signals \cite{horseshoecarvalho}.

Designing a Bayesian CNN model with another activation function also open up a new possibility of finding a more robust model than what we have in this research. Although the experiment setting is on the robustness against adversarial attack, work on \cite{pssilu}, stated that Parametric Activation Functions like Parametric Shifted Sigmoidal Linear Unit (PSSiLU) shows more robustness when combined with adversarial training.

Previous researches has mixed opinion on how Markov Chain Monte Carlo (MCMC) method can be more robust against adversarial attack against the deterministic model. Our research only uses Stochastic Variational Inference (SVI) from \texttt{Pyro} module, therefore it will be also compelling to find out how MCMC methods like Hamiltonian Monte Carlo (HMC) or No-U-Turn Sampler (NUTS) \cite{homanNUTS} performs, despite the increase in inference complexity.

A simpler approximate inference strategies can also be explored, which can include MC Dropout \cite{mcdropout} and Stochastic Weight Averaging-Gaussian (SWAG) \cite{swag}. In doing so, we can see how these simpler methods keeping their performance against SEU injection and how the complexity in terms of model and  training time can be saved.



% TODO:
%1. go through all todos
%2. check early chapters, make sure research question etc. us correct
%3. add more references or remove older ones
%4. cut the findings chapter about 5 to 10 pages. maybe we can remove the neuron attac location part. we also make explanations shorter.
%5. finally, give your recommendation for future works. 


% =============================================================================

% Finally, after the main matter, the back matter is specified.  This is
% typically populated with just the bibliography.  LaTeX deals with these
% in one of two ways, namely
%
% - inline, which roughly means the author specifies entries using the 
%   \bibitem macro and typesets them manually, or
% - using BiBTeX, which means entries are contained in a separate file
%   (which is essentially a database) then imported; this is the 
%   approach used below, with the databased being dissertation.bib.
%
% Either way, the each entry has a key (or identifier) which can be used
% in the main matter to cite it, e.g., \cite{X}, \cite[Chapter 2}{Y}.

\backmatter

\bibliography{PutawaraRevalda_mscthesis}



% -----------------------------------------------------------------------------

% The dissertation concludes with a set of (optional) appendices; these are 
% the same as chapters in a sense, but once signalled as being appendices via
% the associated macro, LaTeX manages them appropriately.

\appendix

\chapter{Appendix}
\label{appx:a}

\iffalse
Content which is not central to, but may enhance the dissertation can be 
included in one or more appendices; examples include, but are not limited
to

\begin{itemize}
\item lengthy mathematical proofs, numerical or graphical results which 
      are summarised in the main body,
\item sample or example calculations, 
      and
\item results of user studies or questionnaires.
\end{itemize}

\noindent
Note that in line with most research conferences, the examiners are not
obliged to read such appendices.
\fi

\textbf{Floating Point Numbers Architecture}

\begin{figure}[h]
\caption{Architecture for storing floating point numbers in this research, which is FP32}
\includegraphics[scale=0.8]{fp32_own.png}
\centering
\label{fig:floatingpointusedIllustrationown}
\end{figure}

\textbf{Sampled Weight Distribution Illustration}

\begin{figure}[H]
\caption{Sampled Weight Distribution before SEU (panel a), Uniform(-0.28,0.64) and after SEU (panel b), Uniform(1000.0,0.64)}
\includegraphics[scale=0.7]{sample_weight_distribution.png}
\centering
\label{fig:sampledweight}
\end{figure}

\pagebreak
\textbf{Train Accuracy of Bayesian CNN models}

\begin{table}[H]
\begin{tabular}{ll|rrrrrrr|r}
\toprule
Model Variant & prior & RWG & WG & relu & relu6 & sigmoid & sin & tanh & average \\
\midrule
Model 0 & Gaussian & \B{0}{.7469} & \B{0}{.7100} & \B{1}{.7503} & \B{1}{.7503} & \B{0}{.7434} & \B{0}{.7172} & \B{0}{.7150} & \B{0}{.7333} \\
Model 0 & Laplace & \B{0}{.6559} & \B{0}{.6056} & \B{0}{.7172} & \B{1}{.7212} & \B{0}{.7178} & \B{0}{.6266} & \B{0}{.6141} & \B{0}{.6655} \\
Model 0 & Uniform & \B{0}{.9216} & \B{1}{.9350} & \B{0}{.9306} & \B{0}{.9300} & \B{0}{.8625} & \B{0}{.9316} & \B{0}{.9131} & \B{0}{.9178} \\
\midrule
Model 1 & Gaussian & \B{0}{.7469} & \B{0}{.7100} & \B{1}{.7506} & \B{1}{.7506} & \B{0}{.7434} & \B{0}{.7172} & \B{0}{.7144} & \B{0}{.7333} \\
Model 1 & Laplace & \B{0}{.6562} & \B{0}{.6044} & \B{0}{.7169} & \B{1}{.7241} & \B{0}{.7184} & \B{0}{.6234} & \B{0}{.6138} & \B{0}{.6653} \\
Model 1 & Uniform & \B{1}{.9344} & \B{0}{.9297} & \B{0}{.9291} & \B{0}{.9275} & \B{0}{.8625} & \B{0}{.9313} & \B{0}{.9131} & \B{0}{.9182} \\
\midrule
Model 2 & Gaussian & \B{0}{.7434} & \B{0}{.6872} & \B{1}{.7441} & \B{0}{.7438} & \B{0}{.7359} & \B{0}{.6875} & \B{0}{.6887} & \B{0}{.7187} \\
Model 2 & Laplace & \B{0}{.6697} & \B{0}{.6128} & \B{0}{.7194} & \B{1}{.7225} & \B{0}{.7169} & \B{0}{.6453} & \B{0}{.6197} & \B{0}{.6723} \\
Model 2 & Uniform & \B{0}{.9319} & \B{1}{.9434} & \B{0}{.9181} & \B{0}{.9297} & \B{0}{.8581} & \B{0}{.9281} & \B{0}{.9275} & \B{1}{.9196} \\
\midrule
Model 3 & Gaussian & \B{0}{.7469} & \B{0}{.7106} & \B{0}{.7503} & \B{1}{.7506} & \B{0}{.7434} & \B{0}{.7172} & \B{0}{.7144} & \B{0}{.7333} \\
Model 3 & Laplace & \B{0}{.6569} & \B{0}{.6034} & \B{0}{.7156} & \B{1}{.7216} & \B{0}{.7188} & \B{0}{.6212} & \B{0}{.6147} & \B{0}{.6646} \\
Model 3 & Uniform & \B{0}{.9256} & \B{1}{.9331} & \B{0}{.9300} & \B{0}{.9278} & \B{0}{.8625} & \B{0}{.9319} & \B{0}{.9228} & \B{0}{.9191} \\
\midrule
 & Average & \B{0}{.7780} & \B{0}{.7488} & \B{0}{.7977} & \B{1}{.8000} & \B{0}{.7736} & \B{0}{.7565} & \B{0}{.7476} & \B{0}{.7717} \\
\bottomrule
\end{tabular}
\caption{Train accuracy of different activation functions across model variants in evaluating Bayesian CNN Models b=0.1}
\label{table:activation_performance_bayesian_train-1}
\end{table}

\begin{table}[H]
\begin{tabular}{ll|rrrrrrr|r}
\toprule
Model Variant & prior & RWG & WG & relu & relu6 & sigmoid & sin & tanh & average \\
\midrule
Model 0 & Gaussian & \B{0}{.6909} & \B{0}{.7678} & \B{1}{.8716} & \B{0}{.8250} & \B{0}{.6994} & \B{0}{.8231} & \B{0}{.7894} & \B{0}{.7810} \\
Model 0 & Laplace & \B{0}{.6678} & \B{0}{.6991} & \B{1}{.8703} & \B{0}{.7831} & \B{0}{.6975} & \B{0}{.7741} & \B{0}{.7706} & \B{0}{.7518} \\
Model 0 & Uniform & \B{0}{.7459} & \B{0}{.7778} & \B{0}{.8684} & \B{0}{.8363} & \B{0}{.7975} & \B{1}{.8694} & \B{0}{.8116} & \B{0}{.8153} \\
\midrule
Model 1 & Gaussian & \B{0}{.6984} & \B{0}{.7684} & \B{1}{.8662} & \B{0}{.8125} & \B{0}{.6966} & \B{0}{.8191} & \B{0}{.7772} & \B{0}{.7769} \\
Model 1 & Laplace & \B{0}{.6781} & \B{0}{.6953} & \B{1}{.8469} & \B{0}{.7863} & \B{0}{.6963} & \B{0}{.7750} & \B{0}{.7656} & \B{0}{.7491} \\
Model 1 & Uniform & \B{0}{.7434} & \B{0}{.7800} & \B{0}{.8744} & \B{0}{.8153} & \B{0}{.7947} & \B{1}{.8756} & \B{0}{.8194} & \B{0}{.8147} \\
\midrule
Model 2 & Gaussian & \B{0}{.7206} & \B{0}{.7803} & \B{1}{.8884} & \B{0}{.8631} & \B{0}{.6944} & \B{0}{.8559} & \B{0}{.8306} & \B{0}{.8048} \\
Model 2 & Laplace & \B{0}{.7013} & \B{0}{.7312} & \B{1}{.8788} & \B{0}{.8500} & \B{0}{.6866} & \B{0}{.7959} & \B{0}{.7978} & \B{0}{.7774} \\
Model 2 & Uniform & \B{0}{.7366} & \B{0}{.7816} & \B{0}{.8862} & \B{0}{.8694} & \B{0}{.8025} & \B{1}{.8903} & \B{0}{.8072} & \B{1}{.8248} \\
\midrule
Model 3 & Gaussian & \B{0}{.6947} & \B{0}{.7769} & \B{1}{.8719} & \B{0}{.8237} & \B{0}{.6950} & \B{0}{.8184} & \B{0}{.7772} & \B{0}{.7797} \\
Model 3 & Laplace & \B{0}{.6809} & \B{0}{.7006} & \B{1}{.8541} & \B{0}{.7816} & \B{0}{.6922} & \B{0}{.7709} & \B{0}{.7653} & \B{0}{.7494} \\
Model 3 & Uniform & \B{0}{.6737} & \B{0}{.7688} & \B{0}{.8747} & \B{0}{.8128} & \B{0}{.7919} & \B{1}{.8759} & \B{0}{.8100} & \B{0}{.8011} \\
\midrule
 & Average & \B{0}{.7027} & \B{0}{.7523} & \B{1}{.8710} & \B{0}{.8216} & \B{0}{.7287} & \B{0}{.8286} & \B{0}{.7935} & \B{0}{.7855} \\
\bottomrule
\end{tabular}
\caption{Train accuracy of different activation functions across model variants in evaluating Bayesian CNN Models b=1.0}
\label{table:activation_performance_bayesian_train-0}
\end{table}


\begin{table}[H]
\begin{tabular}{ll|rrrrrrr|r}
\toprule
Model Variant & prior & RWG & WG & relu & relu6 & sigmoid & sin & tanh & average \\
\midrule
Model 0 & Gaussian & \B{0}{.6453} & \B{0}{.6488} & \B{1}{.8328} & \B{0}{.6822} & \B{0}{.6944} & \B{0}{.6691} & \B{0}{.7797} & \B{0}{.7075} \\
Model 0 & Laplace & \B{0}{.6362} & \B{0}{.6484} & \B{1}{.8228} & \B{0}{.6822} & \B{0}{.6744} & \B{0}{.6647} & \B{0}{.7672} & \B{0}{.6994} \\
Model 0 & Uniform & \B{0}{.6784} & \B{0}{.6969} & \B{1}{.8666} & \B{0}{.6706} & \B{0}{.8169} & \B{0}{.6966} & \B{0}{.8178} & \B{0}{.7491} \\
\midrule
Model 1 & Gaussian & \B{0}{.6512} & \B{0}{.6538} & \B{0}{.7175} & \B{0}{.7025} & \B{0}{.7000} & \B{0}{.6744} & \B{1}{.7806} & \B{0}{.6971} \\
Model 1 & Laplace & \B{0}{.6481} & \B{0}{.6494} & \B{0}{.7056} & \B{0}{.6922} & \B{0}{.6731} & \B{0}{.6669} & \B{1}{.7659} & \B{0}{.6859} \\
Model 1 & Uniform & \B{0}{.6803} & \B{0}{.7078} & \B{1}{.8259} & \B{0}{.6741} & \B{0}{.8222} & \B{0}{.7034} & \B{0}{.8200} & \B{0}{.7477} \\
\midrule
Model 2 & Gaussian & \B{0}{.6872} & \B{0}{.6919} & \B{1}{.8331} & \B{0}{.7356} & \B{0}{.7228} & \B{0}{.6969} & \B{0}{.8122} & \B{0}{.7400} \\
Model 2 & Laplace & \B{0}{.6750} & \B{0}{.6700} & \B{1}{.8134} & \B{0}{.7156} & \B{0}{.7094} & \B{0}{.6634} & \B{0}{.8016} & \B{0}{.7212} \\
Model 2 & Uniform & \B{0}{.7094} & \B{0}{.7212} & \B{1}{.8869} & \B{0}{.7794} & \B{0}{.8541} & \B{0}{.6878} & \B{0}{.8425} & \B{1}{.7830} \\
\midrule
Model 3 & Gaussian & \B{0}{.6575} & \B{0}{.6556} & \B{1}{.8347} & \B{0}{.6959} & \B{0}{.7006} & \B{0}{.6847} & \B{0}{.7809} & \B{0}{.7157} \\
Model 3 & Laplace & \B{0}{.6353} & \B{0}{.6581} & \B{1}{.8187} & \B{0}{.6847} & \B{0}{.6719} & \B{0}{.6847} & \B{0}{.7681} & \B{0}{.7031} \\
Model 3 & Uniform & \B{0}{.6728} & \B{0}{.6956} & \B{1}{.8656} & \B{0}{.6831} & \B{0}{.8284} & \B{0}{.6956} & \B{0}{.8191} & \B{0}{.7515} \\
\midrule
 & Average & \B{0}{.6647} & \B{0}{.6748} & \B{1}{.8186} & \B{0}{.6998} & \B{0}{.7390} & \B{0}{.6823} & \B{0}{.7963} & \B{0}{.7251} \\
\bottomrule
\end{tabular}
\caption{Train accuracy of different activation functions across model variants in evaluating Bayesian CNN Models b=10.0}
\label{table:activation_performance_bayesian_train+1}
\end{table}

\pagebreak
\textbf{Test Accuracy of Bayesian CNN models}

\begin{table}[H]
\begin{tabular}{ll|rrrrrrr|r}
\toprule
Model Variant & prior & RWG & WG & relu & relu6 & sigmoid & sin & tanh & average \\
\midrule
Model 0 & Gaussian & \B{0}{.7500} & \B{0}{.7225} & \B{0}{.7837} & \B{1}{.7850} & \B{0}{.7500} & \B{0}{.7588} & \B{0}{.7475} & \B{0}{.7568} \\
Model 0 & Laplace & \B{0}{.7500} & \B{0}{.7075} & \B{1}{.7700} & \B{0}{.7688} & \B{0}{.7500} & \B{0}{.7438} & \B{0}{.7150} & \B{0}{.7436} \\
Model 0 & Uniform & \B{0}{.9187} & \B{0}{.9363} & \B{0}{.9213} & \B{0}{.9250} & \B{0}{.8413} & \B{1}{.9387} & \B{0}{.8988} & \B{1}{.9114} \\
\midrule
Model 1 & Gaussian & \B{0}{.7500} & \B{0}{.7225} & \B{1}{.7850} & \B{1}{.7850} & \B{0}{.7500} & \B{0}{.7588} & \B{0}{.7475} & \B{0}{.7570} \\
Model 1 & Laplace & \B{0}{.7500} & \B{0}{.7087} & \B{1}{.7700} & \B{1}{.7700} & \B{0}{.7500} & \B{0}{.7475} & \B{0}{.7150} & \B{0}{.7445} \\
Model 1 & Uniform & \B{0}{.9100} & \B{0}{.9313} & \B{0}{.9237} & \B{0}{.9187} & \B{0}{.8387} & \B{1}{.9387} & \B{0}{.8988} & \B{0}{.9086} \\
\midrule
Model 2 & Gaussian & \B{0}{.7500} & \B{0}{.7412} & \B{1}{.7825} & \B{0}{.7788} & \B{0}{.7500} & \B{0}{.7638} & \B{0}{.7512} & \B{0}{.7596} \\
Model 2 & Laplace & \B{0}{.7462} & \B{0}{.7175} & \B{1}{.7987} & \B{0}{.7950} & \B{0}{.7400} & \B{0}{.7512} & \B{0}{.7312} & \B{0}{.7543} \\
Model 2 & Uniform & \B{0}{.9213} & \B{1}{.9500} & \B{0}{.9113} & \B{0}{.9113} & \B{0}{.8363} & \B{0}{.7975} & \B{0}{.9137} & \B{0}{.8916} \\
\midrule
Model 3 & Gaussian & \B{0}{.7500} & \B{0}{.7225} & \B{0}{.7850} & \B{1}{.7863} & \B{0}{.7500} & \B{0}{.7588} & \B{0}{.7488} & \B{0}{.7573} \\
Model 3 & Laplace & \B{0}{.7500} & \B{0}{.7050} & \B{1}{.7700} & \B{0}{.7688} & \B{0}{.7500} & \B{0}{.7425} & \B{0}{.7175} & \B{0}{.7434} \\
Model 3 & Uniform & \B{0}{.9237} & \B{0}{.8662} & \B{0}{.9225} & \B{0}{.9062} & \B{0}{.8387} & \B{1}{.9413} & \B{0}{.9062} & \B{0}{.9007} \\
\midrule
 & Average & \B{0}{.8058} & \B{0}{.7859} & \B{1}{.8270} & \B{0}{.8249} & \B{0}{.7787} & \B{0}{.8034} & \B{0}{.7909} & \B{0}{.8024} \\
\bottomrule
\end{tabular}
\caption{Test accuracy of different activation functions across model variants in evaluating Bayesian CNN Models b=0.1}
\label{table:activation_performance_bayesian_test-1}
\end{table}

\begin{table}[H]
\begin{tabular}{ll|rrrrrrr|r}
\toprule
Model Variant & prior & RWG & WG & relu & relu6 & sigmoid & sin & tanh & average \\
\midrule
Model 0 & Gaussian & \B{0}{.7562} & \B{0}{.7650} & \B{1}{.9137} & \B{0}{.8825} & \B{0}{.8025} & \B{0}{.7612} & \B{0}{.8938} & \B{0}{.8250} \\
Model 0 & Laplace & \B{0}{.7575} & \B{0}{.7538} & \B{1}{.8975} & \B{0}{.8562} & \B{0}{.7675} & \B{0}{.7500} & \B{0}{.8750} & \B{0}{.8082} \\
Model 0 & Uniform & \B{0}{.8063} & \B{0}{.8263} & \B{1}{.9025} & \B{0}{.8712} & \B{0}{.8237} & \B{0}{.8762} & \B{0}{.8662} & \B{0}{.8532} \\
\midrule
Model 1 & Gaussian & \B{0}{.7600} & \B{0}{.7612} & \B{1}{.9137} & \B{0}{.8788} & \B{0}{.8050} & \B{0}{.7512} & \B{0}{.8925} & \B{0}{.8232} \\
Model 1 & Laplace & \B{0}{.7475} & \B{0}{.7550} & \B{1}{.9113} & \B{0}{.8700} & \B{0}{.7650} & \B{0}{.7500} & \B{0}{.8762} & \B{0}{.8107} \\
Model 1 & Uniform & \B{0}{.8325} & \B{0}{.8125} & \B{0}{.9125} & \B{0}{.8762} & \B{0}{.8237} & \B{1}{.9150} & \B{0}{.8875} & \B{1}{.8657} \\
\midrule
Model 2 & Gaussian & \B{0}{.7525} & \B{0}{.7588} & \B{1}{.9100} & \B{0}{.8925} & \B{0}{.7525} & \B{0}{.7825} & \B{0}{.8862} & \B{0}{.8193} \\
Model 2 & Laplace & \B{0}{.7488} & \B{0}{.7588} & \B{1}{.9100} & \B{0}{.8825} & \B{0}{.7488} & \B{0}{.7512} & \B{0}{.8738} & \B{0}{.8105} \\
Model 2 & Uniform & \B{0}{.7400} & \B{0}{.8213} & \B{0}{.9150} & \B{0}{.8875} & \B{0}{.8225} & \B{1}{.9175} & \B{0}{.8662} & \B{0}{.8529} \\
\midrule
Model 3 & Gaussian & \B{0}{.7538} & \B{0}{.7612} & \B{1}{.9200} & \B{0}{.8825} & \B{0}{.7963} & \B{0}{.7625} & \B{0}{.8925} & \B{0}{.8241} \\
Model 3 & Laplace & \B{0}{.7500} & \B{0}{.7488} & \B{1}{.9175} & \B{0}{.8738} & \B{0}{.7625} & \B{0}{.7512} & \B{0}{.8812} & \B{0}{.8121} \\
Model 3 & Uniform & \B{0}{.7362} & \B{0}{.8225} & \B{1}{.9237} & \B{0}{.7512} & \B{0}{.8000} & \B{0}{.9163} & \B{0}{.9075} & \B{0}{.8368} \\
\midrule
 & Average & \B{0}{.7618} & \B{0}{.7787} & \B{1}{.9123} & \B{0}{.8671} & \B{0}{.7892} & \B{0}{.8071} & \B{0}{.8832} & \B{0}{.8285} \\
\bottomrule
\end{tabular}
\caption{Test accuracy of different activation functions across model variants in evaluating Bayesian CNN Models b=1}
\label{table:activation_performance_bayesian_test-0}
\end{table}

\begin{table}[H]
\begin{tabular}{ll|rrrrrrr|r}
\toprule
Model Variant & prior & RWG & WG & relu & relu6 & sigmoid & sin & tanh & average \\
\midrule
Model 0 & Gaussian & \B{0}{.7262} & \B{0}{.7300} & \B{0}{.8550} & \B{0}{.8237} & \B{0}{.8100} & \B{0}{.7425} & \B{1}{.8738} & \B{0}{.7945} \\
Model 0 & Laplace & \B{0}{.7562} & \B{0}{.7550} & \B{0}{.8612} & \B{0}{.8113} & \B{0}{.7712} & \B{0}{.7488} & \B{1}{.8700} & \B{0}{.7962} \\
Model 0 & Uniform & \B{0}{.7612} & \B{0}{.7975} & \B{1}{.9075} & \B{0}{.7600} & \B{0}{.8962} & \B{0}{.7500} & \B{0}{.9062} & \B{0}{.8255} \\
\midrule
Model 1 & Gaussian & \B{0}{.7538} & \B{0}{.7488} & \B{0}{.7913} & \B{0}{.8137} & \B{0}{.8275} & \B{0}{.7475} & \B{1}{.8738} & \B{0}{.7938} \\
Model 1 & Laplace & \B{0}{.7488} & \B{0}{.7462} & \B{0}{.7850} & \B{0}{.7950} & \B{0}{.7662} & \B{0}{.7500} & \B{1}{.8688} & \B{0}{.7800} \\
Model 1 & Uniform & \B{0}{.7488} & \B{0}{.8287} & \B{0}{.9075} & \B{0}{.8137} & \B{0}{.8912} & \B{0}{.7500} & \B{1}{.9150} & \B{0}{.8364} \\
\midrule
Model 2 & Gaussian & \B{0}{.7488} & \B{0}{.7500} & \B{0}{.8562} & \B{0}{.8087} & \B{0}{.7675} & \B{0}{.7500} & \B{1}{.8850} & \B{0}{.7952} \\
Model 2 & Laplace & \B{0}{.7488} & \B{0}{.7512} & \B{0}{.8488} & \B{0}{.7825} & \B{0}{.7550} & \B{0}{.7488} & \B{1}{.8600} & \B{0}{.7850} \\
Model 2 & Uniform & \B{0}{.7775} & \B{0}{.7900} & \B{0}{.8988} & \B{0}{.8688} & \B{0}{.8912} & \B{0}{.7500} & \B{1}{.9062} & \B{1}{.8404} \\
\midrule
Model 3 & Gaussian & \B{0}{.7525} & \B{0}{.7500} & \B{0}{.8575} & \B{0}{.8213} & \B{0}{.8187} & \B{0}{.7488} & \B{1}{.8812} & \B{0}{.8043} \\
Model 3 & Laplace & \B{0}{.7500} & \B{0}{.7562} & \B{0}{.8612} & \B{0}{.8050} & \B{0}{.7700} & \B{0}{.7500} & \B{1}{.8688} & \B{0}{.7945} \\
Model 3 & Uniform & \B{0}{.7750} & \B{0}{.8050} & \B{0}{.9113} & \B{0}{.8037} & \B{0}{.8962} & \B{0}{.7500} & \B{1}{.9125} & \B{0}{.8362} \\
\midrule
 & Average & \B{0}{.7540} & \B{0}{.7674} & \B{0}{.8618} & \B{0}{.8090} & \B{0}{.8218} & \B{0}{.7489} & \B{1}{.8851} & \B{0}{.8068} \\
\bottomrule
\end{tabular}
\caption{Test accuracy of different activation functions across model variants in evaluating Bayesian CNN Models b=10}
\label{table:activation_performance_bayesian_test+1}
\end{table}

\textbf{Robustness of Deterministic and Bayesian CNN models}

\begin{figure}[H]
\caption{Robustness of Deterministic CNN model by model variant}
\includegraphics[scale=0.5]{robmodeldetgran_new.png}
\centering
\label{fig:robdetmodelgran}
\end{figure}

\begin{figure}[H]
\caption{Robustness of Bayesian CNN model by model variant}
\includegraphics[scale=0.5]{robbaymodelgran.png}
\centering
\label{fig:robbaymodelgran}
\end{figure}

%Summary of how bitflips can impact the value stored in weights or bias
\begin{table}[h]
\centering
\begin{tabular}{lr}
\toprule
Bit Index & After Value \\
\midrule
0 &  0.60335773229599\\
\midrule
1 &  -2.0531199724572943e+38\\
3 &  -1.404801691640145e-10\\
6 &  -0.037709858268499374\\
\midrule
10 &  -0.72835773229599\\
15 &  -0.60726398229599\\
21 &  -0.60329669713974\\
\bottomrule
\end{tabular}
\caption{Illustration of value change after SEU (original value = -0.60335773229599)}
\label{table:robbothbitloc2}
\end{table}

\begin{figure}[H]
\caption{Log mean absolute difference against robustness metrics, excluding the data from bit index=1}
\includegraphics[scale=0.5]{logabsdiff2.png}
\centering
\label{fig:logabsdiff2}
\end{figure}

\textbf{GitHub Repository}
\vspace{1em}

The source code used in this research is available at \url{https://github.com/putawararevalda/bnndesignrobustness}.

% =============================================================================

\end{document}
